% !TEX Program = xelatex
% 1. `print`/`nonprint`：该选项控制是否以印刷模式生成文档，印刷模式会自动在论文的前置部分添加必要的空白页，默认为`nonprint`。注意标点符号为英文!
% 草稿使用
\documentclass[print, master, vlined]{DissertUESTC}


\begin{document}
    % 此命令可让LaTeX像Word那样直接堆叠页面内容，而不再默认拉伸段间距，代价是页尾可能会有明显空白
    \raggedbottom
    
    % 封面
	\uestccover{异构星座组网架构及路由机制研究}
				{信息与通信工程}
				{202321011316}
				{卢培亮}
				{章小宁}
				{教授}%{终南山古墓派}
				{信息与通信工程学院}  % 默认将过长的学院名称压缩至下划线宽度


	% 中文扉页，仅研究生用

	\ClsNum{}  % \ClsNum{<分类号>}
	\ClsLv{}  % \ClsLv{<密级>}
	\UDC{}  % \UDC{<UDC号>}
	\DissertationTitle{异构星座组网架构及路由机制研究}  % \DissertationTitle{<题名>}
	\Author{卢培亮}  % \Author{<作者姓名>}
	\Supervisor{章小宁}{教授}{电子科技大学}{成都}  % \Supervisor{<指导教师>}{<职称>}{<单位名称>}{<单位地址>}
	
	\DegLv{硕士}  % \DegLv{<申请学位级别>}，该信息由文档类选项自动确定，需修改默认内容时使用，否则注释即可
	\Major{信息与通信工程}  % \Major{<学科专业>}

	\Date{TODO}{TODO}  % \Date{<论文提交日期>}{<论文答辩日期>}
	\Grant{电子科技大学}{2026年6月}  % \Grant{<学位授予单位>}{<学位授予日期>}
	\Reviewer{todo}{todo}  % \Reviewer{<答辩委员为主席>}{<评阅人>}
	\uestczhtitlepage  % 过长评阅人文本将被自动压缩至下划线宽度，不会超出页面边界	

	% 英文扉页，仅研究生用
	% \uestcentitlepage{<文题>}{<专业>}{<学号>}{<作者>}{<导师>}{<副导师>}{<学院>}，若无副导师，则将“<副导师>”参数留空即可
	\uestcentitlepage{Research on Networking Architecture and Routing Mechanism for Heterogeneous Satellite Constellations}
					{Information and Communication Engineering}
					{202321011316}
					{Peiliang Lu}
					{Prof. Xiaoning Zhang}{}
					{School of Information and Communication Engineering}
	
	% 独创性声明：[签名宽度]{日期}{作者签名图片1}[作者签名图片2 <默认值为作者签名图片1>]{导师签名图片}，仅研究生用
	% \declaration{}{}{}  % 空白独创性声明
	% \declaration[3cm]{2024年08月31日}{authsign}{spvrsign}
	% \declaration[3cm]{2024年08月31日}{authsign}[authsign]{spvrsign}
	\declaration[3cm]{2024年08月31日}{}[]{}
 

	% 开启中文摘要
	\zhabstract
    随着空间信息技术的飞速发展，通导遥一体化建设已成为天地一体化信息网络发展的核心方向，未来的空间网络需同时承载宽带接入、对地观测及天机边缘计算等多样化业务。构建由地球静止轨道、中地球轨道与低地球轨道深度融合的异构星座网络，是满足上述需求的关键基础设施。然而，面对异构网络节点规模庞杂、拓扑结构高度动态变化以及业务服务质量需求差异显著等挑战，传统的组网架构与路由机制在全局管控效率与多维业务保障方面存在明显不足。本文围绕异构星座的组网架构、路由算法及仿真平台验证展开了深入研究。

    针对异构星座网络节点规模大、拓扑高度动态的难题，提出了一种基于软件定义网络的分级多域协同管控架构。采用GEO全局主控、LEO域内协同的分层机制，利用GEO维护全网虚拟拓扑并实施跨域策略，LEO层依据地面网格划分自治域，通过虚拟节点实现物理与逻辑解耦。此外，设计了基于虚拟地址的异构互联机制，通过双层报文封装屏蔽底层卫星移动性，实现了透明互联与无感切换。分析表明，该架构有效降低了信令开销，解决了动态管控难题，提升了系统稳定性。
    
    针对多域网络环境下业务QoS需求多样化及网络拥塞问题，提出了基于虚拟节点分层图模型的分层协同QoS路由机制。域间由GEO基于抽象拓扑规划跨域路径并分解QoS预算，域内设计了基于方向引导与拥塞感知的改进遗传算法，利用类曼哈顿网格特性引入方向引导初始化策略，并构建包含非线性拥塞惩罚项的适应度函数以主动规避热点链路。仿真显示，该机制在满足时延、丢包率等多维约束的同时，显著实现了跨层负载均衡，提升了网络吞吐量。
    
   针对传统仿真真实性不足的问题，设计并实现了基于容器虚拟化技术的异构星座网络仿真平台。平台集成STK与Docker技术，利用STK生成高精度轨道数据，通过Redis驱动虚拟网络，结合Open vSwitch与Linux流量控制模块构建可编程虚拟链路，真实模拟传播时延、带宽限制与动态切换，并部署控制器集群实现分级协同。验证表明，平台突破了协议栈简化局限，支持真实业务端到端传输，提供了高保真的半实物验证环境。
   
	% 中文关键字
	\zhkeywords{异构星座；软件定义卫星网络；虚拟节点；QoS路由；遗传算法；虚拟化仿真}

 
	
	% 开启英文摘要
	\enabstract
	With the development of space-ground integrated information networks, heterogeneous satellite constellations composed of GEO, MEO, and LEO have become critical infrastructures for global seamless coverage. However, the high mobility of large-scale LEO satellites leads to highly dynamic topology changes. Coupled with the differentiated QoS requirements of diverse services, this poses severe challenges to network management architecture and routing mechanisms. Addressing these issues, this thesis investigates the networking architecture, cooperative routing algorithms, and simulation platform construction for heterogeneous constellations. The main contributions are as follows:
    
    Firstly, aiming at the management challenges caused by the large scale and dynamic topology of heterogeneous constellations, a Hierarchical Domain Control Architecture based on SDN is proposed. This architecture adopts a "GEO Global Control and LEO Intra-domain Coordination" mode. A virtual node construction strategy based on geographic grids and a virtual address interconnection mechanism are designed to abstract the high-dynamic physical topology into a logically static virtual network, effectively masking the frequent topology reconstruction and significantly reducing control signaling overhead.
    
    Secondly, addressing the path planning problem under multi-dimensional QoS constraints, a Collaborative Routing Mechanism for Heterogeneous Constellations is proposed. This mechanism employs a strategy of Inter-domain Macro Planning and Intra-domain Micro Optimization. In the intra-domain phase, an Improved Genetic Algorithm based on direction guidance and congestion awareness is introduced. By incorporating geometric distance constraints for initialization and congestion-aware mutation operators, this algorithm overcomes the slow convergence of traditional algorithms and achieves load balancing while meeting QoS constraints.
    
    Finally, to verify the effectiveness of the proposed architecture and algorithms, a Heterogeneous Satellite Network Emulation Platform based on Container Virtualization is designed and implemented. The platform uses STK to drive high-fidelity orbital data and combines Docker containers with Open vSwitch to build a programmable virtual network environment. Simulation results demonstrate that the proposed architecture offers significant advantages in signaling overhead and link stability compared to traditional architectures. Furthermore, the H-IGA routing algorithm effectively reduces end-to-end delay and packet loss rate while improving network throughput under non-uniform traffic loads.

	% 英文关键字
	\enkeywords{Heterogeneous Constellation; Software Defined Networking (SDN); Hierarchical Domain Architecture; QoS Routing; Improved Genetic Algorithm; Network Emulation}
	
	\tableofcontents  % 主目录，必要

    % 论文正文章节
	\chapter{绪论}
    \section{研究背景与意义}

    随着空间信息技术的发展与国家战略需求的升级，天地一体化信息网络的功能正逐渐从单一的通信传输向通信、导航、遥感综合信息服务体系演进。未来的空间网络不仅需要提供全球覆盖的宽带接入服务，还需承担高分辨率对地观测数据的实时回传、天基边缘计算以及高精度时空基准服务等多样化任务。这些业务在时延敏感度、带宽需求及覆盖连续性方面存在显著差异。受限于单一轨道高度的物理特性，仅依靠同步轨道卫星或单一低轨星座难以同时满足上述多维度的服务质量需求。因此，构建由地球静止轨道（GEO）、中地球轨道（MEO）与低地球轨道（LEO）深度融合的异构卫星网络，利用不同轨道的优势互补实现多星座协同，已成为空间信息网络发展的必然趋势。

    在此背景下，全球范围内的大规模低轨星座建设已进入实质性部署阶段。在国外，SpaceX 公司的 Starlink 星座\citess{hui2025review}通过大规模部署卫星及星间激光链路技术，初步实现了全球宽带覆盖。OneWeb 星座\citess{handley2023new} 则侧重于提供高纬度及极地区域的通信服务。在国内，为了构建自主可控的空间基础设施，我国已启动了以国网\citess{gillinger2025developing}和千帆\citess{yu2025analyses}为主力的万级低轨卫星星座建设规划。这些计划共同预示着未来空间网络将呈现出节点规模巨大、轨道类型异构、全域立体覆盖的发展特征。
    
    特别是为应对日益复杂多变的太空任务需求，构建能够同时支持不同类型、不同功能任务的异构星座协同架构，已成为提升系统效能的关键。美军航天发展局提出的扩散型作战人员太空架构（Proliferated Warfighter Space Architecture, PWSA）\citess{10562491}，突破了传统单一任务星座的局限，设计了包含传输、跟踪、监视、导航、作战管理等七大功能层的异构体系。这些层依据任务功能进行划分，既涵盖了提供全球低时延数据传输的通信任务，也囊括了导弹预警跟踪、情报监视侦察以及定位导航授时等感知与保障任务。PWSA 的实践表明，通过构建这种多功能融合的异构星座网络，能够有效打破不同任务间的壁垒，实现通信、感知、指控等多样化空间任务的深度协同与高效处理。
    
    随着星座规模的增长和任务复杂度的提升，传统的基于透明转发或静态预配置的卫星网络协议体系已难以适应新的需求，网络架构正向基于 IP 互联、支持星上处理和智能路由的方向演进。然而，面对 GEO、MEO 与 LEO 深度融合带来的时变拓扑、异构协议兼容性以及多维服务质量需求，传统的分布式路由机制面临收敛速度慢、拥塞控制困难、资源利用率低等挑战。单纯依赖星上计算资源进行全网路由计算，在万级节点规模下将产生巨大的开销与时延。
    
    针对上述问题，引入软件定义网络（SDN）技术\citess{kreutz2014software}，构建控制与转发分离、分层分域的灵活管控架构，成为解决异构星座组网难题的关键技术途径。通过 SDN 架构，可以将复杂的路由决策从资源受限的卫星节点卸载至功能强大的控制器集群，实现对异构星座资源的统一抽象与全局优化。因此，深入研究面向异构星座的 SDN 组网架构及高效 QoS 路由机制，不仅能够提升我国未来空间信息网络的自主可控能力与关键业务保障水平，也为抢占下一代空天信息通信技术的制高点提供了重要的理论支撑与工程验证参考。
    

    \section{国内外研究现状}
    \subsection{异构星座组网与管控架构研究现状}
    随着卫星互联网向大规模、异构化方向发展，传统的基于专用硬件和固定功能的卫星组网架构已难以适应业务需求的动态变化。为了解决异构星座资源利用率低、管控僵化等问题，现有研究引入SDN技术，通过控制与转发分离的思想重构卫星网络架构，推动了组网模式从底层资源虚拟化向多层协同管控的演进。

    针对卫星网络节点众多、拓扑高动态的特点，早期的研究主要致力于构建高效的控制平面，探讨如何将 SDN 技术引入卫星网络以实现灵活性管理。文献\cite{ahmed2018software} 提出了一种基于云架构的卫星接入网框架（SatCloudRAN），探讨了利用 SDN 技术实现卫星网关的集中式控制与带宽资源的按需分配，解决了传统刚性架构下资源调度不灵活的问题。文献\cite{JSTY201702007} 研究了天地一体化网络中的集中管理机制，提出通过 SDN 控制器对天基设备进行统一配置以提升服务质量。文献\cite{袁梦珠2017基于SDN的卫星网络关键技术研究}针对 GEO/LEO 双层场景，将分组管理概念与 SDN 集中控制相结合，设计了由 GEO 卫星作为中转中枢的架构，负责收集 LEO 卫星群状态并将路由决策下发，实现了对大规模节点的批量管理。文献\cite{石晓东2020基于} 进一步细化了这一思路，明确了 GEO 层生成全局视图、LEO 层专注数据转发的功能划分，这种控制与转发的物理分离有效提升了网络的 QoS 性能。
    
    随着研究的深入，针对更复杂的 LEO/MEO/GEO 三层及星地融合场景，管控架构的设计逐渐转向实现全维度的统一协同。文献\cite{bao2014opensan} 提出了 OpenSAN 架构，旨在解决传统卫星网络配置低效的问题。该架构将控制平面部署于由至少三颗 GEO 卫星组成的GEO Group中，统一管理包含 MEO/LEO 在内的整个数据平面，并运行流表匹配协议支持组播等功能，显著减少了对地面信关站的依赖。文献\cite{li2017service} 则提出了名为 SERvICE 的综合天地一体化框架，采用管理平面、控制平面、数据平面的三层逻辑。其管理平面由卫星网络管理中心（SNMC）构成，控制平面分布于 GEO 卫星和地面数据中心，数据平面则涵盖了 MEO/LEO 卫星及地面设备。该框架通过多层控制器的协同部署，实现了对星地异构资源的集中管控与精细化流量工程。

    综上，传统的卫星组网模式已难以应对异构星座轨道时延差异大、节点高动态变化以及多业务并行带来的管控压力。引入 SDN 技术，利用其转控分离与逻辑集中的特性，实现对跨轨道层级资源的统一抽象，不仅能有效降低管控信令的交互开销，更能够为复杂业务场景下的跨域资源调度与 QoS 细粒度保障提供灵活的底层架构支撑。
    
    % 尽管上述研究在架构灵活性与资源利用率上取得了显著进展，但在面对未来万级规模的异构星座时，仍存在局限性。首先，现有的 GEO 主控或地面集中式架构，在面对成千上万颗 LEO 卫星的高频拓扑变化时，控制平面的信令开销巨大，容易形成通信瓶颈。其次，OpenSAN 等架构虽然实现了分层，但缺乏针对地理区域的分域自治机制，当 LEO 卫星高速跨越区域边界时，缺乏高效的控制器迁移与状态同步策略。最后，现有的架构多侧重于同构协议内的管理，对于不同星座间的互联缺乏透明的地址映射与封装机制。因此，设计一种具备分级多域特征、能够适应高动态拓扑变化的 SDN 协同管控架构，是实现异构星座高效组网的关键。

    \subsection{异构星座协同路由技术研究现状}
    % =======================新写========================
    异构星座网络的路由技术是实现多层轨道资源协同与数据高效传输的核心。针对异构卫星节点在覆盖范围、运动特性及链路状态上的差异，现有研究主要沿着分层分组路由与SDN协同管控路由两条技术路线演进，并不断向QoS保障方向发展。
    
    在早期的异构组网研究中，核心思路是利用高轨卫星的广域覆盖特性来辅助低轨卫星进行分层管理，以降低全网路由计算的复杂度。针对多层卫星网络架构，通过收集各层级延迟信息构建路由表、或利用高轨卫星覆盖范围对低轨卫星卫星进行分组管理的策略被相继提出\citess{akyildiz2002mlsr, 10.1145/570790.570809}。在此分层架构基础上，文献\cite{10.1007/s11276-013-0656-z}针对层间链路或节点故障的生存性机制与快速重路由策略，从而增强了异构网络在动态环境下的鲁棒性。

    随着SDN技术的引入，研究重点转向了控制与转发分离的集中式管控架构，以实现异构资源的灵活配置。文献\cite{7017183}设计了一种基于SDN的卫星网络架构，利用GEO卫星收集LEO层的网络状态信息，并将其转发至地面网络运行控制中心（NOCC）进行路由计算。文献\cite{bao2014opensan}也采用了类似的思路，利用GEO卫星辅助底层卫星网络的控制管理和路由计算，实现了控制层面的解耦与协同。针对更复杂的空天地一体化场景，文献\cite{li2017service}基于软件定义的空地一体化通信架构，结合启发式路由算法，实现了针对低时延或高带宽业务的异构资源协同调度。文献\cite{8766918}则构建了基于SDN的三层卫星通信网络模型，并实现了一种自适应路由算法（ARA），利用控制器的全局视角优化跨层路由路径。

    现在业务需求从单一通信向通导遥综合服务演进，单纯的最小跳数算法已难以满足带宽、时延、丢包率等多维 QoS 约束，研究者开始引入启发式算法以解决这一复杂的路径寻优问题。文献\cite{石晓东2020基于} 将 GEO 卫星作为 SDN 控制器，建立多 QoS 约束的网络模型，根据业务需求动态调整链路权重，实现了异构网络资源的统一调控。文献\cite{zhou2007novel} 提出了一种层级分布式 QoS 路由协议，通过虚拟节点法屏蔽 LEO 拓扑的动态性，并将路由问题建模为带宽约束下的最小延迟问题。文献\cite{long2010sustainable} 在此基础上，综合考虑端到端时延、链路利用率和丢包率，尝试利用蚁群算法和遗传算法解决 MEO/LEO 网络中的多目标路径寻优问题。针对网络负载分配问题，文献\cite{冯玺宝2016多层卫星网络} 提出了基于簇首管理的 QoS 路由策略，利用改进的最小跳算法预计算备选路径，同时结合流量预测技术，利用神经网络预测地面接入业务量，以降低星间传输业务的丢包率。文献\cite{1024075130.nh}利用 GEO 卫星作为控制层，通过DQN模型规划LEO卫星的路径，以最小化端到端时延。

    综上，异构星座协同路由机制正由单一目标优化向多维 QoS 约束下的智能协同转变。虽然现有基于 SDN 的管控架构有效提升了异构资源的调度灵活性，但在庞杂的拓扑下如何通过高效算法主动规避拥塞并平衡负载，仍是兼顾全局寻优效率与局部响应实时性的关键难题，也已成为保障未来异构星座网络综合服务能力的研究趋势。
    % =======================新写========================
    
    % 尽管上述研究在异构路由机制上取得了丰富成果，但在面对大规模高动态星座及非均匀业务负载时，部分关键问题仍有待进一步优化。首先，在拓扑管理方面，基于离散时间片的快照路由虽然有效简化了动态性，但传统方法中预计算的路由表往往难以灵活应对时间片内部可能发生的瞬时链路拥塞，容易在热点区域形成流量瓶颈。其次，在算法效率方面，标准的启发式算法在缺乏网络拓扑先验知识引导的情况下，往往面临搜索盲目性大、寻优效率不稳定的挑战。因此，研究一种结合拓扑规律进行方向引导、并具备动态拥塞感知能力的改进路由算法，对于提升异构星座网络的传输效能具有重要意义。

    \section{本论文主要研究内容}
    1.3 本论文主要研究内容
    针对当前异构星座组网与路由机制所面临的挑战，本文依托包含 LEO 感知层、LEO 通信层、MEO 中继层及 GEO 管控层 的四层深度融合网络场景，围绕架构设计、算法优化、平台验证的研究主线，开展以下核心工作：
    
    (1) 设计基于 SDN 的异构星座分级多域协同组网架构。 针对未来通导遥一体化网络中不同轨道层的功能差异，本文构建了以 GEO 为全局管控层、MEO 为跨域中继层、LEO 感知星座为接入与数据层的分级架构。提出了一种GEO 全局主控与LEO 域内协同的 管控模式：在垂直维度上，利用 GEO的广域覆盖实现跨层控制指令下发；在水平维度上，针对大规模 LEO 节点，设计了基于地理网格的 虚拟节点 构建策略与 分域自治 机制。该架构通过物理节点与逻辑身份的解耦，有效屏蔽了底层 LEO 卫星高速运动带来的拓扑剧变，解决了异构星座环境下的控制信令风暴问题。
    
    (2) 提出面向多域动态 QoS 的异构星座协同路由机制。 为解决多层网络中跨域路径规划复杂及局部链路易拥塞的难题，本文设计了一套分层协同路由策略。在域间层面，GEO 控制器基于全局抽象拓扑进行宏观路径规划与 QoS 预算分解；在域内层面，提出了一种 基于方向引导的改进遗传算法。该算法结合类曼哈顿网络的拓扑规律，引入 方向引导初始化策略 以克服标准算法搜索盲目性大的缺陷，并设计了包含 非线性拥塞惩罚项 的适应度函数，在保障端到端时延、带宽等多维 QoS 约束的同时，实现对网络热点链路的主动规避与负载均衡。
    
    (3) 构建基于容器虚拟化技术的异构星座网络仿真平台。 为验证所提架构与协议在真实网络协议栈下的性能，本文设计并实现了一个 高保真半实物仿真平台。该平台不同于传统的数值模拟，采用 STK卫星建模工具作为物理层驱动，生成高精度的卫星轨道与链路通断数据；利用 Docker 容器技术 模拟卫星节点与地面终端的 Linux 网络协议栈，通过 Open vSwitch  与 Linux TC 模块构建可编程的虚拟交换网络。在此基础上，开发了 SDN 控制器接口，实现了从物理场景驱动到路由流表下发的完整闭环，真实模拟了网络数据包在动态拓扑下的封装、转发与排队行为。
    
    \section{本论文组织结构}

    本论文共分为六章，各章节的具体内容安排如下：

    第一章为绪论。 阐述了在通导遥一体化背景下，异构星座组网架构与路由机制的研究背景及意义；系统梳理了国内外在异构星座管控架构、跨层协同路由技术等方面的研究现状，分析了现有研究的局限性；明确了本文的主要研究内容、创新点及篇章结构。
    
    第二章为相关技术背景介绍。 重点介绍了卫星轨道参数与分类、星间链路特性及典型的星座构型；阐述了软件定义网络（SDN）的核心架构及其在卫星网络中的应用模式；综述了卫星网络路由技术的基本原理与分类，为后续架构与算法设计奠定理论基础。
    
    第三章为异构星座组网与管控架构研究。 针对异构星座的时变拓扑与大规模管控难题，提出了一种基于 SDN 的分级多域协同管控架构。详细论述了“GEO 全局主控 + LEO 域内协同”的分层机制、基于地理网格的虚拟节点构建策略以及物理-逻辑地址映射方法，并设计了基于虚拟地址封装的异构网络互联机制。
    
    第四章为异构星座多域协同路由策略设计。 针对多维 QoS 约束下的选路问题，构建了基于虚拟节点的分层图模型。设计了包含全局规划与局部寻优的协同路由框架：在域间层面，制定了基于全局抽象拓扑的路径规划与 QoS 预算分解策略；在域内层面，提出了一种基于方向引导与拥塞感知的改进遗传算法（H-IGA），并对其核心算子与收敛性能进行了详细设计与分析。
    
    第五章为基于容器虚拟化技术的网络仿真平台设计与实现。 针对异构星座网络协议验证的高保真需求，设计并实现了一个基于 STK 与 Docker 的虚拟化网络仿真平台。重点阐述了物理场景数据与虚拟网络载体的动态映射机制、SDN 控制器与底层 OVS 虚拟交换机的交互流程；并在此仿真环境中，对所提管控架构与路由机制进行了功能验证与系统级性能评估。
    
    第六章为总结与展望。 对全文的研究工作与核心结论进行了总结，分析了当前研究存在的不足，并对异构星座组网技术的未来发展方向进行了展望。
    

    \chapter{相关技术背景介绍}
    \section{卫星网络组网架构}
    随着空间技术的快速迭代，卫星通信网络正经历从单一轨道向多层异构、从简单中继向宽带互联网演进的深刻变革。大规模低轨星座的部署使得网络节点数量呈指数级增长，传统的组网模式难以适应日益复杂的网络环境。本节将从大规模低轨星座的发展现状、异构星座的组网形态以及网络管控架构的演进三个维度进行阐述。
    
    低地球轨道（LEO）卫星凭借其低传输时延、低路径损耗以及低制造成本等显著优势，已成为构建下一代空间信息网络的核心载体。相较于传统的地球静止轨道（GEO）卫星，低轨卫星能够有效解决高纬度地区的覆盖盲区问题，并提供与地面光纤网络相媲美的传输体验。近年来，全球范围内掀起了大规模低轨卫星星座的建设热潮，卫星网络呈现出巨型化、宽带化和低轨化的发展特征。在国外，以美国 SpaceX 公司的 Starlink 计划和英国的 OneWeb 计划为代表，商业航天企业正在加速部署由数千甚至上万颗卫星组成的巨型星座。Starlink 计划通过在轨运行海量卫星，利用高密度的卫星覆盖实现全球宽带接入\citess{hui2025review}，OneWeb 则致力于构建数百颗卫星的星座系统，重点服务于高通量低时延的通信需求\citess{handley2023new}。在国内，为了满足国家战略需求并抢占轨道频谱资源，我国也规划了包括 GW-A59 和 GW-2 等子星座在内的庞大星座计划，规划卫星总数达到万颗量级\citess{TDYT202303006}。这些大规模星座的出现，导致网络空间中的节点数量激增，拓扑结构呈现出高度的动态性和复杂性，对现有的组网架构和管控技术提出了严峻挑战。

    为了应对大规模单层 LEO 网络在路由收敛和拓扑管理上的复杂性，研究人员提出了利用不同轨道高度卫星进行分层组网的异构架构。这种架构利用高轨卫星覆盖范围广、相对静止的特点来协助管理高动态的低轨卫星。文献\cite{lee2000satellite}率先提出了多层卫星网络（Satellite over Satellite, SoS）的概念，构建了一种层级化的网络架构，通过高层卫星作为管理节点来降低底层卫星的组网复杂度。在此基础上，文献\cite{akyildiz2002mlsr}提出了针对多层卫星 IP 网络的 MLSR 路由算法，进一步验证了多层架构在提升网络性能方面的优势。文献\cite{chen2005routing}则针对 LEO/MEO 双层架构设计了分层路由协议，利用 MEO 卫星对 LEO 卫星进行分组管理，有效缓解了大规模网络中的路由更新压力。这些早期的异构组网研究为后续的多层星座管控奠定了基础。

    在早期的卫星通信系统中，组网架构通常采用分布式管控模式。卫星节点不仅需要承担业务数据的存储与转发，还需负责链路状态信息的收集、同步以及路由表的计算\citess{JSTY201702007}。然而，面对节点数量成千上万的大规模卫星网络，这种分布式架构暴露出明显的局限性：首先，星上计算资源和存储资源极为有限，难以支撑全网路由计算的高额开销。其次，由于卫星拓扑的高动态变化，全网链路状态信息的频繁泛洪会导致严重的信令风暴，使得路由收敛速度缓慢，难以保障业务的服务质量。最后，软硬件紧耦合的特性使得网络协议升级困难，缺乏灵活性。为了解决分布式架构的痛点，软件定义网络转控分离的思想被引入卫星网络，组网架构随之向集中化、智能化方向演进，并经历了从单层管控向多层分域管控的发展过程。

    将SDN架构引入卫星网络，能够有效应对其固有的挑战，并带来显著优势\citess{jiang2023software}。首先，卫星网络拓扑时变、星间链路间歇连通，而SDN的集中控制能力使其能够基于全局视角进行路由计算与资源分配，从而提升网络管理效率。其次，卫星硬件部署后难以升级，SDN控制与转发分离的特性，以及软硬件解耦的设计，使得通过网络远程进行软件维护与功能更新成为可能，极大增强了网络的灵活性与可演进性。最后，面对由不同轨道、不同类型卫星及地面网络构成的异构融合网络，SDN提供了统一的编程接口与控制平面，能够实现跨域的协同资源调度与高效管理。
    
    SDN 是一种新型网络架构范式，其核心思想是将传统网络设备中紧密耦合的控制平面与数据转发平面进行解耦，通过集中式的逻辑控制器对全网资源进行统一调度与编程管理\citess{bholebawa2018performance}。如图\ref{fig: SDN结构}所示，SDN 架构通常分为三层：顶层为应用层，运行各类网络服务，如负载均衡、安全防护、移动性管理等；中间为控制层，包含一个或多个 SDN 控制器，负责维护全网的全局拓扑视图，并根据应用请求执行路由计算、策略部署和链路状态监控等；底层为由交换机、路由器等网络设备组成的基础设施层，负责匹配策略、转发数据包等。其中，北向接口位于控制器层与应用层之间，其主要功能是为上层网络应用提供对底层网络资源的抽象访问能力，各类网络服务可通过北向接口向 SDN 控制器请求网络状态信息或下发控制指令。当前主流的北向接口多采用RESTful API\citess{richardson2008restful}形式，以及基于领域特定语言的高级编程框架，如Frenetic\citess{foster2011frenetic}、Pyretic\citess{monsanto2013composing}等；而南向接口则连接控制器层与基础设施层，负责将控制器生成的转发规则精确地下发至数据平面设备，并实时采集链路状态、队列负载等运行时信息，目前常用的协议有OpenFlow\citess{mckeown2008openflow}、ForCES\citess{doria2010forces}、OVSDB\citess{pfaff2013open}等。
    \begin{figure}[!htb]
        \centering
        \includegraphics[width=0.7\linewidth]{fig/SDN结构.pdf}
        \caption{软件定义网络三层架构}
        \label{fig: SDN结构}
    \end{figure}
    
    目前，根据控制平面与数据平面中卫星类型的组合方式，可主要归纳为单层、双层与三层结构\citess{jiang2023software}。
     \begin{figure}[!htb]
      \centering
      \subfloat[]{
        \includegraphics[width=0.3\linewidth]{fig/SDN卫星单层.pdf}
        \label{fig: SDN卫星单层}
      }
      \hfill
      \subfloat[]{
        \includegraphics[width=0.3\linewidth]{fig/SDN卫星双层.pdf}
        \label{fig: SDN卫星双层}
      }
      \hfill
      \subfloat[]{
        \includegraphics[width=0.3\linewidth]{fig/SDN卫星三层.pdf}
        \label{fig: SDN卫星三层}
      }
      \caption[软件定义卫星网络结构]{软件定义卫星网络结构。\subref{fig: SDN卫星单层}单层架构；\subref{fig: SDN卫星双层}双层结构；\subref{fig: SDN卫星三层}三层结构}
      \label{fig: 软件定义卫星网络架构}
    \end{figure}

    单层结构是最基础的形态，如图\ref{fig: SDN卫星单层}所示。通常仅由单一类型的卫星（以低轨LEO卫星最为常见）构成数据平面。其控制平面可部署于地面站，如文献\cite{wang2018design}通过地面数据中心管理LEO星座资源；亦可部署于LEO卫星本身，构成带内控制模式，如文献\cite{hu2018leo}及文献\cite{papa2020design}中考虑动态控制器迁移与重配置开销的LEO星上控制器方案，此种方式虽能降低传输时延，但需应对因卫星高速运动导致的控制器动态迁移与拓扑同步开销。

    双层结构则融合了两种不同类型的卫星，如图\ref{fig: SDN卫星双层}所示。最常见的组合是以覆盖广、位置相对静止的高轨GEO卫星作为控制平面，管理作为数据平面的LEO卫星\citess{xiaodong2020multi}。此种架构充分利用了GEO卫星的广域视野与链路稳定性，支持多QoS目标的自适应路由算法。此外，控制平面亦可设于地面站，同时管理GEO与LEO数据平面，如FRBSN\citess{sheng2017toward}；或采用地面站为超级控制器、GEO为主控制器、部分LEO为从控制器的混合分层模式\citess{xu2018software}，以增强网络控制的鲁棒性与覆盖完整性。文献\cite{1021025190.nh}将地面中心作为一级控制器负责全局路由计算，三颗GEO卫星作为二级控制器协助控制信令中转与拓扑收集，而LEO卫星仅充当交换机负责数据转发。
    
    三层结构最为复杂，如图\ref{fig: SDN卫星三层}所示。同时引入了GEO、MEO和LEO卫星，旨在构建一个功能分层、协同运作的立体网络。文献\cite{bao2014opensan}框架以GEO卫星群作为控制平面，统一管理包含GEO、MEO、LEO在内的整个数据平面，显著减少地面站数量。在星地融合场景下，文献\cite{li2017service}将控制功能分布于GEO卫星、地面数据中心与卫星网关，由卫星网络管理中心进行统一协调，数据平面则囊括了MEO/LEO卫星及地面路由器、交换机等设备，以实现端到端的精细化服务质量保障。文献\cite{li2019icn}将主控制器置于GEO，辅助控制器置于MEO以分担负载，LEO则专注于数据转发，并通过协同缓存等技术优化内容检索性能。此外，还有文献\cite{feng2014scheme}、文献\cite{feng2017hetnet}等架构将控制器部署于地面，统一调度多层卫星资源，支持异构网络融合与移动性管理。
    
    \section{卫星网络路由算法}
    1、单层卫星网络路由

    随着低轨卫星网络的快速发展，路由技术作为实现星间数据转发的核心机制，面临着卫星节点高动态性、网络拓扑频繁变化、星座规模不断扩大等挑战。传统的地面网络路由协议难以直接适用于卫星网络环境。目前，卫星网络路由技术主要可分为基于时间虚拟化和基于位置虚拟化两大类别\citess{吕若晨2022卫星网络路由技术研究}，以下对各类型的代表性方法进行分析。
    
    基于时间虚拟化的路由技术是早期卫星网络路由研究的重要方向，其核心思想是利用卫星轨道运动的可预测性，将系统周期划分为离散的时间片，在每个时间片内将动态拓扑视为静态进行路由计算，这类技术主要解决了卫星高速运动带来的拓扑频繁变化问题。文献\cite{werner2002dynamic}首次提出将卫星网络运行周期平均划分为多个持续时间相等的时段，每个时段称为一个拓扑快照。在该时段内，网络拓扑被视为固定不变，并基于此静态拓扑预计算路由表。该方法实现简单、易于管理。文献\cite{王京林2009LEO}对此进行了优化，提出通过删除维持时间较短的星间链路来延长快照的持续时间，进一步提升了系统效率。然而，固定的时间划分可能导致快照切换时刻与实际拓扑变化时刻不同步，造成路由状态与网络实际状态的不匹配。为解决等时间间隔划分的不足，文献\cite{gounder1999routing}中提出了基于星间链路连通性变化的快照划分方法。该方法不再依赖于固定的时间间隔，而是以星间链路的建立或断开作为拓扑变化的触发条件，将系统周期划分为持续时间不等的拓扑快照。这种划分方式更贴近网络拓扑的实际变化，能够更精确地描述网络的动态特性。
    然而，该方法在提升精度的同时，需要实时监测链路状态，计算复杂度较高，且快照持续时间的不确定性增加了系统管理的难度。
    
    总体而言，基于时间虚拟化的路由技术通过将动态拓扑转化为系列静态拓扑，有效简化了路由计算过程，通常由地面控制中心进行路由预计算并上注至卫星网络。然而，随着星座规模的扩大，全网路由表的存储和更新将带来巨大的开销，且静态的路由计算难以适应网络的动态拥塞变化，这些局限性促使研究者转向其他技术路线。为克服时间虚拟化技术的局限性，基于位置虚拟化的路由技术从空间维度提出了新的解决方案。该类技术的基本思路是通过将移动的卫星与固定的地理区域或虚拟节点建立映射关系，实现路由计算的简化，从而在空间上屏蔽卫星网络的动态性。
    
    文献\cite{ekici2000datagram}中提出了基于地面覆盖区域虚拟化的路由方案。该方法按照经纬度将地球表面划分为若干逻辑区域，每个区域被赋予唯一的逻辑地址。在卫星运行过程中，距离区域中心最近的卫星被动态指定为该区域的接入卫星，并继承相应的逻辑地址。这种方法的核心优势在于：路由计算仅需考虑源和目的端所在的固定逻辑区域，而无需关注具体服务卫星的变化，从而有效屏蔽了卫星的运动性。这些方法充分利用了星座拓扑的规则性和周期性，但当地面用户接入卫星切换时，需要进行路由信息的转移，可能引发额外的信令开销，在频繁切换场景下性能会受到一定影响。
    
    为进一步优化系统性能，在后续研究\cite{ekici2002distributed}中提出了数据报路由算法(Datagram Routing Algorithm，DRA)。该方法不再将卫星映射到地面区域，而是将LEO卫星轨道球面划分为若干个虚拟区域，每个区域对应一个虚拟卫星节点，构成固定的虚拟拓扑结构。在该算法中，实际卫星根据其当前位置与虚拟节点建立动态映射关系，路由计算在虚拟拓扑上进行，计算结果再映射回实际卫星。DRA算法设计简单，仅需利用卫星位置信息即可完成路由决策，且LEO卫星间无需进行动态信息交互，大大降低了系统复杂度。然而，该方法也缺乏对网络实时状态的感知能力，难以应对网络拥塞问题，且其路由计算依赖于极轨星座的对称性，不适用于倾斜轨道星座，这限制了其在现代异构星座中的适用性。
    
    基于位置虚拟化的路由技术通过空间映射有效降低了路由计算的复杂度，但其性能在很大程度上依赖于星座结构的规则性，且通常需要预存卫星位置信息。随着星座规模的扩大和轨道类型的多样化，这些方法的适应性和扩展性面临严峻挑战。

    2、多层卫星网络路由
    
    随着未来空间网络向多层异构架构演进，单纯依靠 LEO 层的路由机制已难以满足全球覆盖与多维资源调度的需求。利用 GEO/MEO 卫星的广域覆盖与稳定性优势，辅助底层 LEO 卫星进行分层协同路由成为主流研究方向。

    针对双层或三层卫星网络架构，分层管理路由协议被广泛采用。文献\cite{lee2000satellite}提出的分层卫星路由协议（Hierarchical Satellite Routing Protocol，HSRP），该协议利用MEO层汇聚LEO局部拓扑并下发全网状态，由LEO卫星自行计算路由。该机制减轻了LEO卫星的拓扑维护压力，但MEO层的频繁交互也引入了显著的控制开销与传播时延。文献\cite{akyildiz2002mlsr}针对三层卫星网络提出了多层卫星路由协议（Multi Layer Satellite Routing，MLSR），引入分组抽象思想构建了分级管理架构。该协议采用自顶向下的逐层路由计算策略，由GEO和MEO卫星分别作为管理者计算下层网络的组间路由，通过这种分层分组机制显著提升了系统在大规模网络中的可扩展性。文献\cite{zhou2007novel} 则进一步结合了虚拟节点与分层管理的优势，提出了一种混合路由策略，利用 MEO 卫星维护 LEO 层的虚拟拓扑状态，在降低控制开销的同时提升了路由的稳定性。
    
    文献\cite{long2010sustainable} 等尝试将遗传算法、蚁群算法等智能优化方法应用于多层卫星网络，旨在异构拓扑中搜索满足时延与带宽约束的最优路径。尽管如此，现有算法在面对异构星座的高动态拓扑时，往往面临收敛速度慢、局部拥塞规避能力不足等问题。特别是当 LEO 层出现局部热点流量时，如何利用 MEO/GEO 层进行高效的流量卸载与负载均衡，仍是异构星座路由机制亟待解决的关键问题。文献\cite{周云晖2008卫星网络}提出了一种融合分组管理与时间分片策略的时分路由协议（TDRP）。该协议建立了自下而上的状态汇聚与自上而下的路由分发流程：底层卫星将链路状态逐层上报至MEO与GEO管理者进行域划分及拓扑同步，最终由GEO层基于全局信息统一计算域间路由，并通过分层下发机制完成全网路由表的协同更新。

    针对传统分组式路由机制在网络灵活性与资源全局调度上的局限性，近年来研究者开始尝试将SDN架构引入卫星网络，以期通过控制与转发分离的思想突破现有瓶颈。
    文献\cite{8258968}提出了一种名为SoftSpace的下一代卫星网络软件定义架构，并在该框架下开发了面向体验质量的空间路由机制。该方案利用SDN的逻辑集中与可编程特性，根据用户体验需求动态优化路径选择，旨在解决传统网络难以提供差异化服务及灵活性不足的问题。文献\cite{li2017service}同样利用SDN技术将控制平面集中于GEO卫星，利用其全局视野提出了一种面向多维QoS的路由算法辅以带宽分配策略，为不同优先级业务动态计算最优路径，并引入调整因子以有效规避瓶颈链路，在链路拥塞时优先保障高优先级业务的资源需求。文献\cite{1021025190.nh}提出了一种基于SDN架构的业务优先级路由算法,利用地面与GEO卫星协同管控全局拓扑,实现路由计算与数据转发分离。算法结合虚拟节点模型屏蔽LEO卫星的动态特性，并制定了差异化的拥塞规避策略，确保高优先级业务在链路负载过高时沿最短路径传输，而调度次优先级业务切换至迂回路径，从而在有限的星上资源下有效保障了关键业务的低时延传输需求。文献\cite{9441632}提出了一种基于软件定义LEO-GEO双层卫星架构的QoS感知路由算法（QoSRA）。该算法通过剔除计算过程中的迭代步骤对低复杂度路由算法（LCRA）进行了优化，在降低星上存储开销的同时，利用高优先级队列机制保障了双向会话服务的传输质量，从而实现了优于传统方法的低时延与高吞吐量性能。

    综上所述，卫星网络路由技术正经历从单层简单拓扑向多层异构架构、从刚性分布式控制向软件定义灵活管控演进的趋势。虽然传统的分层路由协议在一定程度上缓解了网络扩展性问题，但面对日益复杂的星上环境与差异化的QoS需求，基于SDN架构的路由机制凭借其控制与转发分离、全局资源感知及灵活调度的优势，展现出更优的适应性与性能潜力，已成为提升未来大规模卫星网络综合服务能力的关键技术路径。
    
    \section{卫星网络仿真平台}
    大规模星座网络具有节点数量庞大、拓扑结构高度动态变化以及链路状态时变等特征，这使得构建真实的物理测试环境不仅成本极高，而且难以复现各种极端的网络场景\citess{jones2018recent}。因此，利用网络仿真平台对星座组网协议、路由算法及网络功能进行验证与评估，成为了相关技术研究中不可或缺的重要手段。目前的卫星网络仿真技术主要分为硬件实物测试、离散事件模拟以及虚拟化仿真三类\citess{1025403180.nh}。硬件实物测试平台通过构建真实的卫星通信硬件组件来模拟网络环境，文献\cite{ivancic2010experience}中研究人员曾将监测星座的在轨卫星作为测试平台，验证了延迟容忍网络协议，文献\cite{komnios2014spice}提出SPICE测试床，通过专业硬件精确模拟卫星通信真实损伤环境的地面实验室平台，用于验证延迟容忍网络协议与算法。然而，这类平台受限于高昂的成本和复杂的环境搭建，难以满足大规模组网的仿真需求。
    
    相比之下，离散事件模拟技术核心在于通过软件模型来抽象和再现网络设备与协议的行为。这类平台通常采用离散事件驱动机制，维护一个按时间排序的事件队列来推进仿真进程。
    代表性工具包括 OPNET、NS3 和 OMNET++ 等。OPNET\citess{chang1999network} 是一款功能全面的商业级网络仿真开发环境，它提供了直观的图形化用户界面和丰富的标准协议库，支持采用分层结构对有线及无线网络组件的行为进行精细建模与性能评估。NS3\citess{riley2010ns} 作为一款广泛使用的开源离散事件网络模拟器，相比前代 NS2\citess{issariyakul2008introduction} 在扩展性和代码设计上有了显著提升，它支持通过 C++ 或 Python 脚本构建网络拓扑，具备优秀的跨层仿真能力，并支持读取和解析 Wireshark 等真实数据包分析工具的捕获文件。OMNET++\citess{varga2010omnet++} 则是一种基于组件的模块化离散事件仿真框架，它不直接提供预置的网络协议模型，而是提供基于 C++ 的基础类库和消息传递机制，供用户灵活构建自定义的简单模块或复合模型，非常适用于验证复杂的通信协议与分布式系统。然而，此类数据包级的仿真器也有明显的局限性，其网络协议栈通常是简化的数学模型而非真实的操作系统内核，这可能导致仿真结果与实际系统存在偏差。此外，这类平台难以直接承载真实的业务流量，限制了其在应用层功能验证方面的能力。

    为了兼顾仿真的真实性与灵活性，基于虚拟化技术的仿真平台逐渐成为研究热点。该类技术利用轻量级虚拟化或全系统虚拟化技术，在物理主机上运行真实的操作系统和网络协议栈。其中，容器技术作为轻量级虚拟化的代表，通过将应用程序、依赖库及配置文件打包为独立的运行单元，实现了应用与底层环境的解耦\citess{aqasizade2025experimental}。Docker 是当前最主流的容器化平台，如图\ref{fig: 传统虚拟化与容器虚拟化对比}它引入了 Docker 引擎来替代传统虚拟机中的虚拟机监视器。与虚拟机需要运行完整的客户操作系统不同，Docker 容器直接共享宿主机的操作系统内核，在进程级进行资源隔离。
    \begin{figure}[!htb]
		\centering
		\subfloat[]{
			\includegraphics[width=0.3\linewidth]{fig/虚拟机与容器的对比图-传统虚拟化.pdf}
			\label{fig: 传统虚拟化}
		}
		% \hfill
		\subfloat[]{
			\includegraphics[width=0.3\linewidth]{fig/虚拟机与容器的对比图-容器虚拟化.pdf}
			\label{fig: 容器虚拟化}
		}
		\caption[]{传统虚拟化与容器虚拟化对比。\subref{fig: 传统虚拟化}传统虚拟化；\subref{fig: 容器虚拟化}容器虚拟化} 
        \label{fig: 传统虚拟化与容器虚拟化对比}
	\end{figure}
 
    表\ref{tab: 两种虚拟化方式}进一步详细对比了这两种虚拟化方式在启动速度、资源开销等关键维度上的差异。
    \begin{table}[htbp]
      \centering
      \caption{虚拟机与 Docker 容器对比}
      \label{tab: 两种虚拟化方式}
      \begin{tabular}{l c c}
        \toprule
        \textbf{对比维度} & \textbf{虚拟机 (VM)} & \textbf{Docker 容器} \\
        \midrule
        启动速度 & 分钟级  & 秒级  \\
        镜像大小 & GB级 (包含完整操作系统) & MB级 (仅应用+依赖) \\
        性能开销 & 较高 (需虚拟化CPU/内存) & 接近原生 (共享内核) \\
        隔离级别 & 强隔离 (硬件虚拟化) & 进程级隔离 (命名空间) \\
        资源占用 & 每台VM占用独立内核 & 共享Host OS内核 \\
        系统支持 & 可运行不同操作系统 & 依赖Host系统内核 \\
        可移植性 & 需迁移整个系统镜像 & 一次构建，到处运行 \\
        管理复杂度 & 需管理OS补丁、更新 & 只管理应用层 \\
        \bottomrule
      \end{tabular}%
    \end{table}
    
    这种架构设计使得容器具有显著的性能优势，启动时间可从分钟级缩短至秒级，资源占用从 GB 级降低至 MB 级，且单台物理主机可支撑的节点数量远超虚拟机。这些特性使得基于 Docker 的虚拟化方案能够以较低的硬件成本实现大规模、高保真的网络拓扑构建，非常契合大规模星座网络仿真的需求。目前，已有多种仿真工具采用了这一技术路线。例如，Mininet 利用 Linux 网络命名空间实现了轻量级网络仿真\citess{kaur2014mininet}，Containernet在此基础上集成了 Docker 容器技术以支持复杂应用部署\citess{peuster2018containernet}，EmuStack 则基于 OpenStack 和 SDN 架构设计\citess{sturmberg2016emustack}。
    
    为了在物理主机内部构建出逼真的卫星网络拓扑，并实现基于 SDN 的灵活管控，除了基础的容器虚拟化技术外，还需要依赖虚拟网络设备接口、虚拟交换机以及控制器等关键组件的协同工作。Docker 提供了 Host、None 和 Bridge 等多种网络模式，其中 桥接模式（Bridge Mode） 是实现容器互联的核心机制。如图 \ref{fig: 容器互联桥接模式} 所示，在此模式下，宿主机上创建一个虚拟网桥（可以是默认的 Docker0 网桥，也可以是自定义的 OVS 网桥），容器通过 Veth-pair 挂载至该网桥上，从而在逻辑上组成一个二层局域网 。
    \begin{figure}[!htb]
        \centering
        \includegraphics[width=0.6\linewidth]{fig/容器桥接模式.pdf}
        \caption{容器互联桥接模式}
        \label{fig: 容器互联桥接模式}
    \end{figure}

    其中，Veth-pair 是 Linux 内核提供的一对虚拟以太网设备，它们总是成对出现，如同两端连接的虚拟网线，一端连接容器独立的网络命名空间，另一端连接宿主机的虚拟交换机，从而打通了容器与外部网络的通信链路。为了实现复杂的网络转发逻辑，Open vSwitch (OVS) 被广泛用作虚拟交换机。与传统的 Linux Bridge 相比，OVS 是一款专为虚拟化环境设计的开源多层交换机，支持标准的 OpenFlow 协议。OVS 采用内核态数据路径与用户态守护进程分离的架构，在保障数据包高速转发的同时，允许控制器通过流表下发精细的转发规则，从而实现链路通断控制、QoS 保障及流量监控等高级功能。图\ref{fig: OVS处理包的过程}展示了OVS的包处理过程，当网卡接收到一个数据包后，该包进入 OVS 的数据路径。
    \begin{figure}[!htb]
        \centering
        \includegraphics[width=0.6\linewidth]{fig/OVS处理包的过程.pdf}
        \caption{OVS处理包的过程}
        \label{fig: OVS处理包的过程}
    \end{figure}
    
    如果数据路径在缓存中成功匹配到对应的流表项，则直接执行流表项中定义的动作。反之，如果发生流表缺失，数据路径会将该数据包上送至 ovs-vswitchd，由其进行进一步的流表匹配\citess{188960}。一旦 ovs-vswitchd 确定了该数据包的处理方式，便会将数据包连同处理指令一并返回给数据路径执行。通常情况下，它还会指示数据路径将这条流表项缓存起来，以便后续快速处理同类数据包。
    
    
 %    \section{卫星通信网络}
 %    \subsection{卫星轨道参数与分类}
 %    如图\ref{fig: 卫星开普勒轨道参数}，地心惯性坐标系（Earth-Centered Inertial, ECI）是轨道力学中描述地球卫星运动最常用的参考系之一。该坐标系以地球质心为原点，X轴指向参考历元的春分点方向，Z轴沿地球自转轴指向北极，Y轴则按右手定则与X、Z轴构成右手坐标系\citess{刘洋2007地心惯性坐标系到质心轨道坐标系的坐标转换方法}。
    
 %    卫星轨道参数用于定量描述卫星在空间中的运动轨迹，是计算卫星位置的基础。可用典型的开普勒轨道六根数唯一确定卫星轨道的形状、大小、空间方向及其在轨道上的瞬时位置\citess{ruth2024master}。长半轴 $a$ 和偏心率 $e$ 描述轨道的形状与大小，直接决定卫星的覆盖范围与运行周期；轨道倾角 $i$ 和升交点赤经 $\Omega$ 定义轨道面的空间定向，决定了星座的几何构型（如 Walker 星座的轨道面分布）；近地点角距 $\omega$ 与真近点角 $V$ 则描述卫星在轨道内的定向与瞬时位置。这些参数共同构成了对卫星轨道全面且精确的描述。

 %    \begin{figure}[!htb]
 %        \centering
 %        \includegraphics[width=0.6\linewidth]{fig/卫星开普勒轨道参数.pdf}
 %        \caption{卫星开普勒轨道参数}
 %        \label{fig: 卫星开普勒轨道参数}
 %    \end{figure}
    
 %    卫星轨道能决定系统的覆盖范围、传输时延及拓扑复杂度。基于不同特征，通常从高度、周期及倾角三个维度进行分类。
    
 %    （1）按轨道高度分类

 %    如图 \ref{fig: 卫星轨道分布} 所示，受大气阻力与范·艾伦辐射带（即地球磁场捕获的高能粒子聚集区，易导致星载电子设备损伤）等物理环境制约\citess{郝选文0空间信息网抗毁路由及网络防攻击技术研究}，通信卫星根据轨道高度的不同主要分为三类\citess{徐双2019软件定义卫星网络关键技术研究}：
 
 %    低地球轨道 (LEO, 500--1500 km)：位于内辐射带下方。单星覆盖小且相对运动快，导致网络拓扑高度动态，需巨型星座实现全球覆盖，但具备低时延、低路损优势。
    
 %    中地球轨道 (MEO, 5000--20000 km)：位于内、外辐射带之间的安全窗口。特性介于 LEO 与 GEO 之间，常用于导航（如 GPS）及区域中继。
   
 %    地球静止轨道 (GEO, $\approx$ 36000 km)：位于赤道上空。单星可覆盖约 1/3 地球且相对静止，适合广播业务；但存在两极盲区，且传输时延大，需具备抗磁层扰动能力。
 %    % [填充填充填充填充填充填充填充填充填充填充填充填充填充填充填充填充填充填充填充填充填充填充填充填充填充填充填充填充填充填充填充填充填充填充填充填充填充填充填充填充填充填充填充填充填充填充填充填充填充填充填充填充填充填充填充填充填充填充填充填充填充填充填充填充填充填充填充填充填充填充填充填充填充填充填充填充填充填充填充填充填充]
 %    \begin{figure}[!htb]
 %        \centering
 %        \includegraphics[width=0.6\linewidth]{fig/卫星轨道分类}
 %        \caption{卫星轨道分类}
 %        \label{fig: 卫星轨道分布}
 %    \end{figure}

 %    （2）按轨道周期分类\citess{vallado2001fundamentals}
    
 %    根据卫星绕地球运行一周所需时间（即轨道周期 T），可将轨道划分为：
    
 %    亚同步轨道：轨道周期小于地球自转周期，包括 LEO 与 MEO，适用于需频繁重访或低时延通信的任务；
    
 %    同步轨道：轨道周期等于地球自转周期，其中若轨道倾角为 0° 且偏心率为 0，则称为地球静止轨道（GEO），卫星在地面观察者视角下静止于赤道上空某一点；
    
 %    超同步轨道：轨道周期大于地球自转周期，多用于轨道转移或特殊科学探测任务，实际运行卫星较少。
    
 %    （3）按轨道倾角分类\citess{larson1999space}
    
 %    轨道倾角 i 定义为轨道平面与地球赤道面之间的夹角（$0^\circ \le i \le 180^\circ$），据此可分为：
    
 %    赤道轨道：$i=0^\circ$，轨道位于赤道平面内;
    
 %    倾斜轨道：$0^\circ < i < 90^\circ$，覆盖范围主要集中于中低纬度，广泛用于通信与导航星座；
    
 %    极轨道：$i\approx 90^\circ$，卫星每圈经过或靠近地球两极，结合地球自转可实现全球覆盖，适用于遥感、气象及高纬通信。

 %    \subsection{卫星链路分类}
 %    根据链路连接的拓扑关系，星间链路（Inter-Satellite Link, ISL）是实现自主组网通常可分为三类：轨道内星间链路、轨道间星间链路以及层间星间链路。
    
 %    （1）轨道内星间链路（Intra-Plane ISL）
    
 %    指同一轨道面（即具有相同轨道高度、倾角和升交点赤经的一组卫星）内相邻卫星之间的通信链路。由于同轨卫星相对位置稳定，此类链路易于建立并维持高可靠性连接，常用于数据在轨道面内的纵向传递。例如，在Walker Delta星座构型\citess{1971Walker}中，每颗卫星通常与前后两颗同轨卫星建立双向Intra-Plane ISL，构成环状拓扑。
    
 %    （2）轨道间星间链路（Inter-Plane ISL）
    
 %    指不同轨道面之间、但处于相近纬度或交汇区域的卫星所建立的链路。由于异轨卫星存在显著的相对运动，链路持续时间较短、几何关系动态变化剧烈，但其是实现跨轨道面数据交换与全球连通性的核心手段。
    
 %    （3）层间星间链路（Inter-Orbit Link, IOL）
    
 %    指部署于不同轨道高度层之间的星间链路，属于多层卫星网络中的跨层连接机制。例如，IOL可用于将LEO用户接入流量汇聚至GEO节点，或由MEO导航卫星向LEO通信卫星提供增强信息。
    
 %    上述三类链路共同构成了卫星网络的空间拓扑。其中，Intra-Plane与Inter-Plane ISL主要支撑单层LEO星座的平面内与平面间连通，而IOL则为多层异构架构提供垂直协同能力。
    
 %    \subsection{卫星星座构型}
 %    现代卫星通信、导航与遥感系统广泛依赖多星协同工作。为实现全球或区域连续覆盖，需合理设计卫星星座构型。其中，Walker倾斜轨道星座 和 极地轨道星座 是两类经典且应用广泛的构型。前者通过数学化参数实现均匀覆盖，后者则凭借高倾角轨道实现对地球两极的观测能力。
    
 %    （1）Walker倾斜轨道星座
    
 %    % Walker倾斜轨道星座是一种基于圆形轨道、等间隔升交点赤经、均匀相位分布的规则星座构型\citess{1971Walker}，如图所示\ref{fig: 倾斜轨道星座示意图}。所有卫星具有相同轨道高度和倾角，但分布在多个轨道平面中。核心参数由$N/P/F$表示，其中N为卫星总数，P为轨道平面数量，F为相位因子（Phasing Factor），表示相邻轨道平面间卫星的相位偏移量（单位：卫星间隔数）。每个轨道平面含 $S = N/P$ 颗卫星，相邻轨道平面升交点赤经差为 $\Delta\Omega = {360^\circ}/P$，相邻平面内首颗卫星的平近点角差为 $\Delta M = 360^\circ \times F/S$。
 %    % 典型卫星系统中，Globalstar（48/8/1，1,414 km，52°）侧重中低纬度区域的语音与数据服务；Galileo（27/3/1，23,222 km，56°）作为中地球轨道导航系统，提供高精度定位。这些系统均基于Walker Delta构型，在轨道高度、倾角与规模上针对各自任务进行了优化。

 %    Walker倾斜轨道星座是一种基于圆形轨道、等间隔升交点赤经、均匀相位分布的规则星座构型\citess{1971Walker}，如图所示\ref{fig: 倾斜轨道星座示意图}。其几何特征由参数组 $N/P/F$ 唯一确定，其中 $N$ 为卫星总数，$P$ 为轨道平面数，$F$ 为相位因子。在该构型中，相邻轨道面的升交点赤经差为 $\Delta\Omega = 360^\circ/P$，相邻平面内卫星的平近点角差为 $\Delta M = 360^\circ \times F/S$（$S=N/P$）。该构型在中低纬度具有稳定的异轨链路几何特征，广泛应用于 Globalstar 通信系统及 Galileo 导航系统。
 
 %     （2）极地轨道星座
     
 %    如图\ref{fig: 极地轨道星座示意图}所示，极地轨道星座由倾角接近 90° 的轨道组成（如 Iridium 系统），卫星飞越两极，结合地球自转可实现全球无缝覆盖。然而，该构型在赤道区域存在拓扑缺陷：相邻轨道在赤道处间距最大，且存在相对运动方向相反的反向缝，增加了星间链路建链与路由切换的复杂度，若卫星密度不足易形成瞬时覆盖盲区。
 %    \begin{figure}[!htb]
	% 	\centering
	% 	\subfloat[]{
	% 		\includegraphics[width=0.4\linewidth]{fig/倾斜轨道星座_stk.png}
	% 		\label{fig: 倾斜轨道星座_stk}
	% 	}
	% 	\hfill
	% 	\subfloat[]{
	% 		\includegraphics[width=0.4\linewidth]{fig/倾斜轨道星座_俯视.pdf}
	% 		\label{fig: 倾斜轨道星座_俯视}
	% 	}
	% 	\caption[不同观察角度下的倾斜轨道星座的示意图]{不同观察角度下的倾斜轨道星座的示意图\subref{fig: 倾斜轨道星座_stk}赤道平面；\subref{fig: 倾斜轨道星座_俯视}极地} \label{fig: 倾斜轨道星座示意图}
	% \end{figure}

 %    （3）极轨与倾斜轨道结合的混合星座构型
    
 %    % 现代大型低轨卫星系统普遍采用极轨道与倾斜轨道相结合的混合星座构型，以在保障全球覆盖的同时，有效提升中低纬度人口密集区的服务容量与链路性能\citess{zhang2022leo}。
 %    % Starlink第二代系统采用混合多倾角轨道构型，由四个轨道层组成。极轨层运行于 560 km 高度，倾角为 97.6°（太阳同步轨道），部署 348 颗卫星，专门保障高纬度及极地地区（纬度 > ±60°）的连续覆盖，三个倾斜轨道层均位于约 550 km 高度，倾角分别为 43°、53.2° 和 70°，共部署 3,888 颗卫星，重点服务于北美、欧洲和亚洲等中低纬度人口密集区域。该设计通过极轨层彻底消除传统倾斜星座在极区的覆盖盲区，同时利用多倾角倾斜层显著提升赤道至中纬度的容量密度，从而实现“从赤道到两极”的高可用、高容量宽带接入。OneWeb初始星座采用 1,200 km 高度、87.9° 倾角的近极轨 Walker 构型（648 颗卫星），其后续规划引入约 50° 倾角的补充轨道层，旨在增强中纬度区域的服务容量。

 %    为了兼顾全球覆盖与热点区域容量，现代巨型星座（如 Starlink Gen2、OneWeb 后续规划）普遍采用极轨层与倾斜层相结合的混合构型\citess{zhang2022leo}。该设计利用极轨层消除高纬度及极区的覆盖盲区，同时利用多层倾斜轨道增加中低纬度人口密集区的卫星重叠度。通过极轨打底、倾斜增强的多层协同，实现了从赤道到两极的高可用、高容量接入。
    
 %    \begin{figure}[!htb]
	% 	\centering
	% 	\subfloat[]{
	% 		\includegraphics[width=0.4\linewidth]{fig/极地轨道星座_stk.png}
	% 		\label{fig: 极地轨道星座_stk}
	% 	}
	% 	\hfill
	% 	\subfloat[]{
	% 		\includegraphics[width=0.4\linewidth]{fig/极地轨道星座_俯视.pdf}
	% 		\label{fig: 极地轨道星座_俯视}
	% 	}
	% 	\caption[不同观察角度下的极地轨道星座的示意图]{不同观察角度下的极地轨道星座的示意图\subref{fig: 极地轨道星座_stk}赤道平面；\subref{fig: 极地轨道星座_俯视}极地} \label{fig: 极地轨道星座示意图}
	% \end{figure}


 %    \section{软件定义卫星网络}
 %    软件定义网络（Software-Defined Networking, SDN）是一种新型网络架构范式，其核心思想是将传统网络设备中紧密耦合的控制平面与数据转发平面进行解耦，通过集中式的逻辑控制器对全网资源进行统一调度与编程管理\citess{bholebawa2018performance}。如图\ref{fig: SDN结构}所示，SDN 架构通常分为三层：顶层为应用层，运行各类网络服务，如负载均衡、安全防护、移动性管理等；中间为控制层，包含一个或多个SDN控制器，负责维护全网的全局拓扑视图，并根据应用请求执行路由计算、策略部署和链路状态监控等；底层为由交换机、路由器等网络设备组成的基础设施层，负责匹配策略、转发数据包等。
 %    其中，北向接口位于控制器层与应用层之间，其主要功能是为上层网络应用提供对底层网络资源的抽象访问能力，各类网络服务可通过北向接口向SDN控制器请求网络状态信息或下发控制指令。当前主流的北向接口多采用RESTful API\citess{richardson2008restful}形式，以及基于领域特定语言的高级编程框架，如Frenetic\citess{foster2011frenetic}、Pyretic\citess{monsanto2013composing}等；而南向接口则连接控制器层与基础设施层，负责将控制器生成的转发规则精确地下发至数据平面设备，并实时采集链路状态、队列负载等运行时信息，目前常用的协议有OpenFlow\citess{mckeown2008openflow}、ForCES\citess{doria2010forwarding}、OVSDB\citess{pfaff2013open}等。

 %    将SDN架构引入卫星网络，能够有效应对其固有的挑战，并带来显著优势。首先，卫星网络拓扑时变、星间链路间歇连通，而SDN的集中控制能力使其能够基于全局视角进行路由计算与资源分配，从而提升网络管理效率。其次，卫星硬件部署后难以升级，SDN控制与转发分离的特性，以及软硬件解耦的设计，使得通过网络远程进行软件维护与功能更新成为可能，极大增强了网络的灵活性与可演进性。最后，面对由不同轨道、不同类型卫星及地面网络构成的异构融合网络，SDN提供了统一的编程接口与控制平面，能够实现跨域的协同资源调度与高效管理。目前，根据控制平面与数据平面中卫星类型的组合方式，可主要归纳为单层、双层与三层结构\citess{jiang2023software}。

 %    单层结构是最基础的形态，如图\ref{fig: SDN卫星单层}所示。通常仅由单一类型的卫星（以低轨LEO卫星最为常见）构成数据平面。其控制平面可部署于地面站，如文献\cite{wang2018design}通过地面数据中心管理LEO星座资源；亦可部署于LEO卫星本身，构成带内控制模式，如文献\cite{hu2018leo}及文献\cite{papa2020design}中考虑动态控制器迁移与重配置开销的LEO星上控制器方案，此种方式虽能降低传输时延，但需应对因卫星高速运动导致的控制器动态迁移与拓扑同步开销。

 %    双层结构则融合了两种不同类型的卫星，如图\ref{fig: SDN卫星双层}所示。最常见的组合是以覆盖广、位置相对静止的高轨GEO卫星作为控制平面，管理作为数据平面的LEO卫星\citess{xiaodong2020multi}。此种架构充分利用了GEO卫星的广域视野与链路稳定性，支持多QoS目标的自适应路由算法。此外，控制平面亦可设于地面站，同时管理GEO与LEO数据平面，如FRBSN\citess{sheng2017toward}；或采用地面站为超级控制器、GEO为主控制器、部分LEO为从控制器的混合分层模式\citess{xu2018software}，以增强网络控制的鲁棒性与覆盖完整性。
    
 %    三层结构最为复杂，如图\ref{fig: SDN卫星三层}所示。同时引入了GEO、MEO和LEO卫星，旨在构建一个功能分层、协同运作的立体网络。文献\cite{bao2014opensan}框架以GEO卫星群作为控制平面，统一管理包含GEO、MEO、LEO在内的整个数据平面，显著减少地面站数量。在星地融合场景下，文献\cite{li2017service}将控制功能分布于GEO卫星、地面数据中心与卫星网关，由卫星网络管理中心进行统一协调，数据平面则囊括了MEO/LEO卫星及地面路由器、交换机等设备，以实现端到端的精细化服务质量保障。文献\cite{li2019icn}将主控制器置于GEO，辅助控制器置于MEO以分担负载，LEO则专注于数据转发，并通过协同缓存等技术优化内容检索性能。此外，还有文献\cite{feng2014scheme}、文献\cite{feng2017hetnet}等架构将控制器部署于地面，统一调度多层卫星资源，支持异构网络融合与移动性管理。
    
 %    % 这些多样化的架构探索表明，软件定义卫星网络的设计需紧密结合具体应用场景与性能目标，在控制集中度、部署成本、传输时延与网络可靠性之间寻求最佳平衡。
 %    \begin{figure}[!htb]
 %        \centering
 %        \includegraphics[width=0.7\linewidth]{fig/SDN结构.pdf}
 %        \caption{软件定义网络三层架构}
 %        \label{fig: SDN结构}
 %    \end{figure}


 %    \section{卫星网络路由}
 %    % 参考 《卫星网络路由协议设计-訾元坤》 《基于QoS保障的卫星网络路由优化研究-陈晓燕》《卫星网络路由技术研究-吕若晨》  
 %    \subsection{低轨卫星网络路由}
 %    随着低轨卫星网络的快速发展，路由技术作为实现星间数据转发的核心机制，面临着卫星节点高动态性、网络拓扑频繁变化、星座规模不断扩大等挑战。传统的地面网络路由协议难以直接适用于卫星网络环境。目前，卫星网络路由技术主要可分为基于时间虚拟化和基于位置虚拟化两大类别\citess{吕若晨2022卫星网络路由技术研究}，以下对各类型的代表性方法进行分析。
    
 %    （1）基于时间虚拟化的路由
    
 %    基于时间虚拟化的路由技术是早期卫星网络路由研究的重要方向，其核心思想是利用卫星轨道运动的可预测性，将系统周期划分为离散的时间片，在每个时间片内将动态拓扑视为静态进行路由计算。这类技术主要解决了卫星高速运动带来的拓扑频繁变化问题。在该技术框架下，根据时间片划分依据的不同，又可进一步细分为两种主要方法。
    
 %    （a）等时间间隔划分。文献\cite{werner2002dynamic}首次提出将卫星网络运行周期平均划分为多个持续时间相等的时段，每个时段称为一个“拓扑快照”。在该时段内，网络拓扑被视为固定不变，并基于此静态拓扑预计算路由表。该方法实现简单、易于管理。文献\cite{王京林2009LEO}对此进行了优化，提出通过删除维持时间较短的星间链路来延长快照的持续时间，进一步提升了系统效率。然而，固定的时间划分可能导致快照切换时刻与实际拓扑变化时刻不同步，造成路由状态与网络实际状态的不匹配。

 %    （b）基于链路通断划分。为解决等时间间隔划分的不足，文献\cite{gounder1999routing}中提出了基于星间链路连通性变化的快照划分方法。该方法不再依赖于固定的时间间隔，而是以星间链路的建立或断开作为拓扑变化的触发条件，将系统周期划分为持续时间不等的拓扑快照。这种划分方式更贴近网络拓扑的实际变化，能够更精确地描述网络的动态特性。
 %    然而，该方法在提升精度的同时，也带来了新的挑战：需要实时监测链路状态，计算复杂度较高，且快照持续时间的不确定性增加了系统管理的难度。
    
 %    总体而言，基于时间虚拟化的路由技术通过将动态拓扑转化为系列静态拓扑，有效简化了路由计算过程，通常由地面控制中心进行路由预计算并上注至卫星网络。然而，随着星座规模的扩大，全网路由表的存储和更新将带来巨大的开销，且静态的路由计算难以适应网络的动态拥塞变化，这些局限性促使研究者转向其他技术路线。

 %    （2）基于位置虚拟化的路由
    
 %    为克服时间虚拟化技术的局限性，基于位置虚拟化的路由技术从空间维度提出了新的解决方案。该类技术的基本思路是通过将移动的卫星与固定的地理区域或虚拟节点建立映射关系，实现路由计算的简化，从而在空间上屏蔽卫星网络的动态性。根据映射对象的不同，该类技术可进一步分为两种典型方法：
    
 %    （a）基于覆盖区域虚拟化。文献\cite{ekici2000datagram}中提出了基于地面覆盖区域虚拟化的路由方案。该方法按照经纬度将地球表面划分为若干逻辑区域，每个区域被赋予唯一的逻辑地址。在卫星运行过程中，距离区域中心最近的卫星被动态指定为该区域的接入卫星，并继承相应的逻辑地址。这种方法的核心优势在于：路由计算仅需考虑源和目的端所在的固定逻辑区域，而无需关注具体服务卫星的变化，从而有效屏蔽了卫星的运动性。这些方法充分利用了星座拓扑的规则性和周期性，但当地面用户接入卫星切换时，需要进行路由信息的转移，可能引发额外的信令开销，在频繁切换场景下性能会受到一定影响。
    
 %    （b）基于星座网络区域虚拟化。为进一步优化系统性能，在后续研究\cite{ekici2002distributed}中提出了数据报路由算法(Datagram Routing Algorithm，DRA)。该方法不再将卫星映射到地面区域，而是将LEO卫星轨道球面划分为若干个虚拟区域，每个区域对应一个虚拟卫星节点，构成固定的虚拟拓扑结构。在该算法中，实际卫星根据其当前位置与虚拟节点建立动态映射关系，路由计算在虚拟拓扑上进行，计算结果再映射回实际卫星。DRA算法设计简单，仅需利用卫星位置信息即可完成路由决策，且LEO卫星间无需进行动态信息交互，大大降低了系统复杂度。然而，该方法也缺乏对网络实时状态的感知能力，难以应对网络拥塞问题，且其路由计算依赖于极轨星座的对称性，不适用于倾斜轨道星座，这限制了其在现代异构星座中的适用性。
    
 %    基于位置虚拟化的路由技术通过空间映射有效降低了路由计算的复杂度，但其性能在很大程度上依赖于星座结构的规则性，且通常需要预存卫星位置信息。随着星座规模的扩大和轨道类型的多样化，这些方法的适应性和扩展性面临严峻挑战。

 %    \subsection{多层卫星网络路由}
 %    多层卫星网络通过整合不同轨道高度（如GEO、MEO、LEO）的卫星，实现了全球覆盖、负载均衡与容错能力的提升，已成为未来空间信息网络的核心架构。然而，其网络拓扑的高动态性和异构性对路由算法设计提出了严峻挑战。为应对这一挑战，研究者提出了多种路由协议，其中，基于分层管理思想的协议尤为关键，它们通过将网络控制功能上移至高层卫星，以简化路由管理并提升网络可扩展性。

 %    文献\cite{lee2000satellite}提出的分层卫星路由协议（Hierarchical Satellite Routing Protocol，HSRP），是面向LEO/MEO双层网络的典型代表。其工作流程为：各LEO卫星周期性地与邻居交换HELLO数据包，以获取局部拓扑状态信息；随后，LEO卫星通过层间链路上报这些信息至MEO卫星；MEO卫星之间通过星间链路交互信息，从而获取LEO层的全网拓扑结构；最终，MEO卫星将整合后的全网拓扑信息分发至各LEO卫星，由LEO卫星自行计算其路由表。HSRP协议有效降低了LEO卫星的计算负担，但MEO层卫星为收集全网信息而进行的频繁交互，会引入一定的控制开销与传播时延。

 %    文献\cite{周云晖2008卫星网络}提出时分路由协议（Time Division Routing Protocol，TDRP），在分组管理思想上融合了时间分片策略。其工作流程是一个自下而上信息收集、自上而下路由分发的分层协同过程。在时隙周期开始时，底层LEO卫星负责收集包括链路时延和有效带宽在内的动态链路状态信息，并根据服务时间最长原则选择MEO卫星作为其域管理者；MEO管理者随后对所属LEO卫星进行逻辑域划分，并计算域内路由，同时通过MEO层间交互完成全局LEO域拓扑信息的同步；MEO层将链路状态和域结构信息继续向GEO层汇聚，由GEO管理者完成MEO域的划分与域内路由计算。在信息收集阶段完成后，流程进入路由计算与分发阶段：GEO层首先基于全局拓扑独立计算层内路由，再结合下层域结构信息逐层推导出MEO域间路由表和LEO域间路由表，最终通过分层下发机制将完整的路由表分发至各MEO和LEO域成员卫星，实现全网路由状态的协同更新。
    
    
    \section{本章小结}
    本章系统阐述了本文研究所涉及的相关技术背景。首先，梳理了卫星网络从单层 LEO 向多层异构架构演进的必然性，以及管控模式从传统分布式向 SDN 分层协同演进的技术脉络；其次，分析了卫星路由技术的发展历程，重点指出了现有单层路由策略在应对异构网络多维 QoS 需求时的局限性；最后，对比了网络数值仿真与半实物模拟技术，明确了基于容器虚拟化的仿真方法在验证真实协议栈性能方面的独特优势。上述研究现状与技术分析，为本文提出的分级多域 SDN 架构、改进遗传算法路由机制以及基于 STK-Docker 的仿真平台设计提供了坚实的理论依据与技术支撑。

    % 参考《基于QoS保障的卫星网络路由优化研究-陈晓燕》
    % 《软件定义卫星网络多控制器部署策略_陈金涛》
    \chapter{异构星座组网与管控架构研究}
    \label{chap: 异构星座组网与管控架构研究}

    \section{异构星座场景概述}
    面向未来通导遥一体化空间信息网络建设的国家重大需求，单一功能的卫星星座已难以满足多样化的业务保障要求。本项目依托某大型异构星座系统背景，构建了一个包含低轨感知层、低轨通信层、中轨中继层及高轨管控层的四层深度融合网络架构。该场景旨在解决海量遥感数据的实时回传与全球用户的宽带接入问题，同时利用多层轨道优势实现网络的高效管控。针对异构星座中不同节点的任务属性，本研究将空间段划分为以下四个逻辑与物理层级：
    
    LEO 感知层，由运行在约 700 km 高度、45° 倾角的 100 颗低轨卫星组成。该层卫星主要搭载高分辨率光学或雷达载荷，负责对地观测与数据采集。作为网络中的主要数据源节点，其具备海量数据的生成能力，但星上存储与传输资源受限，需依托上层通信网络实现数据的准实时回传。
    
    LEO 通信层，由运行在约 1000 km 高度、70° 倾角的 225 颗宽带通信卫星组成。该层采用 Walker 星座构型，通过星间链路（ISL）构建高速骨干传输网络。其核心功能是接入与转发，既负责地面用户终端的宽带接入，也承担底层感知数据的汇聚与接力传输任务。
    \begin{table}[htbp]
      \centering
      \caption{特定异构星座仿真参数配置}
      \label{tab: 特定异构星座仿真参数配置}
      \resizebox{\textwidth}{!}{% 自动缩放表格以适应页面宽度
      \begin{tabular}{l c c c c}
        \toprule
        \textbf{参数} & \textbf{LEO 感知层 (Detect)} & \textbf{LEO 通信层 (LEO)} & \textbf{MEO 中继层 (MEO)} & \textbf{GEO 管控层 (GEO)} \\
        \midrule
        功能定位 & 数据源/感知 & 传输/接入 & 跨域中继 & 全局管控 \\
        轨道类型 & 倾斜圆轨道 & 倾斜圆轨道 & 倾斜圆轨道 & 地球静止轨道 \\
        轨道高度 (km) & $\approx 700$ & $\approx 1,000$ & $\approx 7,000$ & $\approx 35,786$ \\
        轨道倾角 & $45^\circ$ & $70^\circ$ & $45^\circ$ & $0^\circ$ \\
        卫星总数 ($N$) & 100 & 225 & 8 & 3 \\
        轨道面数 ($P$) & 10 & 15 & 4 & 1 \\
        面内星数 ($S$) & 10 & 15 & 2 & 3 \\
        相位因子 ($F$) & 1 & 1 & 1 & - \\
        RAAN 间隔 & $36^\circ$ & $24^\circ$ & $90^\circ$ & - \\
        \bottomrule
      \end{tabular}%
      }
    \end{table}
    
    MEO 中继层，由运行在约 7000 km 高度的 8 颗中轨卫星组成。相较于 LEO 卫星，MEO 卫星单星覆盖范围更大，可见时间窗口更长。在网络架构中，该层作为跨域协作的稳定中继节点，用于承载跨洋或跨洲的长距离跳数传输，有效降低纯 LEO 多跳传输带来的累积时延与丢包风险。
    
    GEO 管控层，由 3 颗地球静止轨道卫星组成，实现对地球表面的全时全域覆盖。在本架构中，GEO 层不仅可作为应急通信的兜底传输手段，更核心的作用是作为全局主控制器的部署载体。利用其广域视野与相对静止特性，负责全网拓扑维护、跨域路由规划及控制信令的统一下发。
   
    基于上述异构星座设计，本研究利用 STK (Systems Tool Kit)\citess{wall2024systems} 软件构建了高保真的轨道动力学仿真场景。场景中各层星座的具体轨道参数配置如表\ref{tab: 特定异构星座仿真参数配置}所示。该异构星座模型的3D和2D视图，分别如图\ref{fig: STK_3D}和图\ref{fig: STK_2D}所示。
    \begin{figure}[!htb]
		\centering
		\subfloat[]{
			\includegraphics[width=0.8\linewidth]{fig/STK_3D.png}
			\label{fig: STK_3D}
		}
		\hfill
		\subfloat[]{
			\includegraphics[width=0.8\linewidth]{fig/STK_2D.png}
			\label{fig: STK_2D}
		}
		\caption[仿真星座模型的3D和2D视图]{异构星座模型的3D和2D视图。\subref{fig: STK_3D}3D视图；\subref{fig: STK_2D}2D视图} \label{fig: 星座模型的3D和2D视图}
	\end{figure}

    \section{异构星座网络分级多域协同管控架构}
    针对异构星座网络节点规模大、拓扑动态变化快以及业务需求多样化的特点，传统的集中式或完全分布式控制架构难以同时满足全局优化与实时响应的需求。为此，本节提出一种分级多域协同管控架构。该架构在垂直维度上采用分级控制，以分离全局策略制定与局部网络维护；在水平维度上采用多域划分，将大规模网络划分为若干自治区域进行管理。 首先本节先分析异构星座网络控制器部署方案。

    1、控制器部署方案分析
    
    在软件定义卫星网络架构中，控制器部署策略主要分为星上单层控制、地面控制、星地联合控制以及星上多层控制四种方案。如表\ref{tab:controller-deployment}所示，各类方案在覆盖范围、计算能力、网络时延和适应动态性方面各有优势\citess{万颖2022基于}。
    \begin{table}[htbp] % 根据排版需求选择合适的选项，如 [h]、[t]、[b] 等
        \centering
        \caption{控制器部署方案对比}
        \label{tab:controller-deployment}
        \begin{tabular}{llcccc}
            \toprule
            \textbf{部署位置} & \textbf{部署方案} & \textbf{覆盖范围} & \textbf{计算能力} & \textbf{网络时延} & \textbf{网络动态性} \\
            \midrule
            地面 & 单层 & 部分 & 较高 & 较低 & 低 \\
            LEO & 单层 & 全球 & 较低 & 一般 & 高 \\
            GEO & 单层 & 全球 & 较低 & 较高 & 中 \\
            地面+GEO & 多层 & 全球 & 较高 & 较高 & 中 \\
            GEO+LEO & 多层 & 全球 & 适中 & 适中 & 高 \\
            \bottomrule
            \end{tabular}
    \end{table}

    通过对比分析发现，对于本项目所涉及的大规模异构星座场景，单一层级或依赖地面的控制方案均存在显著瓶颈。 地面集中控制受限于星地链路的长传播时延与窄带宽，难以应对 LEO 网络的分钟级拓扑重构，易导致控制指令失效及链路拥塞；LEO 分布式控制虽感知敏捷，但受限于单星算力与频繁的控制器切换，难以支撑全网全局优化；而星地联合控制本质上仍依赖地面决策，无法满足天基网络独立自治的需求。

    相比之下，GEO 与 LEO 结合的星上多层部署方案利用轨道层级差异实现了最佳的功能互补。该方案利用 GEO 卫星的广域覆盖与相对静止特性弥补了 LEO 层的全局视野缺失，同时利用 LEO 层的贴近用户特性弥补了 GEO 层的时延劣势。MEO 虽然覆盖大，但其对于 LEO 的相对运动仍然存在，不如 GEO 静止，因此不承担控制职能，适合作为独立的中继节点受控于顶层。综上所述，采用 GEO 作为全局控制节点、LEO 作为区域控制节点的星上分级部署策略，在摆脱地面依赖的同时有效兼顾了全局优化能力与局部响应速度，是适合本研究场景的合理选择。
    
    2、分级多域协同管控机制
    
    基于上述部署方案，本架构建立的分级多域管控架构如图 \ref{fig: 分层管控架构} 所示。该机制通过垂直维度的分级控制与水平维度的多域划分，实现了全网资源的协同调度。
    \begin{figure}[!htb]
        \centering
        \includegraphics[width=0.8\linewidth]{fig/分层管控架构}
        \caption{异构星座分级多域管控架构物理示意图}
        \label{fig: 分层管控架构}
    \end{figure}

    在垂直维度上，架构依据时效性与覆盖范围将控制平面划分为两级。一级管控由部署在 GEO 轨道的全局控制器构成，利用其相对静止且覆盖广的特性，负责维护全网的宏观抽象拓扑。对于 MEO 卫星，由于其作为关键的中继资源且数量较少，由一级管控层直接进行状态监控与调度，不参与下层域控制器的划分。二级管控则由部署在 LEO 层的域控制器构成，侧重于处理局部性、高实时性的网络控制任务，负责维护本域内的微观物理拓扑。

    在水平维度上，架构主要针对大规模 LEO 星座带来的管理复杂度问题，将海量低轨节点在空间上划分为若干个自治管理域。依据地面地理网格，将物理位置投影在同一地理区域内的 LEO 卫星群划分为一个逻辑域。无论卫星如何运动，只要位于该地理网格覆盖范围内，即归属该域管理。域内的拓扑变化仅在域内进行广播和处理，无需频繁上报至全局控制器，从而显著降低了全网控制信令开销。

    为了明确各层级的交互逻辑与功能定义，管控架构的逻辑视图如图 \ref{fig: 异构星座协同管控架构逻辑示意图} 所示。系统在逻辑上划分为一级全局管控平面、二级域管控平面以及数据转发平面，各平面通过标准接口形成感知与执行的闭环。
    \begin{figure}[!htb]
        \centering
        \includegraphics[width=0.7\linewidth]{fig/异构星座协同管控架构逻辑示意图.pdf}
        \caption{异构星座分级多域管控架构逻辑示意图}
        \label{fig: 异构星座协同管控架构逻辑示意图}
    \end{figure}
    
    一级管控即全局管控平面，核心职能在于全网资源统筹。该平面维护全网虚拟拓扑视图，根据业务需求计算跨域长距离传输的宏观路径，并统筹分配全网资源，将端到端 QoS 指标分解为分段预算。特别地，为了确保关键中继资源的调度时效性，该平面通过直控链路直接负责 MEO 中继节点的调度与状态监控，不经由下层域控制器转发。全局控制器通过层间接口向二级域控制器下发宏观路由策略与资源配置参数，并接收各域上报的聚合状态快照。

    二级管控即域管控平面，负责多域自治管理。域控制器周期性感知域内物理链路的通断与拥塞状态，依据全局下发的宏观策略，利用 IGA 算法计算域内微观路径，并生成 OpenFlow 流表下发至底层转发节点，同时管理域内虚拟节点的映射与迁移。在多域协同方面，当业务涉及跨域传输时，相邻的域控制器通过东西向接口交互边界链路状态，或请求全局层协调 MEO 中继节点进行跨域接力，从而保障业务传输的平滑切换。
    
    最下层为数据转发平面，由大规模 LEO 卫星节点群与 MEO 中继节点共同构成。LEO 节点群依据接收到的流表执行数据包的高速转发，并实时监测底层链路状态，将故障或拥塞告警上报至所属的域控制器。MEO 中继节点则作为跨层协作枢纽，主要承担跨域长距离传输任务，其状态直接上报至全局管控平面，从而形成完整的感知与执行闭环。
    

    \section{基于虚拟节点的异构星座网络管理域构建}
    \label{sec: 区域构建}
    针对 LEO 星座高速运动导致的网络拓扑高频变化问题，本节提出一种基于地面地理网格的虚拟节点构建策略。该策略的核心思想是将地球表面划分为若干固定的地理区域，每个区域抽象为一个逻辑静止的地面管理域，当卫星移动出某区域、下一颗卫星覆盖该区域时，虚拟节点映射关系自动切换，由新卫星继承原卫星的逻辑地址和网络状态。无论物理卫星如何运动，其逻辑身份由其当前覆盖的地面区域决定，从而将高动态的物理网络映射为逻辑静止的虚拟网络，有效屏蔽了底层的物理动态性。考虑到LEO通信层与LEO感知均面临卫星高速运动带来的拓扑频繁重构挑战，本节提出的虚拟节点构建策略将统一应用于这两层网络。
    
    \subsection{地面虚拟管理域的划分} %改为 地面虚拟管理域的划分
    \label{sec: 地面虚拟管理域的划分}
    区域划分不仅是控制域的物理边界定义，更是全网负载均衡与控制稳定性的基础。传统的均匀经纬度网格在低轨卫星场景下存在显著缺陷：一方面，受地球几何形状影响，经线在两极逐渐收敛，导致高纬度区域网格的物理面积急剧缩小；另一方面，网格面积的缩小会导致卫星穿越该区域的频率过高，进而引发频繁的控制器切换与巨大的信令交互开销，降低了网络管控的效率。为此，本文建立一种适应轨道特性的变步长网格划分模型。该模型旨在确保单一网格内拥有足够的卫星节点冗余度，并通过纬度自适应调整网格跨度以实现物理面积的相对均衡，及逻辑拓扑的规整性与边界对齐。具体划分策略如下：

    （1）经度方向：统一划分，为避免高低纬度交界处出现边界错位，全网采用统一的经度步长 $\Delta \lambda$。这确保了所有逻辑域在纵向上具有连续贯通的边界，简化了跨域路由的邻居维护逻辑。
  
    （2）纬度方向：变步长划分，依据 LEO 星座的轨道倾角特性与覆盖密度分布，采用非均匀的纬度跨度。在卫星密集的低纬度核心区采用较小的纬度跨度，以细化管控粒度。在中高纬度及极地区域，适当扩大纬度跨度，以适应卫星相对密度的变化。
    
    通过这种划分，将地球表面离散化为集合 $\mathcal{R} = \{R_1, R_2, ..., R_N\}$，其中每个 $R_i$ 代表一个具有近似相等物理面积和卫星覆盖密度的逻辑控制域。在确立地理网格后，核心挑战在于如何实现物理节点与逻辑区域的解耦。本架构基于虚拟机点机制，将网络标识与地理位置绑定，而非与特定卫星绑定。每个划分好的区域$R_i$被赋予一个永久的逻辑标识ID（如IP虚拟地址段）。该标识在SDN控制平面中是静止不变的，作为路由计算的锚点。物理卫星$S_k$并不永久拥有某个IP，而是作为资源载体。当$S_k$运行进入区域$R_i$时，它通过SDN南向接口向控制器注册，并动态加载$R_i$对应的虚拟身份配置，成为该区域的执行节点。对于状态继承，当$S_k$飞离$R_i$时，它卸载该身份；与此同时，新进入的卫星$S_{k+1}$立即继承该虚拟身份。这种机制使得GEO主控制器仅需维护$N$个静态虚拟节点的拓扑关系，而无需感知底层成千上万颗LEO卫星的具体编号变化，从而极大地降低了层间控制链路的信令开销。对于LEO通信层和LEO感知层，系统分别维护两套独立的虚拟IP地址空间，互不干扰。

    在每个虚拟区域$R_i$内，可能同时覆盖多颗LEO卫星。为了执行域内的流表管理和局部路由计算，需从当前覆盖集合$S_{cov}(t)$中选举出一颗卫星作为域控制器（DomainController，DC）。不同于地面网络侧重于带宽资源的选举，本方案基于剩余服务时间最大化原则进行DC角色绑定：
    \begin{equation}
        DC_i(t)=\arg\max_{s\in S_{cov}(t)}\left(\alpha \cdot T_{dwell}(s,R_i)+\beta \cdot \frac{1}{L_{geo}(s)}\right)
    \end{equation}
    
    其中，$T_{dwell}$表示卫星$s$在区域$R_i$内剩余的驻留时间，该值越大，控制器角色越稳定，避免了频繁迁移带来的流表震荡。$L_{geo}$表示该LEO卫星与上方GEO主控制器或MEO中继星的链路时延。由于本架构采用分层管控，域控制器必须保持与上层的顺畅连接。其中，$\alpha$ 和 $\beta$ 为非负权重系数，且 $\alpha + \beta = 1$。这两个参数可根据实际业务场景对网络性能的偏好进行灵活配置：若业务对控制器的稳定性要求较高，可增大 $\alpha$ 以优先选择驻留时间长的节点；若业务对时延敏感，则可增大 $\beta$。在本文后续的仿真分析中，综合考虑稳定性与时延的平衡，取 $\alpha=0.6, \beta=0.4$。

    通过上述构建方案，将高动态的 LEO 星座转化为一个逻辑静止、物理轮替的虚拟网络。针对本项目中包含 225 颗卫星的 LEO 通信星座，举例说明具体的网格划分策略设计如下。 纬度带划分：将南北半球各划分为两个非均匀区间。低纬度带定义为 $0^{\circ} \sim 40^{\circ}$，跨度 $40^{\circ}$，中高纬度带定义为 $40^{\circ} \sim 90^{\circ}$，跨度 $50^{\circ}$，全网共形成 4 个纬度带。经度区划分：设定全球统一的经度步长 $\Delta \lambda = 60^{\circ}$，将全球沿经度方向均分为 6 个区间。
        综合上述划分，地球表面被离散化为集合 $\mathcal{R} = \{R_1, R_2, ..., R_{24}\}$，全网由 24 个 边界完全对齐的逻辑控制域组成。地面虚拟管理域的划分如示意图\ref{fig: LEO通信星座分域3D和2D示意图}所示。根据仿真测算，在该划分方案下，平均每个网格内的卫星覆盖数约为 9.4 颗，既通过变步长策略适应了轨道分布，又保证了逻辑拓扑的规则性。
     
    \begin{figure}[!htb]
		\centering
		\subfloat[]{
			\includegraphics[width=0.8\linewidth]{fig/LEO分域2D示意图-24个域.png}
			\label{fig: LEO分域2D示意图}
		}
		\hfill
		\subfloat[]{
			\includegraphics[width=0.8\linewidth]{fig/LEO分域3D示意图-24个域.png}
			\label{fig: LEO分域3D示意图}
		}
		\caption[LEO通信星座分域3D和2D示意图]{LEO通信星座分域3D和2D示意图。\subref{fig: LEO分域2D示意图}3D示意图；\subref{fig: LEO分域3D示意图}2D示意图} \label{fig: LEO通信星座分域3D和2D示意图}
	\end{figure}

    \subsection{地面虚拟管理域的维护} % 地面虚拟管理域的维护
    在基于虚拟节点的异构星座架构中，区域维护的核心挑战在于物理节点的高速运动导致其与逻辑区域的映射关系频繁更替。为了保障控制平面的连续性，本节设计了一种基于轨道预判的主动式区域维护机制，重点解决域控制器（DC）的角色迁移与跨域流表的一致性更新问题。

    不同于移动自组网（MANET）中节点运动的随机性，LEO卫星严格遵循轨道动力学规律。本方案利用高精度星历数据，在时间维度上引入预切换时间窗($T_{win}$)。对于任意区域$R_m$内的卫星$S_k$，域控制器周期性计算其剩余驻留时间$T_{rem}$。当$T_{rem}<T_{win}$时，触发区域维护流程。时间窗的大小需满足状态同步的时延约束：
    \begin{equation}
        T_{win}\ge\frac{D_{state}}{B_{ISL}}+T_{prop}+T_{proc}
    \end{equation}
    
    其中$D_{state}$为控制器需同步的状态数据量（包括全域流表、拓扑视图、QoS预算等），$B_{ISL}$为星间链路带宽。该预判机制避免了硬切换带来的服务中断风险。当现任控制器 $DC_{old}$ 即将飞离区域边界时，必须将控制权移交给候任控制器 $DC_{new}$。为实现零宕机迁移，本文提出一种迁移策略如下。具体流程如图\ref{fig: 域控制器平滑迁移时序图}所示。
    \begin{figure}[!htb]
        \centering
        \includegraphics[width=1.0\linewidth]{fig/域控制器平滑迁移时序图.pdf}
        \caption{域控制器平滑迁移时序图}
        \label{fig: 域控制器平滑迁移时序图}
    \end{figure}
    
    （1）基于\ref{sec: 地面虚拟管理域的划分}节的选举算法确定$DC_{new}$后，$DC_{old}$通过轨道内链路向$DC_{new}$发起状态同步请求。$DC_{new}$启动SDN控制器进程，建立与$DC_{old}$的数据同步通道，实时拉取当前的全局拓扑快照和活动流表项。
    
    （2）状态同步完成后，进入双活阶段。此阶段原控制器（$DC_{old}$）尚未下线，仍保持与上层GEO及下层成员卫星的活跃连接，继续履行控制职能，保障存量业务不中断；候任控制器（$DC_{new}$）完成流表与拓扑数据的全量同步，随时准备接管流量；
    $DC_{old}$向GEO全局控制器发送位置更新信令，声明虚拟节点$VN_m$的物理映射即将变更。此时，GEO更新其南向接口映射表，将发往$VN_m$的控制指令同时下发给$DC_{old}$和$DC_{new}$。

    （3）当卫星实际越过逻辑边界时刻，$DC_{new}$正式接管该区域的虚拟IP地址，向域内所有成员卫星广播成为域控制器的通告。$DC_{old}$卸载控制逻辑，降级为普通转发节点或进入下一个区域成为新成员。
    
    （4）$DC_{new}$向GEO反馈接管完成信号，GEO停止向$DC_{old}$发送指令，完成迁移闭环。

    \section{基于虚拟地址的异构星座互联机制}
    异构星座的互联面临着动态拓扑高频变化、多轨道协议差异、跨域寻址复杂等核心挑战。LEO 卫星的高速运动导致网络拓扑呈分钟级重构，MEO/GEO 卫星的异构协议栈易造成路由策略冲突，而终端节点跨域通信时物理地址与逻辑位置的强耦合性更会导致路由效率急剧下降。针对上述问题，本文提出基于虚拟地址的异构星座互联机制，通过构建逻辑地址与物理拓扑的解耦架构，实现动态环境下的稳定互联。
    
    \subsection{地址标识体系与报文封装}
    为有效屏蔽低轨卫星高速运动带来的网络拓扑动态性，实现异构网络间的稳定互联，本文提出一种物理与逻辑分离的地址标识体系。该策略的核心在于解耦节点的物理实体与逻辑位置，通过在 LEO 卫星节点引入两套并行的地址标识，并配合隧道技术重构数据包格式，从而构建稳定的骨干路由架构。
    
    在本架构中，LEO 卫星节点需同时维护物理与虚拟两类地址标识，两者分别承载不同的网络职能。物理 IP 地址作为卫星节点的固有属性标识，与卫星硬件实体深度绑定。该地址具有永久唯一性，且不随卫星空间位置的变化而改变。在功能上，物理 IP 主要作为特定卫星的寻址锚点：一方面支持地面运控中心进行状态监控、指令上注及故障诊断等管理操作；另一方面，在拓扑结构相对稳定的同层协作或邻域通信场景中，作为底层逐跳路由的寻址依据。与物理标识相对应，虚拟 IP 地址则是地面管理域的逻辑标识，与前文节所述的地面地理网格建立一一映射关系。该地址在逻辑拓扑中保持绝对静止，作为跨层骨干路由的稳定锚点。物理卫星本身并不永久持有虚拟 IP，而是采用动态绑定机制实现地址复用：当 LEO 卫星运行至某地面网格覆盖范围内时，自动获取并加载该网格对应的虚拟 IP；当卫星飞离该区域时，立即解除绑定并卸载地址。这种机制确保了无论底层物理节点如何更替，上层网络感知的逻辑拓扑始终保持稳定。

    基于上述双重地址体系，本文设计了一种双层报文结构以承载跨域通信业务。该报文由内层物理报头和外层虚拟报头嵌套组成。其中，内层物理报头保留了原始 IP 报头的完整信息，包含源节点和目的节点的本地物理 IP 地址，这保证了数据包在源端生成和末端交付时，对现有终端设备和应用层协议保持完全透明。外层虚拟报头则在跨域传输过程中由边缘节点动态添加，包含源节点和目的节点所属地面管理域的虚拟 IP 地址。这一设计使得 MEO/GEO 等高层节点无需解析内层信息，仅依据外层虚拟地址即可完成路由转发。具体的数据包字段定义如表\ref{tab:packet_header_structure}所示。
    \begin{table}[htbp]
        \centering
        % 表格标题
        \caption{数据包结构定义}
        \label{tab:packet_header_structure}
        \begin{threeparttable}
            % 定义列格式：l代表左对齐，c代表居中。
            % 第三列使用 p{<宽度>} 可实现自动换行，防止描述过长撑爆版面，此处设为 8cm
            \begin{tabular}{clp{8cm}}
                \toprule
                \textbf{报头层级} & \textbf{字段名称} & \textbf{字段描述} \\
                \midrule
                % 外层虚拟报头部分
                \multirow{2}{*}{外层虚拟报头} & \texttt{src\_vir} & 源节点所处 LEO 域的逻辑虚拟地址 \\
                 & \texttt{dest\_vir} & 目的节点所处 LEO 域的逻辑虚拟地址 \\
                \midrule
                % 内层物理报头部分
                \multirow{2}{*}{内层物理报头} & \texttt{src\_phy} & 源节点的本地物理地址 \\
                 & \texttt{dest\_addr} & 目的节点的本地物理地址 \\
                \midrule
                % 数据载荷部分
                数据载荷 & \texttt{payload} & 上层应用数据载荷 \\
                \bottomrule
            \end{tabular}
        \end{threeparttable}
    \end{table}

    数据包在全网的传输过程，实质上是双重地址在不同阶段交替发挥作用的过程。为清晰阐述这一机制，本文将数据包的生命周期划分为四个关键阶段，如图\ref{fig: 数据包封装解封过程}所示。
    \begin{figure}[!htb]
        \centering
        \includegraphics[width=0.9\linewidth]{fig/数据包封装解封过程.pdf}
        \caption{数据包生命周期示意图}
        \label{fig: 数据包封装解封过程}
    \end{figure}
    
    （1）源端生成阶段，当源节点（如地面终端或执行观测任务的卫星）产生数据时，仅依据本地物理网络视图进行处理。源节点将数据载荷封装在标准的 IP 报文中，头部仅包含源物理地址与目的物理地址。此时数据包处于原始物理包状态，尚未携带任何虚拟位置信息，确保了源端设备的协议兼容性。
    
    （2）入口封装阶段，当数据包到达源节点所在域的 LEO 接入卫星时，系统执行路由策略检测。若判定业务流需要跨层传输，接入卫星随即依据源节点和目的节点的地理位置解析出对应的虚拟 IP 地址，并将其作为外层报头附加在原始数据包之前，完成隧道封装。
    
    （3）跨域传输阶段，封装后的数据包进入异构网络层，中间的转发节点仅依据外层虚拟目的地址进行路由决策。由于虚拟地址锚定于静止的地面网格，中间节点无需感知底层 LEO 卫星的快速移动与切换，从而实现了数据包在动态拓扑下的稳定透传。
    
    （4）出口解封与交付阶段，当数据包到达目的地面域的任意一颗卫星（即目的域入口节点）时，该节点识别出虚拟目的地址属于本地域，虚拟隧道的传输任务即告终止。该节点随即剥离外层虚拟报头，将数据包还原为原始物理状态。 此时，系统转入域内物理路由流程：若目的终端直接连接于当前节点，则直接下发；若目的终端连接于同域内的其他卫星，则依据内层物理目的地址，通过星间链路将数据包转发至末端接入卫星，最终通过星地链路精准交付给目的节点。
    

    \subsection{典型业务传输场景}
    基于上述封装协议，系统根据业务流向的不同，自适应地切换单层传输与跨层传输两种模式，以平衡传输效率与协议兼容性。
    
    （1）单层传输模式
    
    当业务流局限于单一轨道域，例如低轨观测星产生数据，经同轨道的其他观测星跳传，回传至同一覆盖区内的地面站时，数据流无需经过异构网络。如图\ref{fig: 单层传输场景}所示，数据沿 D1 $\rightarrow$ D3 $\rightarrow$ D4 $\rightarrow$ D6 $\rightarrow$ D5 $\rightarrow$ 地面站F1的路径传输。
    在此模式下，系统沿用传统的物理编址路由，各卫星节点配置固定的物理地址（例如 34.172.x.x 网段）。由于同层 LEO 卫星间的相对运动较小，拓扑结构虽有变化但邻居关系相对可预测，直接使用物理 IP 进行逐跳转发足以保障连通性。此过程不涉及虚拟地址的封装与解封装操作，避免了额外的报头开销，实现了轻量级的高速转发。
    \begin{figure}[!htb]
        \centering
        \includegraphics[width=0.7\linewidth]{fig/单层传输场景_01.pdf}
        \caption{单层传输场景}
        \label{fig: 单层传输场景}
    \end{figure}

    （2）跨层传输模式
    
    当业务流需跨越不同轨道层或协议域，例如LEO观测星执行对地观测任务，将采集数据经由LEO通信星、MEO通信星进行转发，长距离回传至异地地面站时，需要进行跨层传输。如图\ref{fig: 跨层传输场景}所示，低轨观测星D1将数据发送至LEO通信星，再经由MEO中中继转发，最后到达地面站。在此场景中，物理拓扑是动态的，但逻辑区域是静态的。例如对于物理编址，LEO节点可使用35.172.x.x网段，MEO可使用36.172.x.x网段，GEO可使用37.172.x.x网段，这些地址用于节点间的链路层通信；对于虚拟层编址，系统在10.0.x.x网段构建虚拟叠加网络，每个虚拟IP唯一对应地面的一个地理栅格（即\ref{sec: 区域构建}节划分的区域）。
    \begin{figure}[!htb]
        \centering
        \includegraphics[width=0.9\linewidth]{fig/跨层传输场景_01.pdf}
        \caption{跨层传输场景}
        \label{fig: 跨层传输场景}
    \end{figure}
    
    该场景中，观测星D2生成封装有物理报头（源D2\_PhyIP、宿F2\_PhyIP）的原始数据包，并发往接入卫星L1。随后L1识别跨域流量，依据网格映射关系为数据包封装外层虚拟报头（源L1\_VirIP/10.0.10.0，目F2\_VirIP/10.0.20.0）。封装后的数据包沿 L1$\rightarrow$L2$\rightarrow$L3$\rightarrow$M1 上传至MEO层中继。中间节点仅依据虚拟报头路由，构建起屏蔽底层拓扑动态变化的逻辑隧道。最后数据包经 M1$\rightarrow$L4$\rightarrow$L5 抵达目的域入口节点L4。L4剥离外层虚拟报头以还原物理包，并依据内层F2\_PhyIP将其通过星地链路精准交付至地面站F2，实现了全网路由的逻辑统一与物理透明。通过这种本地地址到虚拟地址的映射机制，本架构屏蔽了异构星座中物理节点的移动性，实现了路由策略在逻辑层面的统一与稳定。

    \section{仿真实验与结果分析}
    % 参考《基于QoS保障的卫星网络路由优化研究-陈晓燕》
    % 《软件定义卫星网络关键技术研究-徐双》

    % \textbf{GEO总控和所看到域控制器、MEO的覆盖时长、}
    % TODO

    \subsection{仿真环境与参数配置}
    \label{sec: STK参数配置}
    为了验证本文提出的GEO 全局主控 、LEO 域内自控、 MEO 中继转发的组网架构有效性，本节采用 STK (Systems Tool Kit) 构建高保真轨道动力学仿真场景，并结合Matlab进行网络性能数据的统计与分析。

    仿真场景构建了一个包含低轨观测、低轨通信、中轨和高轨的三层异构卫星通信网络。其中，Detect为观测类卫星，分布在10个轨道面，每个轨道面10颗，卫星命名规则为Detectxxyy，xx表示轨道面编号，yy为当前轨道面内卫星编号，探测类卫星也可以作为信息传输的中间节点，也可作为信息发送方和接收方；LEO为低轨通信卫星，作为用户接入与数据转发层，分布在15个轨道面，每个轨道面15颗，命名规则为LEOxxyy，xx表示轨道面编号，yy为当前轨道面内卫星编号；MEO 卫星作为跨层协作的稳定中继层，分布在4个轨道面，每个轨道面2颗，命名规则为MEOxy，x表示轨道面编号，y为当前轨道面内卫星编号；GEO 卫星作为全局策略管控层，共3颗。场景中各类星座节点的轨道参数及分布同表\ref{tab: 特定异构星座仿真参数配置}所示。此外，依据地面式的控制架构中，地面站做为从控制器采集卫星状态并汇聚至地面的网络操作控制中心，仿真场景在全球主要区域及岛屿部署部署48个地面站\citess{zhu2017software}，作为本文的对比基准。
    
    % \begin{table}[htbp]
    %   \centering
    %   \caption{星座主要仿真参数配置}
    %   \label{tab: stk_config}
    %   \begin{tabular}{c c c c c}
    %     \toprule
    %     \textbf{参数} & \textbf{Detect} & \textbf{LEO} & \textbf{MEO} & \textbf{GEO} \\
    %     \midrule
    %     轨道类型          & 倾斜圆轨道                     & 倾斜圆轨道                     & 倾斜圆轨道                      & 地球静止轨道 \\
    %     轨道高度          & $\approx 700\,\text{km}$       & $\approx 1000\,\text{km}$      & $\approx 7000\,\text{km}$       & $\approx 35786\,\text{km}$ \\
    %     倾角              & $45^\circ$                     & $70^\circ$                     & $45^\circ$                      & $0^\circ$ \\
    %     偏心率            & $\approx 0$                    & $\approx 0$                    & $\approx 0$                     & $0$ \\
    %     总卫星数 ($N$)    & 100                            & 225                            & 8                               & 3 \\
    %     轨道面数 ($P$)    & 10                             & 15                             & 4                               & 1 \\
    %     相位因子 ($F$)    & 1                              & 1                              & 1                               & N/A \\
    %     相邻轨道面 RAAN 间隔 & $36^\circ$                  & $24^\circ$                     & $90^\circ$                      & N/A \\
    %     \bottomrule
    %   \end{tabular}
    % \end{table}
        
    \subsection{仿真结果与分析}
    
    1、链路特性分析 % 链路之间的可见性、距离变化
    
    为了深入揭示异构星座网络拓扑的时空演进规律，本文选取了网络中具有代表性的 Detect、LEO、MEO 卫星及地面信关站作为观测节点，对其在 24 小时仿真周期内的链路连通状态进行了可视化分析。图\ref{fig: 典型链路可见性}展示了不同类型链路的可见性甘特图。其中，LEO层的卫星LEO0101维持着前后左右四个方向的稳定邻居连接（LEO0102/LEO0115/LEO0201/LEO1515），逻辑上形成了近似规则的网格状结构；与其他轨道间邻居（如LEO0301）呈现出明显的周期性通断特征，这是由于异轨卫星在纬度运动中相对距离不断变化所致。而对于层间链路（如 LEO0101、Detect0101 与上层 MEO、GEO 卫星之间的连接），也表现出鲜明的周期性通断特征。这反映了低轨卫星在高速绕地运行过程中，只能在特定的时间窗口内进入高轨节点的覆盖范围。此外，图中底部的星地链路（北京地面站 Facility\_Beijing 对 Detect 及 LEO 节点的访问）呈现出单次过境时间短、切换频率高的特点，单次可见窗口通常仅维持数分钟。
    \begin{figure}[!htb]
        \centering
        \includegraphics[width=0.8\linewidth]{fig/链路-典型链路可见性01.pdf}
        \caption{典型链路可见性分析}
        \label{fig: 典型链路可见性}
    \end{figure}
    
    （2）距离变化分析
    本文提取了仿真周期内典型观测节点（Detect0101 与 LEO0101）与其他各轨道层节点及地面站之间的距离变化曲线，如图\ref{fig: 典型链路距离变化}所示。
    同层轨道面内链路（如 Detect0101-To-Detect0102/0110）的物理距离展现出极高的稳定性，其数值长期维持在固定范围且波动极小，是网络中最为稳固的传输路径。相比之下，同层异轨链路（如 Detect0101-To-Detect0201/1001 等）则呈现出高频率、大幅度的周期性波动，反映了异轨卫星在极区交叉与赤道扩散过程中的剧烈相对运动。对于跨层链路，受轨道高度差异影响，曲线表现出显著的分层特征：与 MEO 卫星的距离在 5000 km 至 15000 km 之间呈规则的抛物线状起伏；而与 GEO 卫星的距离则维持在 35000 km 以上的高位，且波动周期更长、曲线更为平缓。此外，星地链路（To-Facility\_Beijing）仅在有限的可见窗口内出现，表现为极短时间内距离的急剧下降与上升。
    \begin{figure}[!htb]
		\centering
		\subfloat[]{
			\includegraphics[width=0.7\linewidth]{fig/链路-Detect0101-典型链路距离变化-01.pdf}
			\label{fig: 链路-Detect0101-典型链路距离变化}
		}
		\hfill
		\subfloat[]{
			\includegraphics[width=0.7\linewidth]{fig/链路-LEO0101-典型链路距离变化-01.pdf}
			\label{fig: 链路-LEO0101-典型链路距离变化}
		}
		\caption[典型链路距离变化]{典型链路距离变化。\subref{fig: 链路-Detect0101-典型链路距离变化}Detect0101视角；\subref{fig: 链路-LEO0101-典型链路距离变化}LEO0101视角} 
      \label{fig: 典型链路距离变化}
	\end{figure}
    
    2、架构性能分析
    
    为了全面验证本文提出的分层多域管控架构（Hierarchical Domain Control Architecture，记为HDC-Arch）在解决大规模异构星座管控难题上的有效性，本节以传统的地面集中式管控架构（Ground Control Architecture，记为 GCC-Arch）和基于动态邻居关系的 LEO 分布式管控架构（LEO Distributed Control Architecture，记为 LDC-Arch）作为基线方案，从时延、稳定性及信令开销等维度展开对比分析。

    % 新写的--开始===============================================================
    图\ref{fig: 平均控制回路时延}对比了三种架构在 120 分钟仿真周期内的全网平均控制回路时延。实验结果表明，GCC-Arch受限于地面站的地理分布，时延长期维持在 180ms 以上的高位，且存在明显的覆盖盲区惩罚。LDC-Arch虽然利用星间链路将物理时延降低至 105ms 左右，但其曲线呈现出剧烈的锯齿状波动。这是因为 LEO 卫星的高速运动导致控制器频繁发生硬切换，每次切换引发的路由重算与信令交互带来了不可忽视的时延抖动。相比之下，本文提出的 HDC-Arch展现出了显著的稳定性优势。虽然受限于 GEO 层的物理距离，其平均时延略高于 LDC-Arch，但其时延波动极小。得益于虚拟编址与软切换机制，HDC-Arch 屏蔽了底层物理拓扑的高频变化，仅产生微小的状态迁移开销，消除了 LDC-Arch 中因频繁重选控制器而导致的剧烈震荡。
    \begin{figure}[!htb]
        \centering
        \includegraphics[width=0.7\linewidth]{fig/平均控制回路时延.pdf}
        \caption{三种架构的全网平均控制回路时延随时间变化}
        \label{fig: 平均控制回路时延}
    \end{figure}

    图\ref{fig: 控制信令开销}展示了三种架构在120分钟仿真周期内的全网控制信令开销对比。结果表明，本文提出的HDC-Arch在维持全网10\%实时管控的前提下，平均信令开销仅为每秒506个信令包，相比经过工程优化（引入消息聚合与骨干泛洪机制）的分布式架构LDC-Arch降低了近45\%。这种显著优势主要得益于HDC-Arch采用了虚拟节点技术，当卫星高速穿越地理网格边界时，通过继承虚拟身份保持了逻辑拓扑的恒定，从而避免了LDC-Arch因物理与逻辑归属错位而引发的全网路由重收敛和信令泛洪。此外，虽然地面集中式架构GCC-Arch的统计数值最低，但这主要是受限于地面站的地理分布，导致大量处于盲区的卫星无法上报状态，其实质是以牺牲全网可控性为代价的伪低开销。综上所述，HDC-Arch在保证管控覆盖率的同时，有效解决了大规模异构星座因高动态性带来的信令风暴问题，体现了架构设计的合理性与高效性。
    \begin{figure}[!htb]
        \centering
        \includegraphics[width=0.7\linewidth]{fig/控制信令开销.pdf}
        \caption{三种架构的控制信令开销随时间变化}
        \label{fig: 控制信令开销}
    \end{figure}

    \begin{figure}[!htb]
        \centering
        \includegraphics[width=0.7\linewidth]{fig/控制中断率.pdf}
        \caption{三种架构的控制中断率随时间变化}
        \label{fig: 控制中断率}
    \end{figure}

    图\ref{fig: 控制中断率}展示了三种架构在 120 分钟仿真周期内的控制链路中断率（Control Link Outage Rate）变化情况。该指标反映了网络中处于失控或路由不可达状态的卫星比例，直接衡量了管控系统的可靠性。从统计结果可以看出，GCC-Arch（地面集中式）的中断率最高，平均在 4.8\% 左右，且伴随明显的周期性波动。这表明尽管部署了 48 个地面站，受限于地理分布和地球自转的影响，仍有部分卫星在飞越深海或极地时会短暂脱离地面站的覆盖范围，导致管控链路间歇性中断。LDC-Arch（LEO 分布式）的中断率维持在 2.0\% 左右。虽然其物理上实现了全覆盖，但由于 LEO 卫星的高速运动，平均每分钟发生的数十次网络拓扑切换会频繁触发分布式路由的重收敛过程。在路由表更新的间隙，相关节点的链路处于瞬时不可用状态。相比之下，本文提出的 HDC-Arch 表现最为稳定，平均中断率仅为 0.51\%。HDC-Arch 成功屏蔽了物理拓扑的高频变化，实现了逻辑连接的无缝切换。虽然仍存在极微小的切换处理时延，但其链路可靠性相比 LDC-Arch 提升了近 4 倍，相比 GCC-Arch 提升了近 10 倍，能够更好地满足大规模星座对连续管控的要求。

    图\ref{fig: 三种架构业务端到端时延}展示了三种架构在典型业务链路下的平均端到端时延对比，仿真结果表明 HDC-Arch 在全场景下均具有显著优势，其平均时延较 GCC-Arch 和 LDC-Arch 分别降低了 50\%-65\% 和 30\%-50\%。在 Beijing-New York 等跨洲长途链路中，HDC-Arch 利用 MEO 卫星大幅减少物理跳数，将时延稳定控制在 80ms 以下，显著优于受限于星地回路的 GCC-Arch和受多跳累积影响的 LDC-Arch；而在 Beijing-Shanghai 等区域短途链路中，该架构通过虚拟节点机制屏蔽底层拓扑动态性，有效消除了传统架构中因切换与信令交互产生的基础时延，实现较低时延传输，验证了分层多域架构在物理路径优化与控制开销削减上的可用性。
    \begin{figure}[!htb]
        \centering
        \includegraphics[width=0.7\linewidth]{fig/三种架构业务端到端时延.pdf}
        \caption{三种架构业务的端到端时延统计}
        \label{fig: 三种架构业务端到端时延}
    \end{figure}

    
    % 新写的--结束===============================================================
    
    
 %    (1) 信令开销反映了控制平面在维护网络拓扑和下发流表时的负载压力，图\ref{fig: 架构-信令开销}展示了网络累积信令开销的增长趋势。GC-Arch呈现出陡峭的线性增长趋势,这是由于 LEO/Detect 卫星高速运动，频繁进出地面信关站的覆盖范围，导致链路建立与拆除的信令交互从未停止，网络始终处于高负载状态。HDC-Arch呈现出平缓的低斜率增长,得益于本文提出的虚拟网格映射机制，HDC-Arch有效屏蔽了卫星在同一网格内的微小运动，仅在节点跨越网格边界时触发位置更新。这种机制显著抑制了信令风暴的产生，证明了架构在大规模网络中优异的可扩展性。
 %    \begin{figure}[!htb]
 %        \centering
 %        \includegraphics[width=0.7\linewidth]{fig/架构-信令开销.png}
 %        \caption{控制信令开销随时间的变化}
 %        \label{fig: 架构-信令开销}
 %    \end{figure}
    
 %    （2）控制链路可用性，定义为整个星座中处于已建立有效控制链路的卫星数量占卫星总数的比例。图\ref{fig: 架构-控制可用性-空间}展示了控制可用性随维度的空间分布特性。 GC-SDN在南纬 30° 至 60° 区域（对应广阔的南大洋）的可用性出现断崖式下跌，最低不足 10\%。这是由于地面信关站无法部署在公海，导致该区域成为物理盲区。HDC-Arch利用 GEO 卫星的全球视场，成功填补了这一盲区，在除极点外的中低纬度区域保持了约 60\% 的稳定连接率（图中橙色填充区域即为覆盖增益）。这证明了 HDC-Arch消除了对地面设施地理分布的依赖，实现了全球管控服务的均质化。
 %    \begin{figure}[!htb]
	% 	\centering
	% 	\subfloat[]{
	% 		\includegraphics[width=0.7\linewidth]{fig/架构-控制可用性-空间.png}
	% 		\label{fig: 架构-控制可用性-空间}
	% 	}
	% 	\hfill
	% 	\subfloat[]{
	% 		\includegraphics[width=0.7\linewidth]{fig/架构-控制可用性-时间.png}
	% 		\label{fig: 架构-控制可用性-时间}
	% 	}
	% 	\caption[控制连续可用性的趋势变化]{控制连续可用性的趋势变化。\subref{fig: 架构-控制可用性-空间}随时空间变化；\subref{fig: 架构-控制可用性-时间}随时间变化} \label{fig: 架构-控制可用性}
	% \end{figure}
    
 %    图\ref{fig: 架构-控制可用性-时间}展示了 120 分钟仿真周期内控制链路可用性的时变曲线。在仿真初期，GC-Arch的控制链路可用性因卫星飞越盲区而长期徘徊在 15\% 的低位，随后在飞入地面站区域后才出现陡峭拉升，表明地面管控架构难以提供连续稳定的服务保障；而HDC-Arch 得益于 GEO 层的广域静止覆盖，有效屏蔽了地表海陆分布差异，曲线全程保持平滑演进且无断崖式下跌，随着卫星从初始位置向中低纬度移动并进入最佳视场，响应率稳步攀升并收敛至 85\% 的稳态水平，在长周期运行中维持可连续的服务能力。
    
    
 %    （3）控制回路时延定义为控制器与受控卫星节点之间完成一次信令交互所需的往返时间，图\ref{fig: 架构-控制回路时延}展示了时延概率密度分布。GC-Arch 拥有较低的平均时延44ms，因为其依赖星地链路，LEO 卫星高度低，信号传播路径短，因此具备极低的物理传输时延，但其连接频繁中断（时延趋于无穷大）；而 HDC-Arch 将主控制器部署于 GEO 轨道，控制信令需跨越 GEO-LEO 层间链路，光速传播限制导致了较高的物理时延基数，但也是为了换取广域全时覆盖所必须付出的物理代价。下文的端到端仿真将证明，这种控制层面的慢，换来的是业务层面的稳。
 %    \begin{figure}[!htb]
 %        \centering
 %        \includegraphics[width=0.7\linewidth]{fig/架构-控制回路时延.png}
 %        \caption{控制回路时延概率密度分布}
 %        \label{fig: 架构-控制回路时延}
 %    \end{figure}

 %    (4) 业务建立时延，定义为从源终端发起业务请求开始，到网络控制平面完成路由计算并下发流表，从而建立端到端传输路径所需的时间。为了评估该指标，本节通过蒙特卡洛方法生成 10000 个随机业务请求，评估架构在实际业务场景中的响应能力与服务质量。图\ref{fig: 架构-业务建立时延-累积分布}展示了两种架构下业务建立时延的累积分布函数（CDF），在低时延响应方面，SH-SDN约50\%的业务请求在300ms内通过星上分层控制器完成处理，而GC-SDN仅能处理23\%的请求。并且，SH-SDN的CDF曲线始终领先于GC-SDN，表明其更快的业务建立速度。尤其在高时延区间前，SH-SDN已完成半数业务建立，而GC-SDN仍有近80\%未完成。
 %    \begin{figure}[!htb]
	% 	\centering
	% 	\subfloat[]{
	% 		\includegraphics[width=0.7\linewidth]{fig/架构-业务建立时延-累积分布.png}
	% 		\label{fig: 架构-业务建立时延-累积分布}
	% 	}
	% 	\hfill
	% 	\subfloat[]{
	% 		\includegraphics[width=0.7\linewidth]{fig/架构-业务建立时延-空间.png}
	% 		\label{fig: 架构-业务建立时延-空间}
	% 	}
	% 	\caption[业务建立时延分析]{业务建立时延分析。\subref{fig: 架构-业务建立时延-累积分布}业务建立时延累积分布图；\subref{fig: 架构-业务建立时延-空间}业务建立时延随空间变化趋势} \label{fig: 架构-业务建立时延}
	% \end{figure}
 %     图\ref{fig: 架构-业务建立时延-空间}刻画了平均业务建立时延随纬度的变化情况。GC-Arch 表现出显著的地理不平衡性，其时延随纬度降低呈阶梯状上升，在南极区域达到最大值，反映了集中式架构对地面基站部署位置的高度依赖。相比之下，HDC-Arch 在全纬度范围内的平均时延较基准方案降低了约 $40\% \sim 60\%$，且曲线波动更为平缓。这说明分级管控架构利用星间协同感知与分布式处理机制，显著增强了异构卫星网络对全球尺度业务的均衡支撑能力，有效缓解了因地面基础设施分布不均导致的接入性能差异。
    
    \section{本章小结}
    本章针对异构星座网络面临的信令风暴与地理盲区两大核心痛点，设计了基于分层多域管控架构的组网管控体系。首先，提出了GEO全局主控加LEO域内协同的分级管控策略，利用GEO/MEO/LEO的轨道差异优势，实现了控制功能的合理分级与负载分担。其次，设计了基于虚拟节点的区域构建策略与地址映射机制，将高动态的物理拓扑抽象为逻辑静止的虚拟网络，有效屏蔽了底层卫星的高速运动特性。同时，提出了基于虚拟地址封装的异构互联机制，解决了物理地址与逻辑位置解耦的跨层通信难题。最后，从多个维度对比了HDC-Arch与传统地面控制架构的性能。结果表明，HDC-Arch在信令开销降低、全球覆盖均衡性提升以及业务响应实时性方面均表现出显著优势，尤其在低时延业务建立方面，约50\%的请求可在300ms内完成处理，上述分析验证了架构的理论优越性，具体的工程实现与协议栈性能将在第五章阐述。本章提出的组网架构与管控机制为大规模异构卫星网络提供了稳定的逻辑底座，有效规避了物理拓扑剧变带来的管控难题，这为后续第四章开展基于多维QoS约束的复杂路由算法研究奠定了坚实的架构基础。

    \chapter{异构星座多域路由算法研究}
    % 参考注意查重、抄袭
    % 《基于改进遗传算法的多约束QoS路由算法研究-葛君伟》
    % 《LEO卫星网QoS遗传算法路由协议-张雪东》
    % 《基于多目标约束遗传算法的 SDN 路径增强算法-周睿》
    % 《LEO卫星星座QoS路由问题研究-赵祥（我的变异规则参考它的）》
    % 《基于智能优化的低轨卫星网络路由算法研究-郝苗苗》
    \section{引言}
    在第三章中，本文针对异构星座网络规模巨大、拓扑高动态及多维资源受限的特点，设计了GEO全局主控加LEO域内协同的分层多域管控架构。该架构通过引入虚拟节点与SDN技术，实现了控制平面与数据平面的解耦，以及全局策略制定与局部快速执行的分离。然而，架构的搭建仅为网络管控提供了基础平台，如何在这一架构之上，设计高效的路由机制以保障端到端业务QoS，仍面临严峻挑战。
    
    首先，LEO卫星的高速运动导致星间链路频繁通断，物理拓扑的分钟级重构要求路由算法具备极高的收敛速度与鲁棒性。其次，随着未来天地一体化网络承载业务的多样化，不同业务（如时延敏感的应急通信、带宽敏感的遥感数据回传）对网络资源的需求存在显著差异，单一的路由度量指标已无法满足多维QoS约束。最后，由于卫星节点分布不均及地面信关站的地域限制，网络流量极易在局部区域形成漏斗效应，导致部分链路拥塞甚至瘫痪，这对路由策略的负载均衡能力提出了更高要求。

    针对上述问题，本章提出一种面向多域动态QoS的异构星座路由策略。该策略充分利用分层架构优势，采用宏观规划与微观寻优相结合的思路：在域间层面，利用GEO控制器的全局视野基于抽象拓扑进行跨域路径规划与QoS预算分配；在域内层面，设计基于方向引导与拥塞感知的改进遗传算法（Improved Genetic Algorithm, IGA），在满足多维QoS约束的前提下，实现对拥塞路径的主动规避与网络流量的及时调度。

    \section{异构星座多域网络 QoS 路由问题建模}
    为了屏蔽低轨卫星网络拓扑的高动态性并降低路由计算复杂度，本节基于第三章构建的虚拟地面虚拟管理域的划分，建立基于虚拟节点的分层图模型，并针对多业务需求构建多约束QoS路由优化模型。
    \subsection{基于虚拟节点的异构星座分层图模型构建}
    \label{sec: 异构星座分层图模型}
    为了在屏蔽低轨卫星高动态特性的同时，保留异构网络层间协作与流量状态的实时性，本文基于第三章提出的虚拟区域划分策略，构建了一个包含逻辑节点与物理节点的异构分层图模型。该模型由GEO主控制器维护，作为域间路由规划的全局视图基础。
    
    定义异构星座网络的拓扑图为 $G(t) = (V, E(t), W(t))$，其中 $t$ 表示离散的时间片索引。
    
    1、节点集合 $V$
    
    针对GEO、MEO、LEO三层卫星在覆盖范围与运动特性上的显著差异，节点集合 $V$ 采用虚实结合的构建方式，由以下三部分组成：
    \begin{equation}
        V = V_{LEO}^{vir} \cup V_{MEO}^{phy} \cup V_{GEO}^{phy}
    \end{equation}
    
    LEO虚拟节点集合 ($V_{LEO}^{vir}$)：针对数量庞大且高速运动的LEO卫星，采用逻辑网络与物理网络解耦的设计思想。将全球覆盖区域划分为 $N$ 个逻辑网格，每个网格抽象为一个静止的虚拟节点 $v_{leo}^i (i=1, \dots, N)$。物理LEO卫星与虚拟节点之间存在动态映射关系，当卫星飞越某区域时，自动继承该虚拟节点的逻辑地址信息。这种映射有效屏蔽了LEO同层内的物理拓扑变化。
    
    MEO物理节点集合 ($V_{MEO}^{phy}$)：MEO卫星作为关键中继层，其数量相对较少且轨道较高。为了精确计算层间链路的几何可见性，直接将其建模为物理节点 $v_{meo}^j$。
    
    GEO物理节点集合 ($V_{GEO}^{phy}$)：GEO卫星相对地面静止，既是存储全局拓扑 $G(t)$ 的控制节点，也可以参与数据转发的物理转发节点 $v_{geo}^k$。将其显式纳入图模型，可为高可靠性要求或时延不敏感业务提供稳定的兜底路径。

    2、边集合 $E(t)$
    
    为了准确描述不同性质链路的拓扑特征，本文依据连接节点的属性，将边集合 $E(t)$ 划分为虚边与实边两类：
    \begin{equation}
        E(t) = E_{vir}(t) \cup E_{phy}(t)
    \end{equation}
    
    (1) 虚边集合 ($E_{vir}$)
    
    定义凡是连接端点中包含LEO虚拟节点的链路统称为虚边。此类边描述的是逻辑上的连通性，其物理载体随卫星运动而更替。包含：a）域间虚边 ($v_{leo}^i \leftrightarrow v_{leo}^j$)，连接相邻的LEO虚拟节点。逻辑连接相对静止，物理上对应当前时刻覆盖这两个区域的卫星间的星间链路（ISL）。b）层间接入虚边 ($v_{leo}^i \leftrightarrow v_{meo/geo}^k$)，连接静止的LEO虚拟节点与运动的MEO/GEO物理节点。此类边表示当前区域内的LEO卫星与上层卫星建立的层间链路。
    
    (2) 实边集合 ($E_{phy}$)
    
    定义连接两个物理节点（MEO或GEO）的真实链路为实边，此类边基于视距可见性建立。包括连接MEO与GEO卫星的层间链路($v_{meo} \leftrightarrow v_{geo}$)，以及GEO或MEO星座内部的星间链路 ($v_{geo} \leftrightarrow v_{geo}, v_{meo} \leftrightarrow v_{meo}$)。
    
    3、边权重集合$W(t)$
    
    针对上述两类边，权重的计算机制存在区别。定义边 $e$ 在时刻 $t$ 的权重向量为 $w_e(t) = [B_e(t), D_e(t), L_e(t)]$，分别表示该边的带宽、时延、丢包率。
    
    (1) 针对虚边 ($e \in E_{vir}$)，此类边连接了至少一个虚拟节点，其物理载体通常由多条底层物理链路构成（链路集合记为 $\mathcal{L}_e$），需采用状态聚合机制计算综合权重：
    
    带宽聚合：取集合 $\mathcal{L}_e$ 中的最小值（瓶颈带宽），即 $B_e(t) = \min_{l \in \mathcal{L}_e} B_l(t)$，以反映该逻辑通道的最差传输能力。
    
    时延聚合：取集合 $\mathcal{L}_e$ 的平均值，即 $D_e(t) = \frac{1}{|\mathcal{L}_e|} \sum_{l \in \mathcal{L}_e} D_l(t)$，以反映区域间的平均连通水平。
    
    丢包率聚合：取集合 $\mathcal{L}_e$ 中的最大值，即 $L_e(t) = \max_{l \in \mathcal{L}_e} L_l(t)$，采用保守策略以确保QoS的可靠性上限。
    
    (2) 针对实边 ($e \in E_{phy}$)，此类边连接两个物理节点，对应真实存在的单条链路，采用直接测量，权重直接来源于控制器采集的实时物理状态。

    综上所述，通过对节点维度的虚实结合、链路维度的分类构建以及权重维度的差异化计算，定义了异构星座分层图模型 $G(t)$ ， 该模型示意图如图\ref{fig: 异构星座分层图模型构建}所示。
    
    顶层的GEO卫星与中间的MEO卫星通过黑色的实边构建起网络，对应模型中的 $E_{phy}$ 集合。底层蓝色网格平面代表逻辑划分的LEO虚拟区域，绿色虚线表示连接相邻区域的域间虚边，表示逻辑拓扑的静态特征，橙色点划线代表连接底层虚拟节点与上层物理节点的层间接入虚边。
    \begin{figure}[!htb]
        \centering
        \includegraphics[width=0.8\linewidth]{fig/异构星座分层图模型构建.pdf}
        \caption{异构星座分层图模型示意图}
        \label{fig: 异构星座分层图模型构建}
    \end{figure}

    \subsection{业务QoS度量和路由约束模型}
    \label{sec: 业务QoS度量和路由约束模型}
    为了对异构星座中的多约束路由问题进行形式化分析，本文将网络设备间的连接抽象为图的路径。基于 前文构建的图模型 $G(V, E)$，假设源节点为 $s$，目的节点为 $t$，多约束QoS路由问题可描述为：在图 $G$ 中寻找一条从 $s$ 到 $t$ 的路径 $P(s,t)$，使得该路径上所有链路的组合权值满足特定业务的QoS约束。
    
    1、 QoS度量参数的分类特性
    
    网络中的QoS参数依据其在路径上的合成规则，可分为可加性、凹性和可乘性三类。不同类型的参数具有不同的数学特性及处理策略。
    
    （1）可加性度量参数
    
    此类参数的路径总权值等于路径上所有组成链路的参数值之和，其典型代表为时延、跳数。由于寻找满足多个可加性约束的路径通常属于NP-Complete问题，也是QoS路由算法优化的核心难点，因此在模型中需通过累积计算进行严格约束：
    \begin{equation}
        w(P) = \sum_{e \in P} w(e)
    \end{equation}
    
    （2）凹性度量参数
    
    此类参数的路径总权值由路径上所有链路参数值的最小值（即瓶颈值）决定，其典型代表为带宽。凹性参数具有优良的筛选特性，在路由计算前，可以通过拓扑剪枝预先剔除网络中不满足带宽需求（$w(e) < B_{req}$）的链路，从而显著降低算法的搜索空间：
    \begin{equation}
        w(P) = \min_{e \in P} w(e)
    \end{equation}
    
    （3）可乘性度量参数
    
    此类参数的路径总权值等于路径上所有链路参数值的乘积，其典型代表为丢包率或可用性。在实际计算中，为了降低复杂度，通常采用取对数的方式将其转化为可加性参数；或者在低丢包率网络环境（$L \ll 1$）下，利用近似公式（$1-\prod(1-L_e) \approx \sum L_e$）将其直接作为可加性参数处理：
    \begin{equation}
        w(P) = \prod_{e \in P} w(e) \approx \sum_{e \in P} w(e)
    \end{equation}

    2、基于业务特征向量的路由约束
    
    结合上述度量特性，本文考虑带宽、时延、丢包率这三个典型的约束条件。定义任意业务请求 $R_k$ 为一组特征参数集合：$R_k = \{ B_{req}^{(k)}, D_{req}^{(k)}, L_{req}^{(k)}, \omega_{pri}^{(k)} \}$。其中，$B_{req}$ 、$D_{req}$ 、$L_{req}$ 分别为业务对带宽、时延、丢包率的需求。有效路径 $P(s, d)$ 必须满足以下三个不等式约束。
    
    (1) 瓶颈带宽约束
    
    路径上所有链路的剩余可用带宽 $B_e(t)$ 需满足业务最小带宽需求。在算法实现中，利用凹性特征，凡是不满足 $B_e(t) \ge B_{req}^{(k)}$ 的链路将被视为断开：
    \begin{equation}
        \min_{e \in P(s,d)} B_e(t) \ge B_{req}^{(k)}
    \end{equation}
    
    (2) 累积时延约束
    
    路径端到端总时延需小于业务容忍上限：
    \begin{equation}
        \sum_{e \in P(s,d)} D_e(t) \le D_{req}^{(k)}
    \end{equation}
    
    (3) 可靠性约束
    
    路径的端到端丢包率需满足业务可靠性要求。采用加性近似处理：
    \begin{equation}
        \sum_{e \in P(s,d)} L_e(t) \le L_{req}^{(k)}
    \end{equation}
    
    综上所述，通过对QoS参数的分类处理，可以将复杂的多约束问题进行降维，带宽约束通过预处理解决，丢包率约束转化为加性形式，从而将原问题简化为受带宽限制（剪枝后）的最小代价路径问题。

    \subsection{包含拥塞惩罚项的优化目标函数}
    
    在满足\ref{sec: 业务QoS度量和路由约束模型}节所述硬性 QoS 约束的可行路径集合中，路由算法的核心任务是寻找综合代价最小的最优路径。为了解决传统线性代价函数对高负载链路不敏感导致的网络拥塞问题，并为后续遗传算法提供具备拥塞规避导向的评价标准，本文采用一种包含非线性拥塞惩罚项的全局优化目标函数。综合路径代价函数 $Cost(P)$ 定义如下：
    \begin{equation}
        Cost(P) = \underbrace{ \left( \alpha \cdot \frac{D(P)}{D_{max}} + \beta \cdot \frac{L(P)}{L_{max}} \right) }_{\text{基础QoS代价}} \times \underbrace{ \left( 1 + \lambda \cdot e^{\gamma \cdot \rho_{max}(P)} \right) }_{\text{非线性拥塞惩罚因子}}
    \end{equation}
    
    其中，基础QoS代价该部分反映路径在时延和丢包率方面的基础质量：$D(P)$ 和 $L(P)$ 分别为路径的端到端累积时延和丢包率，$D_{max}$ 和 $L_{max}$ 为归一化参考常数。$\alpha$ 和 $\beta$ 为权重系数（$\alpha + \beta = 1$），根据 \ref{sec: 业务QoS度量和路由约束模型}节中业务特征向量的优先级参数 $\omega_{pri}^{(k)}$ 动态调整。例如，对于时延敏感业务，增大 $\alpha$ 的值以偏向低时延路径。
    
    
    非线性拥塞惩罚因子利用指数函数特性，对接近满载的路径施加高额惩罚。其中，$\rho_{max}(P)$ 表示路径瓶颈负载率，定义为路径上所有链路负载率的最大值，即 $\rho_{max}(P) = \max_{e \in P} (\frac{Traffic(e)}{Capacity(e)})$。在域间路由中，$\rho$ 对应虚拟链路的聚合负载率，指导跨域路径绕开热点区域。在域内路由中，$\rho$ 对应物理链路的实际负载率，指导遗传算法规避拥塞的LEO节点。$\lambda$ 为拥塞惩罚系数，用于调节非线性惩罚项在综合代价中的权重幅值，决定了算法对网络拥塞的规避强度。$\gamma$ 为拥塞敏感度因子，该参数控制惩罚增长的陡峭程度。当 $\rho_{max}(P)$ 较低时，指数项值较小，代价主要取决于基础QoS。当 $\rho_{max}(P)$ 升高并接近拥塞阈值（如 0.8）时，指数项 $e^{\gamma \cdot \rho_{max}}$ 迅速增大，导致总代价 $Cost(P)$ 剧增。
    
    此优化目标函数使得任何包含严重拥塞链路的路径都会产生极高的代价值。在后续章节设计的改进遗传算法（IGA）中，该函数将直接转化为适应度函数，配合拥塞感知变异策略，迫使种群向低负载、高QoS的区域进化，尽量保证QoS保障与负载均衡的统一。

    \section{异构星座多域路由总体框架}
    基于第三章提出的分级多域管控架构及前文构建的异构星座分层图模型，本节提出一种分层规划与协同执行的路由机制。该机制旨在解决大规模异构星座中全网最优性与局部实时性之间的矛盾，通过将复杂的端到端路由问题分解为域间宏观规划与域内精细选路两个子问题，在降低控制平面计算复杂度的同时提升网络响应速度。

    1、协同路由功能架构
    
    图\ref{fig: 多域协同路由功能架构图}展示了协同路由功能架构，该架构包含全局路由规划层、域内路由决策层与数据转发层，这三个层次通过纵向的路由信令交互与横向的数据流转处理，共同构成了一个立体化的路由管控体系。
    \begin{figure}[!htb]
        \centering
        \includegraphics[width=0.8\linewidth]{fig/多域协同路由功能架构图.pdf}
        \caption{多域协同路由功能架构图}
        \label{fig: 多域协同路由功能架构图}
    \end{figure}
    
    逻辑上位于路由架构顶端的是域间路由规划层，该层部署于 GEO 卫星节点，其核心任务是构建并维护全网的抽象虚拟拓扑$G(t)$。利用 GEO 卫星的广域覆盖特性，该层通过虚拟化技术屏蔽底层物理节点的微观状态波动，仅关注逻辑域之间的连接关系与聚合状态。基于这种全局视野，该层利用域间路由算法求解由虚拟节点组成的跨域路径序列，并引入QoS 预算分解策略，将端到端的总时延、丢包率等硬性约束，根据各域的聚合状态能力，拆解为各域必须满足的局部 QoS 预算，为下层的域内路由决策提供了明确的优化边界与约束条件。
    
    域内路由决策层位于中间层，部署于各 LEO 域控制器，其核心职能是通过域内状态感知机制，实时监控本域内的物理拓扑变化 $G_{phy}^{(i)}$ 与链路负载情况。在接收到上层下发的跨域路径指令与局部 QoS 预算后，该层启动域内路径寻优过程，利用改进遗传算法在域内的物理拓扑中进行精细化搜索，旨在寻找能够满足上层 QoS 预算并主动规避拥塞链路的物理路径片段，从而完成从逻辑路径到物理路径的映射。
    
    最底层为数据转发层，由全体 LEO/MEO 卫星节点构成，构成了实际承载业务流量的数据平面。为了最大程度提升数据传输效率，该层剥离了复杂的路由计算功能，专注于数据包的转发处理。各物理节点通过南向接口接收决策层下发的流表生成规则，并基于流表执行匹配与动作操作，实现业务数据在星间链路上的高效传输。

    2、协同路由流程
    
    协同路由过程由业务请求触发，具体流程如图\ref{fig: 跨域协同路由时序图}所示，主要分为四个阶段。

    阶段一为业务接入与判决，当源终端产生通信需求时，首先向其所属的接入域控制器发送携带QoS向量的业务请求。接入域控制器接收请求后，利用本地缓存的拓扑信息进行域内判决。若目的节点位于本逻辑域内，则直接触发域内路由算法完成路径建立；若检测到目的节点跨域，则将路由请求转发至 GEO 主控制器，启动跨域协作流程。

    阶段二为宏观规划与预算分解， GEO 主控制器作为全局核心，获取各域上报的虚拟拓扑快照。通过运行域间路由算法，GEO主控在宏观维度计算跨域路径，并根据端到端 QoS 约束进行预算分解，随后同步向相关的域下发局部 QoS 预算指标。
    
    阶段三为微观执行与流表分发，各域控制器在接收到局部预算后，执行域内路由算法，在满足局部预算的条件下计算各自域内的最优物理路径。路径计算就绪后，各控制器向 GEO 主控反馈就绪确认，之后 GEO 控制器统一下发流表至相关节点，并通知源终端路径已成功建立。
    
    阶段四为数据传输，在控制平面完成端到端路径构建及流表配置后，业务进入数据平面。源终端按照预设路径，通过建立好的星间链路跨越多个逻辑域，实现与目的终端的数据传输。
    \begin{figure}[!htb]
        \centering
        \includegraphics[width=0.85\linewidth]{fig/跨域协同路由时序图.pdf}
        \caption{跨域协同路由时序图}
        \label{fig: 跨域协同路由时序图}
    \end{figure}
    
    \section{异构星座多域网络域间路由}
    1、域间路由算法设计
    
    全局路由规划层的核心任务是计算出域间路径。不同于底层物理网络的快速变化，域间路由计算基于域间抽象拓扑视图，如图 \ref{fig: 域间抽象拓扑视图} 所示。该视图是对异构星座分层图模型 $G(t)$ 的扁平投影，路由计算涉及的节点数量被压缩至 $N$ 个虚拟区域与 $M$ 个高层物理节点，降低了网络规模与计算复杂度。
    \begin{figure}[!htb]
        \centering
        \includegraphics[width=0.9\linewidth]{fig/域间抽象拓扑视图.pdf}
        \caption{域间抽象拓扑视图}
        \label{fig: 域间抽象拓扑视图}
    \end{figure}
    
    
    针对 \ref{sec: 业务QoS度量和路由约束模型} 节定义的多维 QoS 约束，本文设计了一种基于多约束剪枝的 Dijkstra 算法。该算法通过预先剔除不满足带宽要求的链路，并在剩余拓扑中寻找综合代价最小的逻辑路径，从而实现跨域资源的全局优化。域间路由算法伪代码如算法\ref{alg: 域间路由算法}所示，执行逻辑如下。
    
   第 1至6 行对应拓扑剪枝阶段。在时刻 \( t \)，GEO控制器获取全局拓扑状态，首先根据业务请求 \( R_k \) 的瓶颈带宽需求 \( B_{req} \) 对图 \( G(t) \) 进行预处理。初始化有效边集合 \( E_{valid} \) 为空，遍历图中的所有边 \( e \)，若某条边的剩余带宽 \( B_e(t) < B_{req} \)，则将其忽略；反之，若 \( B_e(t) \ge B_{req} \)，则将其加入 \( E_{valid} \)。这一步利用带宽的凹性特征，确保了后续搜索空间中的所有路径在带宽维度上均是可行的。
    
    第 7至9 行执行拥塞感知权值计算。基于剪枝后的有效边集合 \( E_{valid} \)，利用\ref{sec: 业务QoS度量和路由约束模型}节定义的包含拥塞惩罚项的优化目标函数，计算集合中每一条边 \( e(u,v) \) 的综合权值 \( W_{uv} \)。该权值 \( W_{uv} \) 将时延、丢包率等QoS指标与链路负载率 \( \rho_e \) 耦合为一个标量，作为Dijkstra算法中的距离度量。
    \begin{algorithm}[!htbp]
      % 标题
      \caption{域间路由算法}
      \label{alg: 域间路由算法}
    
      % 输入输出定义
      % 注意：根据您的模板，\Input 和 \Output 已经定义好了
      \Input{
      1) $t$ 时刻的全局混合时变图 $G(V, E)$；\\
      2) 业务请求 $R_k = \{ B_{req}, D_{req}, L_{req}, \omega_{pri} \}$；\\
      3) 源虚拟节点 $v_s$，目的虚拟节点 $v_d$；\\
      4) 权重参数 $\alpha, \beta, \lambda, \gamma$。
      }
      \Output{最优跨域路径 $P_{opt}$。}
    
      % --- 步骤 1: 拓扑剪枝 ---
      \tcc{1. 拓扑剪枝}
      初始化有效边集合 $E_{valid} \leftarrow \emptyset$\;
      \For(\tcp*[f]{遍历图中每一条边}){每一条边 $e(u, v) \in E$}{
      \eIf{边剩余带宽 $B_e(t) \ge B_{req}$}{
      $E_{valid} \leftarrow E_{valid} \cup \{e\}$ \tcp*[f]{保留满足带宽需求的边}
      }{
      \textbf{continue} \tcp*[f]{剪除带宽不足的链路}
      }
      }
    
      % --- 步骤 2: 权值计算 ---
      \tcc{2. 拥塞感知权值计算}
      \For{有效集合中的每一条边 $e \in E_{valid}$}{
      利用公式计算包含非线性惩罚的综合权值：\\
      $W_{uv} \leftarrow \left( \alpha \frac{D_e}{D_{max}} + \beta \frac{L_e}{L_{max}} \right) \cdot (1 + \lambda e^{\gamma \cdot \rho_e})$\;
      }
    
      % --- 步骤 3: 初始化 ---
      \tcc{3. 算法初始化}
      \For{图中每一个节点 $v \in V$}{
      $dist[v] \leftarrow \infty$, $prev[v] \leftarrow null$\;
      }
      $dist[v_s] \leftarrow 0$\;
      初始化优先队列 $Q$，并将 $(v_s, 0)$ 入队\;
    
      % --- 步骤 4: 路径搜索 ---
      \tcc{4. 最短路径搜索}
      \While{优先队列 $Q$ 非空}{
      取出当前距离最小的节点 $u \leftarrow Q.pop()$\;
    
      \If(\tcp*[f]{已到达目的节点}){$u == v_d$}{
      \textbf{break}\;
      }
    
      \For{$u$ 在 $E_{valid}$ 中的每一个邻居节点 $v$}{
      计算备选路径代价 $alt \leftarrow dist[u] + W_{uv}$\;
      \If{$alt < dist[v]$}{
      $dist[v] \leftarrow alt$\;
      $prev[v] \leftarrow u$\;
      $Q.update(v, alt)$\;
      }
      }
      }
    
      % --- 步骤 5: 回溯 ---
      \tcc{5. 路径回溯与输出}
      根据前驱数组 $prev[\cdot]$ 从 $v_d$ 回溯至 $v_s$，构建路径序列 $P_{opt}$\;
      \textbf{return} $P_{opt}$\;
    
    \end{algorithm}
    
    第 10至25 行为最短路径的搜索与回溯过程。以源虚拟节点 \( v_s \) 为起点，目的虚拟节点 \( v_d \) 为终点，在由 \( E_{valid} \) 构成的加权子图上运行最短路算法。算法输出累积代价 \( \sum W_{uv} \) 最小的节点序列，即为最优跨域路径 \( P_{opt} \)。
    
    % 步骤4：时延与丢包率校验 TODO-待考虑
    % 尽管步骤2中的代价函数倾向于选择低时延、低丢包的路径，但Dijkstra算法本身只能保证“代价最小”，不能保证“硬约束满足”。因此，需对生成的路径 \( P_{ske} \) 进行校验：
    % \[ \sum_{e \in P_{ske}} D_e(t) \le D_{req}^{(k)} \quad \text{且} \quad \sum_{e \in P_{ske}} L_e(t) \le L_{req}^{(k)} \]
    % 若校验通过，则输出路径；若校验失败（极少情况，通常发生在网络极度拥塞时），则触发QoS协商机制或拒绝该业务请求。
    

    2、跨域QoS预算分解策略
    % 是否有重协商（Re-negotiation）或回退机制？目前的流程图似乎是一次性下发
    \label{sec: 跨域QoS预算分解策略}
    
    跨域路径调度中，GEO主控制器需将业务流的端到端QoS需求合理划分至所经各逻辑域节点，以指导各域内控制器执行路径规划。QoS预算分解的目标在于，在满足端到端约束的前提下，使每个域具备足够的资源余度以适配动态网络状态变化。本策略重点考虑最大时延、最小带宽、最大丢包率三类典型QoS指标。在跨域路径规划完成后，GEO主控制器需将业务流的端到端QoS需求合理分配至所经过的每个逻辑域，以确保各域内的局部路径调度具备可行性，并具备一定的鲁棒性。不同QoS指标具有不同的分配特性，具体策略如下。
    
    (1) 时延预算分配
    
    延迟为典型的可加性指标。假设抽象骨干路径上包含 $K$ 个逻辑域，业务的端到端最大时延预算为 $D_{req}$。为避免将跨域传输开销误计入某一域的本地预算，本方案将全路径的时延开销划分为跨域传输与域内规划两部分：
    \begin{equation}
        D_{req} = D_{cross} + \sum_{i=1}^{K} D_{local}^{(i)}
    \end{equation}
    
    其中，$D_{cross}$ 表示骨干路径上所有跨域边界边（连接不同逻辑域的链路）的传输时延总和；$D_{local}^{(i)}$ 表示待分配给第 $i$ 个逻辑域用于域内路径规划的本地时延预算。在进行实际分配时，GEO控制器首先从规划路径中识别所有属于跨域边界边的集合 $E_{cross}$，计算其物理传播时延总和作为 $D_{cross}$。随后，从业务给定的最大端到端时延预算中剔除该跨域部分，得到剩余可分配时延 $D_{rem}$：
    \begin{equation}
        D_{rem} = D_{req} - \sum_{e \in E_{cross}} d(e)
    \end{equation}
    
    若 $D_{rem} < 0$，则拒绝该请求。否则，将剩余的 $D_{rem}$ 按照路径中各域内边界边的抽象时延占比进行加权分配。设第 $i$ 个域在虚拟拓扑中的聚合时延估值为 $d_{est}^{(i)}$，则分配给该域的时延预算为：
    \begin{equation}
        D_{local}^{(i)} = D_{rem} \times \frac{d_{est}^{(i)}}{\sum_{j=1}^{K} d_{est}^{(j)}}
    \end{equation}
    即为分到第 $i$ 个域的时延预算。该策略确保了拓扑复杂或跨度较大的域能获得更多的时延规划余量。
    
    (2) 带宽预算分配
    
    带宽约束通常表现为端到端路径中的瓶颈特性（凹性指标）。设第 $i$ 个域内路径的最小可用带宽为 $B_{local}^{(i)}$，则路径所需的端到端最小带宽 $B_{req}$ 应满足：
    \begin{equation}
        B_{req} \le \min_{1 \le i \le K} \{ B_{local}^{(i)} \}
    \end{equation}
    
    因此，带宽预算不需要进行拆分。在控制器下发路径任务时，直接将 $B_{req}$ 作为各域控制器的带宽下限要求。各域仅需在本地路径规划中，确保局部路径的瓶颈带宽不低于此阈值即可。
    
    (3) 丢包率预算分配
    
    为保障传输可靠性，路径中每个域的丢包率必须受到严格控制。不同于时延的加权拆分，本策略要求路径中每个域的丢包率 $L_{local}^{(i)}$ 均不高于端到端需求 $L_{req}$，即满足：
    \begin{equation}
        L_{local}^{(i)} \le L_{req}, \quad \forall i \in \{1, \dots, K\}
    \end{equation}
    
    所以，GEO控制器在下发各域丢包率预算时，直接将 $L_{req}$ 作为上限阈值，要求各域内路由算法确保本地路径的丢包率不高于该值，从而以一种保守且鲁棒的方式满足全路径可靠性要求。

    
    \section{异构星座多域网络域内路由}
    \subsection{概述}
    在异构星座多域网络的分层路由架构中，域间路由与域内路由承担着截然不同的功能定位。GEO层的域间路由侧重于在宏观逻辑层面保障业务流跨域传输的连续性，那么LEO层的域内路由则必须深入物理层面，解决如何在具体的时变拓扑中寻找满足严苛QoS要求的可行物理路径。

    对于域内路由而言，业务流通常需要同时满足时延、带宽、丢包率及负载均衡等多维QoS指标。在图论中，这种寻找满足多个相互冲突约束条件的路径问题被称为多约束最优路径问题。理论研究已证明，当独立的加性度量约束数量超过两个时，该问题即属于NP难问题，无法在多项式时间内通过传统的精确算法求得最优解。因此，本文选用遗传算法（Genetic Algorithm, GA）作为域内路由的求解框架。遗传算法作为一种模拟生物进化机制的通用启发式搜索策略，通过种群的迭代进化来逼近最优解，适合求解此类解空间复杂、约束条件非线性的组合优化问题。
    
    然而，标准遗传算法在初始化和变异阶段通常采用完全随机的策略，在面对卫星网络时存在搜索盲目性，导致算法收敛速度慢、且容易产生包含环路的无效解。为了克服这一缺陷，本项目利用低轨卫星网络独特的类曼哈顿网格拓扑结构，构建了一套结构感知的改进遗传算法。其核心改进思想在于将拓扑的空间规律融入进化的环节中，重点在于：
    
    （1）方向引导的初始化策略。利用源宿节点的轨道与相位坐标构建几何约束，限制初始种群的生成范围，使其不再是全网漫游，而是具有明确的导向性，从而大幅提高初始解的质量。
    
    （2）拥塞感知的变异算子。针对随机变异效率低的问题，引入基于局部拓扑信息的修复机制。算法主动识别路径中的拥塞瓶颈并进行局部重构，在进化过程中实现对网络热点链路的主动规避。

    \subsection{染色体编码与方向引导初始化}
    1、物理坐标映射与序列编码

    低轨卫星网络（如Walker星座或极轨星座）虽然整体随时间高速运动，但其节点间的微观连接关系表现出高度的确定性与规律性，这为将球面拓扑抽象为二维网格提供了物理基础。具体而言，卫星节点间的星间链路主要包含两类稳定的连接特征：
    
    （1）同轨链路，同一轨道面内的卫星运行速度与方向一致，相对位置保持静止。因此，任意卫星节点 $(o, p)$ 与其前后相邻卫星 $(o, p-1)$、$(o, p+1)$ 之间的连接是永久且固定的。
    
    （2）异轨链路，相邻轨道面的卫星虽然存在相对运动，但在中低纬度区域，卫星通常与相邻轨道的同相位或近相位卫星保持稳定的可视关系。因此，节点 $(o, p)$ 通常与左右相邻轨道的卫星 $(o-1, p)$、$(o+1, p)$ 保持周期性稳定的连接。
    
    基于上述两类高度规律的连接关系，LEO网络在逻辑上呈现出规则的类曼哈顿二维栅格结构。为了兼顾遗传算法的处理效率与拓扑感知的计算需求，本文建立物理坐标与逻辑ID的双射映射机制。
    假设星座包含 $N_{orb}$ 个轨道面，每个轨道包含 $N_{sat}$ 颗卫星。定义节点逻辑ID $v \in \{1, 2, \dots, N_{total}\}$，其映射关系为：
    \begin{equation}
        v = (o - 1) \times N_{sat} + p
    \end{equation}
    
    其中，$1 \le o \le N_{orb}, 1 \le p \le N_{sat}$。通过该映射，算法既能在染色体操作中利用整数运算的高效性，又能随时反解出节点的空间坐标以计算几何距离。在此基础上，采用基于节点序列的可变长编码。一条染色体 $Ch_k$ 直接由路径上的逻辑ID序列构成：
    \begin{equation}
        Ch_k = [ v_s, v_1, v_2, \dots, v_m, v_d ]
    \end{equation}
    
    图\ref{fig: 域内拓扑视图}展示了域内拓扑视图。在该$3 \times 3$的拓扑栅格中，各节点依据公式被分配唯一的逻辑 ID $v \in \{1, 2, \dots, 9\}$。以源节点(1,1)至目的节点(3,3)的业务为例，图示蓝色线条描绘了一条典型的端到端候选路径，其经过的节点坐标序列为 $(1,1) \rightarrow (1,2) \rightarrow (2,2) \rightarrow (2,3) \rightarrow (3,3)$。基于逻辑 ID 的个体编码，该个体在算法内部被编码为可变长染色体序列 $Ch_k = [1, 2, 7, 8, 13]$。这种编码方式保留了路径的拓扑连接信息，也便于后续算子直接对路径片段进行操作。
    \begin{figure}[!htb]
        \centering
        \includegraphics[width=0.75\linewidth]{fig/域内拓扑视图.pdf}
        \caption{域内拓扑视图}
        \label{fig: 域内拓扑视图}
    \end{figure}
    
    2、基于几何距离约束的初始化策略
    
    种群初始化是遗传算法迭代的起点，在传统算法中，初始路径通常通过随机游走生成。然而，在规模庞大的卫星网络中，纯随机策略极易产生大量迂回、成环甚至不可达的劣质个体，导致算法收敛缓慢。为了提高初始种群的质量，本文结合LEO卫星网络的类曼哈顿网格结构，采用一种基于几何距离约束的启发式初始化策略。该策略利用卫星节点的轨道面与相位编号信息，引导路径生成的总体趋势向目的节点收敛。具体执行步骤如下：
    
    (1) 定义最短跳数边界。对于网络中的任意两个节点 $u$ 和 $v$，利用其逻辑坐标（轨道号 $O$，相位号 $P$）计算两点间的理论最小跳数（即曼哈顿距离）$H_{min}(u, v)$。该距离反映了理想情况下的最短路径开销。
    
    (2) 构建导向性邻居集合。在路径生成过程中，设当前节点为 $v_{curr}$，其所有物理邻居构成的集合为 $N(v_{curr})$。算法根据邻居节点相对于目的节点 $v_d$ 的距离变化趋势，将邻居划分为两类：
    
    一是趋近邻居，$N_{app} = \{ n \in N(v_{curr}) \mid H_{min}(n, v_d) < H_{min}(v_{curr}, v_d) \}$。此类节点在空间上更接近目的节点，沿着该方向前进能有效减少剩余跳数。
    
    二是发散邻居，$N_{div} = \{ n \in N(v_{curr}) \mid H_{min}(n, v_d) \ge H_{min}(v_{curr}, v_d) \}$。此类节点会导致路径背离目的方向或在同层迂回。
    
    (3) 概率导向选择。为了兼顾解的质量与种群多样性，本文采用概率导向选择机制确定下一跳节点 $v_{next}$。设定趋近概率阈值 $P_{guide}$（例如 0.8）：以 $P_{guide}$ 的概率从 趋近邻居集合 $N_{app}$ 中随机选择一个节点；以 $1 - P_{guide}$ 的概率从 发散邻居集合 $N_{div}$ 中随机选择一个节点。通过这种策略，初始种群中的大部分路径将被限制在以源、宿节点为对角的矩形区域及其周边范围内。这不仅显著剔除了无效的冗余搜索空间，保证了初始个体的高适应度，同时也保留了少量发散路径以维持种群的基因多样性，防止算法过早陷入局部最优。

    在计算节点间的距离时，利用上述映射关系反解出当前节点 $v_{curr}$ 的坐标 $(o_c, p_c)$ 与目的节点 $v_d$ 的坐标 $(o_d, p_d)$，从而计算卫星曼哈顿距离 $H_{min}$，该公式考虑卫星轨道的周期性（即0号轨道与N号轨道相邻），确保了方向引导的准确性。：
    \begin{equation}
        H_{min}(v_{curr}, v_d) = \min(|o_c - o_d|, N_{orb} - |o_c - o_d|) + \min(|p_c - p_d|, N_{sat} - |p_c - p_d|)
    \end{equation}
    
    图\ref{fig: 方向引导种群生成}所示的局部拓扑寻优过程进行具体说明。 设当前路径步进至节点 Sat 2-2，此时目的节点位于空间坐标 (5, 5)。 根据曼哈顿距离公式，当前节点到目的节点的理论最小距离为 6。在确定下一跳节点时，算法首先扫描物理邻居集合并进行空间演化趋势划分。 计算结果显示，轨道编号增加方向（Sat$_{3\text{-}2}$）或星位编号增加方向（Sat$_{2\text{-}3}$）的邻居节点均能使剩余距离由 6 降至 5，因此被归入趋近邻居集合 $N_{app}$（如图中绿色节点所示）； 反之，背离目的方向的邻居（如图中灰色节点所示）则被归入发散邻居集合 $N_{div}$。依据概率导向选择机制，以极高概率（$P_{guide}$）引导路径向绿色节点延伸，从而将初始路径的生成范围严格限制在以源、宿节点为对角的矩形几何区域内。 
    \begin{figure}[!htb]
        \centering
        \includegraphics[width=0.6\linewidth]{fig/方向引导种群生成.pdf}
        \caption{基于方向引导的种群初始化示意图}
        \label{fig: 方向引导种群生成}
    \end{figure}

    \subsection{适应度函数构建}
    适应度函数是遗传算法中引导种群进化的核心驱动力。针对多约束QoS路由问题，适应度函数的设计不仅需要量化路径的综合性能，更需有效处理多维约束条件。本文设计适应度函数方式如下。
    
    为了确保在宏观规划与微观执行层面上QoS评价标准的一致性，域内路由的目标函数应与\ref{sec: 业务QoS度量和路由约束模型}节定义的综合代价函数保持数学同构。路径 $P_k$ 的原始代价 $Cost(P_k)$ 同样由时延、丢包率及非线性拥塞项构成：
    \begin{equation}
    \label{formula: 适应度函数}
        Cost(P_k) = \left( \alpha \frac{D(P_k)}{D_{max}} + \beta \frac{L(P_k)}{L_{max}} \right) \cdot (1 + \lambda e^{\gamma \cdot \rho_{max}(P_k)})
    \end{equation}
    
    在此模型中，带宽指标虽然未直接作为加性项出现在公式中，但通过拥塞项 $\rho_{max}(P_k)$ 得到了隐式优化。最小化 $\rho_{max}$ 等效于最大化链路的剩余带宽裕度，这符合卫星网络负载均衡的实际需求。由于遗传算法遵循适应度最大化原则，需建立从代价到适应度的反比映射。定义基础适应度 $F_{base}(P_k)$ 为，其中 $\varepsilon$ 为防止分母为零的极小常数：
    \begin{equation}
        F_{base}(P_k) = \frac{1}{Cost(P_k) + \varepsilon}
    \end{equation}
    
    在多约束路由问题中，传统的强罚函数策略（即对违反任意约束的解赋予极小值或直接剔除）容易导致搜索过程中的有效基因丢失，特别是在网络资源受限、可行解空间狭窄的场景下，算法极易陷入局部最优或无法收敛。为此，本文采用一种自适应软罚函数策略。该策略区分物理硬约束与服务质量软约束，允许轻微违反 QoS 约束的降级解在种群中以一定概率存活，从而引导种群逐步向可行解边界逼近。
    
    首先，定义路径 $P_k$ 在 QoS 指标（如时延）上的相对违规度 $\Delta_{delay}$：
    \begin{equation}
        \Delta_{delay}(P_k) = \max \left( 0, \frac{D(P_k) - D_{local}}{D_{local}} \right)
    \end{equation}
    
    其中 $D_{local}$ 为根据 \ref{sec: 跨域QoS预算分解策略} 节分解得到的域内时延预算。同理可定义丢包率的违规度 $\Delta_{loss}$。基于此，构造动态惩罚系数 $\eta(P_k)$。对于物理硬约束（如带宽不足），由于其不可通过算法优化弥补，仍保持强惩罚；对于 QoS 软约束，采用反比例衰减惩罚：
    \begin{equation}
        \eta(P_k) =\begin{cases}1, & \text{若满足所有约束} \\sigma_{hard}, & \text{若 } B_{min}(P_k) < B_{req} \quad \\frac{\sigma_{soft}}{1 + \kappa \cdot (\Delta_{delay} + \Delta_{loss})}, & \text{若违反 QoS 软约束}\end{cases}
    \end{equation}
    
    % 其中，$B_{min}(P_k)$ 为路径最小剩余带宽；$\sigma_{hard}$ 为物理阻断惩罚系数（取值 $10^{-10}$），用于极大概率淘汰物理不可达路径；$\sigma_{soft}$ 为基础降级系数（本文取 0.1），$\kappa$ 为敏感度调节因子（本文取 5.0）。
    其中，$B_{min}(P_k)$ 为路径最小剩余带宽；$\sigma_{hard}$ 为物理阻断惩罚系数，用于极大概率淘汰物理不可达路径；$\sigma_{soft}$ 为基础降级系数，$\kappa$ 为敏感度调节因子。
    最终，修正后的适应度函数$Fitness(P_k)$ 定义为：
    \begin{equation}
        Fitness(P_k) = F_{base}(P_k) \cdot \eta(P_k)
    \end{equation}
    
    通过该机制，当时延或丢包率轻微超标时，路径仅会受到适度惩罚，其适应度值仍显著高于严重拥塞或物理断连的路径。这使得IGA 能够在高负载网络环境下实现尽力而为的寻路能力，在保证物理连通性的前提下，最大程度降低路由请求的阻塞率。

    \subsection{选择与交叉操作}
    在种群初始化完成后，算法通过选择操作筛选优质个体进入交配池，并利用交叉算子产生新的子代路径。由于卫星网络具有严格的拓扑约束，传统的遗传算子极易产生物理不可达的非法解。因此，本文设计了混合选择策略与双模交叉机制，以平衡算法的全局收敛速度与解的合法性。

    1、混合选择策略
    
    选择操作的主要目的是将父代中的优良基因遗传至子代种群。选择策略的有效性直接决定了算法的收敛性能。为了在保留全局最优解的同时维持种群多样性，本文采用精英保留与锦标赛选择相结合的混合策略。该策略联合亲代种群与进化过程中产生的中间种群进行筛选，具体执行流程如图\ref{fig: 混合选择策略流程图}所示。
    \begin{figure}[!htb]
        \centering
        \includegraphics[width=1\linewidth]{fig/选择策略流程图.pdf}
        \caption{混合选择策略流程图}
        \label{fig: 混合选择策略流程图}
    \end{figure}

    （1）适应度评估与排序：计算当前亲代种群中所有个体的适应度值，并按从优到劣进行排序；
    
    （2）精英保留策略：直接选择亲代种群中适应度最高的 $k$ 个精英个体，将其无损复制至下一代种群，确保最优基因不丢失；
    
    （3）锦标赛选择与进化：为了填充下一代种群的剩余空位，采用锦标赛策略。每次从亲代种群中随机抽取 $m$ 个个体，选取其中适应度最高者作为父本进入配对池。
        
    （4）种群填充：被选入交配池的父代经交叉、变异操作生成子代后，将子代加入下一代种群。重复步骤（3），直至下一代种群规模达到预设值。

    2、交叉策略
    
    低轨卫星网络呈现出的二维栅格拓扑结构，节点间的物理连接受到严格限制。传统的单点或多点交叉算子通常假设基因位可自由拼接，直接应用于卫星网络路径规划时，极易产生断链、死路或回环等非法路径。因此，本文设计的交叉操作必须在满足网络物理连通性的前提下，实现路径片段的有效重组。针对上述约束，本文设计两种互补的交叉策略，中间节点交换法与方向序列重组法。
    
    （1）中间节点交换法。该策略基于路径的拓扑重合特性，通过识别两个父代路径中的共同节点作为锚点，实现路径片段的交换。假设父代路径 $A$ 与父代路径 $B$ 均经过某一中间节点 $M$。算法以节点 $M$ 为分割点，将 $A$ 的前半段路径（源节点至 $M$）与 $B$ 的后半段路径（$M$ 至目的节点）进行拼接，从而构造出新的子代路径。如图\ref{fig: 交叉-中间节点交换法}所示，父代路径 $A$ 与 $B$ 分别为，$$P_A = [(1,1) \to (1,2) \to (2,2) \to (2,3) \to \mathbf{(3,2)} \to (4,2) \to (4,4)]$$$$P_B = [(1,1) \to (2,1) \to (3,1) \to \mathbf{(3,2)} \to (3,3) \to (3,4) \to (4,4)]$$

    算法遍历节点序列，识别出共同节点 $M(3,2)$ 作为锚点。以此为分割中心，执行双向交换操作，生成两个子代路径。
    
    (a) 子代路径 1 ($P_{child1}$)：继承 $P_A$ 的前半段与 $P_B$ 的后半段。$$P_{child1} = P_A[(1,1)\dots(3,2)] + P_B[(3,3)\dots(4,4)]$$$$= [(1,1) \to (1,2) \to (2,2) \to (2,3) \to \mathbf{(3,2)} \to (3,3) \to (3,4) \to (4,4)]$$
    
    (b) 子代路径 2 ($P_{child2}$)：继承 $P_B$ 的前半段与 $P_A$ 的后半段。$$P_{child2} = P_B[(1,1)\dots(3,2)] + P_A[(4,2)\dots(4,4)]$$$$= [(1,1) \to (2,1) \to (3,1) \to \mathbf{(3,2)} \to (4,2) \to (4,4)]$$
    
    该方法能够天然保证路径在拼接点处的连续性，有效保留了父代双方的优良结构特征。
    \begin{figure}[!htb]
        \centering
        \includegraphics[width=0.85\linewidth]{fig/交叉-中间节点交换法.pdf}
        \caption{中间节点交换法示意图}
        \label{fig: 交叉-中间节点交换法}
    \end{figure}

    （2）方向序列重组法。当父代路径间缺乏公共节点时，直接交换节点会导致断链。此时，本文采用将路径转换为方向序列的方式进行隐式交叉，并引入修复机制处理边界越界问题。首先将节点坐标序列转换为由移动方向（如东 E、南 S、西 W、北 N）组成的方向序列。在方向序列的指定位置切分并交换子串，生成新的方向序列，最后再将其反解为坐标路径。若重组后的方向序列导致路径越出网络边界或无法到达终点，则启动目标导向修复策略，即截断非法部分，按最短曼哈顿距离补齐通往目的节点的链路。
    如图\ref{fig: 交叉-方向序列重组法}所示，假设存在两个父代路径的方向序列：$$D_A = [S, S, E, E, E, S]$$$$D_B = [E, S, S, S, E, E]$$设定切分点为索引 2，将同时生成两个子代路径。
    
    (a) 子代路径1 ,保留 $D_A$ 的前缀与 $D_B$ 的后缀，得到初始序列 $[S, S, S, S, E, E]$。基于源节点 $(1,1)$ 推演，当序列执行到第 4 步时到达 $(5,1)$，超过最大轨道编号 4。算法在越界前一节点 $(4,1)$ 处截断，并计算其到目的节点 $(4,4)$ 的曼哈顿距离偏差，需向东移动 3 步。通过插入修复序列 $[E, E, E]$，生成合法路径：$$P_{child1} = [(1,1) \to (2,1) \to (3,1) \to (4,1) \xrightarrow{\text{修复}} (4,2) \to (4,3) \to (4,4)]$$
    
    (b) 子代路径 2 ($P_{child2}$) 推演：保留 $D_B$ 的前缀 ($[E, S]$) 与 $D_A$ 的后缀 ($[E, E, E, S]$)，得到初始序列 $[E, S, E, E, E, S]$。路径推演至第 5 步时到达节点 $(2,5)$，超过最大相位编号 4。算法在节点 $(2,4)$ 处触发边界修复机制。此时当前节点 $(2,4)$ 与目的节点 $(4,4)$ 的偏差为向南 2 步，因此截断后续非法序列并插入 $[S, S]$，生成合法路径：$$P_{child2} = [(1,1) \to (1,2) \to (2,2) \to (2,3) \to (2,4) \xrightarrow{\text{修复}} (3,4) \to (4,4)]$$通过上述双向交叉与修复机制，算法在一次操作中成功产生了两个结构差异显著且物理可达的新个体，有效提升了种群的搜索效率。
    \begin{figure}[!htb]
        \centering
        \includegraphics[width=0.85\linewidth]{fig/交叉-方向序列重组法.pdf}
        \caption{方向序列重组法示意图}
        \label{fig: 交叉-方向序列重组法}
    \end{figure}
    
    无论是采用何种交叉策略，生成的子代路径均需经过合法性校验，包括：邻接有效性，路径中相邻节点对必须在卫星网络拓扑中存在实际物理链路；无环性，路径序列中不得包含重复节点；连通性，路径必须完整覆盖从源节点至目的节点的全程。

    \subsection{变异操作}
    变异操作是维持种群基因多样性、防止算法早熟收敛的关键机制。然而，传统遗传算法通常采用固定的变异概率和随机节点重置策略。在多约束卫星网络环境中，这种缺乏状态感知的变异方式存在盲目性与拓扑依赖性：一方面，随机扰动极易破坏已收敛的低负载路径，或生成的子代路径反而加剧了局部拥塞；另一方面，基于规则网格几何（如平行四边形法则）的变异算子在节点故障导致拓扑不规则时，容易生成物理不可达的非法路径。针对上述问题，本节采用一种基于拥塞感知的自适应局部重路由变异策略。该策略引入链路负载作为反馈信号，实现对拥塞路径的主动规避与修复。
    
    1、拥塞度量与自适应触发机制
    
    为了在保持种群多样性与修复劣质个体之间取得平衡，本策略摒弃固定的变异概率，构建基于Sigmoid函数的自适应概率模型。首先，定义个体 $k$ 的路径瓶颈度。设路径 $P_k$ 包含链路集合 $E_k$，其最大拥塞度 $\rho_{max}(k)$ 定义为路径上所有链路负载率的峰值：
    \begin{equation}
        \rho_{max}(k) = \max_{e \in E_k} \left( \frac{L_{traffic}(e)}{C_{total}(e)} \right)
    \end{equation}
    
    其中，$L_{traffic}(e)$ 为链路当前承载的总流量，$C_{total}(e)$ 为链路物理带宽容量。基于该度量，个体 $k$ 的变异概率 $P_{mut}(k)$ 服从以下自适应分布：
    \begin{equation}
        P_{mut}(k) = P_{base} + \frac{P_{max} - P_{base}}{1 + exp(-\eta \cdot (\rho_{max}(k) - \rho_{th}))}
    \end{equation}
    
    式中各项参数定义如下。$P_{base}$为基准变异率，对低负载路径保留微小的随机变异概率，防止种群早熟。$P_{max}$为最大变异率，当路径严重拥塞时，强制提高变异概率以触发修复机制。$\rho_{th}$为拥塞阈值，当路径负载超过此阈值时，变异概率显著上升。$\eta$ 为敏感度系数，用于控制概率曲线的陡峭程度，$\eta$ 值越大，对拥塞状态的反馈越灵敏。
    
    2、执行流程
    
    当个体 $k$ 被选中进行变异时，若 $\rho_{max}(k) \ge \rho_{th}$，则执行拥塞规避操作。该操作旨在局部拓扑中构建无拥塞旁路，具体步骤如下。
    
    (1) 遍历路径 $P_k$，定位负载率最高的链路 $l_{hot} = (u, v)$。以该链路为中心，向源节点方向回溯 $m$ 跳确定起始锚点 $v_{start}$，向目的节点方向延伸 $n$ 跳确定终止锚点 $v_{end}$（通常取 $m=n=1$，以限制局部搜索范围）。
    
    (2) 构建局部子拓扑 $G_{temp}$。在原拓扑 $G$ 的基础上，将热点链路 $l_{hot}$ 及其关联拥塞节点标记为不可达状态（即从图中移除或赋予无穷大权值）：$G_{temp} = G \setminus \{l_{hot}\}$，使得后续的寻路算法在规划时必须绕开当前的拥塞区域。
    
    (3) 在处理后的拓扑 $G_{temp}$ 中，搜索从 $v_{start}$ 到 $v_{end}$ 的备选子路径集合。考虑到变异操作在进化迭代中的高频触发特性，为降低计算复杂度并加速局部收敛，此处不复用\ref{sec: 业务QoS度量和路由约束模型}节复杂的非线性全局代价函数，而是采用轻量级线性代价度量：
    \begin{equation}
        Cost_{local}(e) = \frac{1}{B_{res}(e)} + \omega \cdot D(e)
    \end{equation}
    % 引入时延项旨在施加几何拓扑约束，防止局部重构产生过长的迂回路径；舍弃丢包率项是基于丢包与拥塞的强相关性，通过最大化剩余带宽已能隐式降低排队丢包风险，避免引入随机统计指标增加搜索开销。
    
    其中，$B_{res}(e)$ 为链路剩余带宽，$D(e)$ 为链路传播时延，$\omega$ 为权重系数。依据上述代价度量择优选取一条作为 $P_{sub}^{new}$。
    
    (4) 利用新生成的无拥塞子路径 $P_{sub}^{new}$ 替换原路径中从 $v_{start}$ 到 $v_{end}$ 的旧片段，形成新个体 $P'_k$。最后对新路径进行连通性校验，确保其满足物理可达性约束。
    
    通过上述机制，变异算子不再是单纯的随机扰动，而是转化为具备自愈能力的局部优化过程，该变异操作示意图如图\ref{fig: 变异操作示意图}所示。
    \begin{figure}[!htb]
        \centering
        \includegraphics[width=0.85\linewidth]{fig/变异操作示意图.pdf}
        \caption{变异操作示意图}
        \label{fig: 变异操作示意图}
    \end{figure}

    \subsection{终止条件}
    针对LEO卫星网络拓扑的高动态性与计算资源受限特征，遗传算法需要在有限的时间窗口内输出可行解，以平衡路由质量与计算时效性。本文采用最大迭代次数与进化停滞检测相结合的双重判定机制来控制算法流程。
    
    首先，设定最大迭代代数 $T_{max}$ 作为算法运行的硬性时间边界，当当前迭代次数 $t$ 达到 $T_{max}$ 时（$t \ge T_{max}$），强制结束迭代以防止因搜索耗时过长导致路由结果滞后于拓扑变化。同时，利用进化停滞窗口 $W_{stag}$ 监测种群的收敛状态，若全局最优适应度在连续 $W_{stag}$ 代内的改善幅度低于判定阈值 $\varepsilon$（即满足 $Fitness_{best}(t) - Fitness_{best}(t - W_{stag}) < \varepsilon$），则判定算法已进入收敛停滞状态并提前终止，以避免无效的冗余计算。
    
    只要满足上述任一条件，算法即刻停止并对当前种群中的最优染色体进行物理节点解码，生成最终的域内路由路径。

    %  算法复杂度分析。主要分析两点：时间复杂度：跟 种群规模 $M$、最大迭代次数 $T$、节点数 $N$ 的关系。主要开销在：适应度计算（遍历路径）、排序（快排）、交叉变异（线性扫描）。结论通常是 $O(T \cdot M \cdot N)$ 这种多项式级别的。空间复杂度：存种群需要的内存。

    \section{仿真实验}
    本小节将详细介绍异构卫星网络路由算法的仿真实验环境、参数设置及对比分析结果。实验旨在验证本文所提多域协同的分层路由机制（记为H-IGA），在多层异构、高动态及非均匀业务负载场景下的性能优势，特别是针对链路拥塞的缓解能力。
    \subsection{仿真环境与参数设置}
    1、仿真场景
    
    仿真实验基于 Python 语言开发，采用 STK (Systems Tool Kit) 软件生成高精度的卫星轨道星历与链路可见性数据。仿真场景构建了与\ref{sec: STK参数配置}一致的异构天地一体化网络场景，包含100 颗低轨探测卫星、225 颗低轨通信卫星、8 颗中轨卫星、3 颗地球同步轨道卫星、48 个全球分布的地面站，全网共计 384 个节点。仿真利用 STK 导出的真实星历数据，以 10秒 为时间步长，动态计算节点位置与链路可见性。为了模拟真实的异构网络特性，不同层级间的链路带宽与通信距离限制采用了非均匀配置，如表\ref{tab: 异构网络链路参数配置}所示。
    
    \begin{table}[htbp]
        \caption{异构网络主要链路参数配置}
        \centering
        \small
        \begin{tabular}{lccc}
        \toprule
        \textbf{链路类型} & \textbf{带宽 (Mbps)} & \textbf{传播时延区间 (ms)} & \textbf{基础丢包率 (\%)} \\
        \midrule
        Ground $\leftrightarrow$ Detect & 100 $\sim$ 150 & 2.3 $\sim$ 8.5 & 0.10 $\sim$ 0.50 \\
        Ground $\leftrightarrow$ LEO    & 100 $\sim$ 150 & 3.3 $\sim$ 10.2 & 0.10 $\sim$ 0.50 \\
        Detect $\leftrightarrow$ Detect & 40 $\sim$ 50 & 1.0 $\sim$ 15.0 & 0.001 $\sim$ 0.01 \\
        LEO $\leftrightarrow$ LEO       & 40 $\sim$ 60 & 1.5 $\sim$ 18.0 & 0.001 $\sim$ 0.01 \\
        MEO $\leftrightarrow$ MEO       & 200 $\sim$ 500 & 20.0 $\sim$ 60.0 & 0.001 $\sim$ 0.01 \\
        Detect $\leftrightarrow$ LEO    & 40 $\sim$ 50 & 1.0 $\sim$ 12.0 & 0.005 $\sim$ 0.05 \\
        Detect/LEO $\leftrightarrow$ MEO       & 100 $\sim$ 200 & 25.0 $\sim$ 55.0 & 0.001 $\sim$ 0.01 \\
        LEO $\leftrightarrow$ GEO       & 50 $\sim$ 80 & 119.0 $\sim$ 130.0 & 0.01 $\sim$ 0.05 \\
        \bottomrule
        \end{tabular}
        \label{tab: 异构网络链路参数配置}
    \end{table}

    为了量化链路负载对传输性能的非线性影响，本文基于经典的BPR 函数形式来构建链路代价模型，以平滑且连续地描述从轻载到过载状态下排队时延的演变规律。链路的总时延 $D_{total}$ 由传播时延 $D_{prop}$ 和动态排队时延 $D_{dyn}$ 组成：
    \begin{equation}
        D_{total} = D_{prop} \times (1 + \alpha \cdot u^\beta)
    \end{equation}
    
    其中 $u$ 为链路利用率 ($u = \text{UsedBW} / \text{Capacity}$)，$\alpha$ 和 $\beta$ 为拥塞敏感系数，仿真中分别设置为 1.0 和 6.0，以模拟高负载下时延的剧烈增长。丢包率 $L$ 同样基于负载率动态计算，当链路过载时，丢包率随负载率迅速上升：
    \begin{equation}
    L(u) =
    \begin{cases}
      0.1\% + u \times 0.1\% & u \le 1.0 \\
      0.2\% + (u-1.0) \times 80\% & u > 1.0 
    \end{cases}
    \end{equation}
    
    2、业务流量设置
    
    实验构建热点区域与背景流量混合的非均匀业务场景，以测试算法在局部高负载下的性能。对于业务分布，80\%的业务源节点集中在北美西海岸 (Lat: $30^{\circ}\sim40^{\circ}$, Lon: $-125^{\circ}\sim-115^{\circ}$)，宿节点集中在东亚区域 (Lat: $30^{\circ}\sim40^{\circ}$, Lon: $130^{\circ}\sim145^{\circ}$)，这种跨洋长距离传输强迫流量汇聚于有限的 MEO节点，制造激烈的跨层资源竞争；20\%的业务再全球范围内随机选取的源宿节点对，用于模拟背景通信负载。并且实验定义了四种不同 QoS 需求的业务流，如表\ref{tab: 多级 QoS 业务类型定义}所示。
    \begin{table}[htbp]
        \caption{多级 QoS 业务类型定义}
        \centering
        \begin{tabular}{lcccc}
        \toprule
        \textbf{业务类型} & \textbf{带宽 (Mbps)} & \textbf{时延容限 (ms)} & \textbf{丢包率容限 (\%)} & \textbf{优先级} \\
        \midrule
        指令业务  & 0.1 & 100 & \ 1.0 & 极高 \\
        视频业务  & 10 & 300 & 2.0 & 高 \\
        话音业务  & 2 & 150 & 2.0 & 中 \\
        遥感业务  &  35 & 2500 & 5.0 & 低 \\
        \bottomrule
        \end{tabular}
        \label{tab: 多级 QoS 业务类型定义}
    \end{table}

    3、算法参数设置

    为了评估本文提出的多域协同分层路由（H-IGA）的性能，实验选取了DSP 算法和SGA算法作为对比基准。DSP 算法是将异构分层网络视为扁平化的全网拓扑，仅以链路的静态传播时延为权重。在仿真中，所有类型的业务分组按照 Dijkstra 算法计算出的物理最短路径进行传输。 SGA算法 采用与 H-IGA 相同的多域分层路由架构，但在域内寻优阶段使用未改进的标准遗传算法。与 H-IGA 相比，SGA 不包含方向引导初始化机制，且其适应度函数仅考虑时延与丢包的线性加权，缺乏对链路拥塞程度的非线性惩罚项。H-IGA 的核心进化参数设置如表\ref{tab: H-IGA算法参数设置}所示。针对不同业务的 QoS 需求，H-IGA 在计算适应度时采用动态权重系数。对于时延优先型业务，时延权重 $\alpha=0.8$，丢包权重 $\beta=0.2$；对于吞吐量优先型业务，时延权重 $\alpha=0.2$，丢包权重 $\beta=0.8$。
    \begin{table}[htbp]
      \caption{H-IGA算法参数设置}
      \centering
      \begin{tabular}{lc}
        \toprule
        \textbf{参数名称} & \textbf{数值} \\
        \midrule
        种群规模  & 30 \\
        最大迭代次数 & 25 \\
        交叉概率 ($P_c$) & 0.8 \\
        变异概率 ($P_m$) & 0.2 \\
        方向引导概率 ($P_{guide}$) & 0.7 \\
        拥塞敏感度 ($\gamma$) & 2.0 \\
        拥塞感知权重 ($\lambda$) & 1.0 \\
        \bottomrule
      \end{tabular}
      \label{tab: H-IGA算法参数设置}
    \end{table}

    \subsection{仿真结果分析}
    基于前述仿真环境与参数设置，本节首先从平均端到端时延、分组丢包率及网络吞吐量三个核心维度，对比分析了 H-IGA 与基准算法（DSP、SGA）在不同并发流规模下的整体性能；随后，深入探讨了 H-IGA 算法在多业务并发场景下，针对不同 QoS 需求业务的差异化保障能力。
    
    1、算法整体性能对比分析
    
    图 \ref{fig: 算法对比-时延} 展示了三种路由算法在网络并发流数从 10 至 600 变化过程中的平均时延表现。在轻负载阶段，DSP 算法凭借其静态最短路径策略，取得了最低的平均时延（约 31ms 至 39ms）。这是因为在网络自由流状态下，LEO 层的物理路径最短，传输时延最小。H-IGA 和 SGA 由于采用了分层架构，部分流量跨层经过 MEO 链路，导致基础时延略高于 DSP（约 37ms 至 47ms），但在业务容许范围内。随着并发流数量增加，DSP 算法的时延曲线呈现出剧烈的指数级增长，由于缺乏负载感知能力，持续将流量注入带宽受限的 LEO或Detect 瓶颈链路，产生拥塞惩罚。SGA 算法虽然表现优于 DSP，在重负载情况下平均时延约为78ms，这表明分层架构本身能够通过 MEO 层分流一定程度的压力，但由于缺乏方向引导策略，SGA 在高维解空间中的搜索效率较低，难以快速找到全局最优的避拥路径。H-IGA 在重负载情况下，其平均时延仍稳定在 63ms，相比 DSP 降低了约 40\%。这表明 H-IGA 通过拥塞感知机制，有效地将非延时敏感业务引至 MEO 层，缓解了 LEO 层的排队压力。
     \begin{figure}[!htb]
        \centering
        \includegraphics[width=0.7\linewidth]{fig/三种算法对比-时延.pdf}
        \caption{三种算法平均端到端时延随全局负载的变化}
        \label{fig: 算法对比-时延}
    \end{figure}
    
    图 \ref{fig: 算法对比-丢包率} 展示了平均分组丢包率随并发流数的变化趋势。在轻负载阶段，三种算法的丢包率均维持在 约0.4\% 的较低区间，此时的丢包主要源于大气信道损耗而非网络排队。随着网络进入高负载区，DSP 算法的丢包率迅速攀升，并在重负载下达到约7\%，表明 DSP 因固守 LEO 最短路径，导致局部链路负载率激增，进而触发了严重的队列溢出。SGA 算法虽然通过分层分流将丢包率控制在 3\% 左右，但 H-IGA 算法进一步展现了较高的可靠性，其在重负载下的丢包率仅为 2\%，显著低于 DSP 和 SGA。这一优势得益于 H-IGA 充分利用跨层长路径来规避 LEO 层的局部热点，并且域内寻路时还具备拥塞感知的能力。虽然这种策略增加了路径跳数，但有效避免了拥塞导致的队列溢出，在保障数据完整性方面实现了较优的权衡。
    \begin{figure}[!htb]
        \centering
        \includegraphics[width=0.7\linewidth]{fig/三种算法对比-丢包率.pdf}
        \caption{三种算法平均丢包率延随全局负载的变化}
        \label{fig: 算法对比-丢包率}
    \end{figure}
    
    图 \ref{fig: 算法对比-吞吐量} 进一步对比了三种算法在网络吞吐量方面的表现，该指标直接反映了网络资源调度的极限能力。随着业务流数量的增加，DSP 算法的吞吐量增长最早出现饱和瓶颈。特别是在重负载阶段，其吞吐量停滞在 3050 Mbps 左右。这一现象表明，受限于LEO 层路径选择，物理链路的带宽瓶颈限制了全网容量的释放，导致新增流量无法被有效传输。相比之下，SGA 算法虽然通过分层架构提升了一定的传输能力，但 H-IGA 实现了接近线性的吞吐量增长。相比 DSP，H-IGA 的网络吞吐量提升了约70\%。这一显著提升表明了 H-IGA 的多域协同机制将大带宽需求业务可以跨层进行立体分流，从而极大提升了异构网络在高负载下的资源利用效率。
    \begin{figure}[!htb]
        \centering
        \includegraphics[width=0.7\linewidth]{fig/三种算法对比-吞吐量.pdf}
        \caption{三种算法网络吞吐量随全局负载的变化}
        \label{fig: 算法对比-吞吐量}
    \end{figure}
    
    2、多业务 QoS 保障能力分析
    
    为了进一步验证 H-IGA 针对不同 QoS 需求业务的差异化保障能力，本文提取了场景中四种业务类型在不同负载下的详细性能数据进行分析。结果显示，算法成功执行了基于业务特性的分级调度策略。

    图 \ref{fig: HIGA-不同业务时延} 展示了不同类型业务的端到端时延表现。 结果显示，不同业务的时延特征呈现出显著的分层现象。指令与话音业务的时延曲线始终稳定在 50ms 左右，且未随网络负载增加而显著波动，表明算法成功为延时敏感业务分配了 LEO 层的低时延路径。相比之下，遥感业务的平均时延上升至 98ms 左右，表明经由 MEO 层中继传输。这证实了H-IGA能够识别业务需求，主动将延时容忍度较高的业务调度至高层网络，从而降低了底层链路的负载压力。
    \begin{figure}[!htb]
        \centering
        \includegraphics[width=0.7\linewidth]{fig/HIGA-时延.pdf}
        \caption{四类业务平均端到端时延随全局负载的变化}
        \label{fig: HIGA-不同业务时延}
    \end{figure}
    
    图 \ref{fig: HIGA-不同业务丢包率} 展示了不同业务的传输可靠性。得益于拥塞感知机制，指令与话音业务的丢包率被控制在极低水平，表明关键业务有效避开了网络中的拥塞节点。对于占据主要带宽的遥感业务，虽然其跨层传输增加了路径长度，但丢包率依然处于较低的可接受范围。这说明 MEO 层的带宽资源有效支撑了大容量数据的传输，避免了流量集中在 LEO 层瓶颈链路而引发严重的排队溢出。
    \begin{figure}[!htb]
        \centering
        \includegraphics[width=0.7\linewidth]{fig/HIGA-丢包率.pdf}
        \caption{四类业务平均分组丢包率随全局负载的变化}
        \label{fig: HIGA-不同业务丢包率}
    \end{figure}
    
    图 \ref{fig: HIGA-不同业务吞吐量} 展示了多业务并发下的网络容量分布情况。遥感业务的吞吐量呈现出稳健的增长趋势，贡献了网络中绝大部分的有效流量。这一结果证明 H-IGA 算法能够合理调度多层网络资源，利用 MEO 层的大带宽特性承载高吞吐量业务，同时保留 LEO 层资源以保障延时敏感业务。这种基于业务特征的差异化调度策略，既满足了不同业务的传输需求，也提升了异构网络的整体资源利用效率。
    \begin{figure}[!htb]
        \centering
        \includegraphics[width=0.7\linewidth]{fig/HIGA-吞吐量.pdf}
        \caption{四类业务吞吐量随全局负载的变化}
        \label{fig: HIGA-不同业务吞吐量}
    \end{figure}


    \section{本章小结}
    本章针对异构星座网络拓扑高动态、业务需求多样化及网络资源受限的挑战，系统研究了面向多域动态 QoS 的协同路由算法。

    首先，为了解决路由计算复杂度高的问题，构建了基于虚拟节点的异构星座分层图模型，通过虚实结合的节点映射机制，有效屏蔽了底层 LEO 星座微观拓扑的频繁变化，为上层路由规划提供了稳定的逻辑视图。在此基础上，提出了包含全局路由规划层、域内路由决策层与数据转发层的分层协同路由架构，将复杂的端到端路由问题解耦为域间骨干规划与域内精细选路两个子问题，实现了全局优化与局部实时的平衡。
    
    其次，针对多维 QoS 约束下的选路难题，分别设计了域间与域内路由算法。在域间层面，利用基于多约束剪枝的 Dijkstra 算法，实现了基于抽象拓扑的骨干路径规划与 QoS 预算的智能分解；在域内层面，提出了基于方向引导与拥塞感知的改进遗传算法（H-IGA）。该算法通过引入基于几何距离约束的初始化策略，克服了传统遗传算法盲目搜索导致的收敛慢问题，并利用拥塞感知变异算子，实现了对网络热点链路的主动规避与自我修复。
    
    最后，通过 STK 与 Python 仿真，在非均匀业务负载场景下对所提算法进行了性能验证。仿真结果表明，相比于传统最短路算法与未改进的遗传算法，H-IGA 算法在保障多类业务差异化 QoS 需求的同时，显著降低了高负载下的网络时延与分组丢包率，大幅提升了异构星座网络的整体吞吐量与资源利用率。

	\chapter{基于容器虚拟化技术的网络仿真平台}
    \section{引言}
    前文第三章提出了面向异构星座的分层多域管控架构，第四章设计了基于改进遗传算法的多维 QoS 路由机制。虽然通过数值分析验证了上述架构与算法在理论层面的有效性，但在实际的异构卫星网络环境中，网络协议的交互、各种物理链路的动态变化以及真实业务流的传输特性往往更为复杂。单纯的数学建模或数值仿真难以真实反映网络协议栈的处理开销、控制信令的交互时延以及网络拥塞时的动态行为。因此，构建一个能够模拟真实网络环境、支持真实协议栈运行且具备大规模节点承载能力的高保真仿真平台，是对本文提出的架构与算法进行工程验证的关键环节。

    % TODO: 补充引用一些参考文献（看周益辉的论文）
    目前，卫星网络仿真主要依赖两类工具：一类是基于离散事件模拟的传统仿真软件（如 NS-3、OPNET），另一类是基于硬件实物的半实物测试床。传统仿真软件虽然能够模拟大规模节点，但其网络协议栈通常是简化的数学模型，难以支持真实的应用层业务和 SDN 控制器的直接部署，且随着节点数量增加，仿真效率会显著下降。硬件测试床虽然逼真度最高，但造价昂贵、部署周期长，难以灵活模拟拥有数千颗卫星的巨型星座及其频繁的拓扑变化。
    
    针对上述局限性，本章提出并实现了一种基于容器虚拟化 的异构卫星网络仿真平台。该平台利用 STK 软件生成高精度的卫星轨道与链路数据，作为物理世界的驱动源；利用 Docker 容器技术的轻量级特性，在有限的计算资源下模拟大规模卫星节点与地面终端，提供真实的 Linux 内核协议栈环境；利用 Open vSwitch (OVS) 和 Linux Traffic Control (TC) 机制构建可编程的虚拟链路，模拟时变的网络拓扑与信道特征。
    
    本章将详细阐述该仿真平台的总体架构设计与关键技术实现，重点解决 STK 数据与虚拟网络的动态映射、大规模节点的轻量化构建以及跨域协同路由的工程实现等难题，并在此平台上对本文提出的分层管控架构和 QoS 路由算法进行全面的性能测试与评估。

    \section{仿真平台总体架构}
    平台采用基于 SDN 的分层式虚拟化架构，设计遵循控制与转发分离、仿真驱动与网络业务分离的原则，将系统划分为四个核心平面：管理编排面、控制面、数据转发面以及外部数据支撑面。总体架构设计如图\ref{fig: 仿真平台整体架构}所示。
    \begin{figure}[!htb]
        \centering
        \includegraphics[width=0.85\linewidth]{fig/平台整体架构.pdf}
        \caption{仿真平台整体架构}
        \label{fig: 仿真平台整体架构}
    \end{figure}
 
    管理编排面负责底层虚拟化资源的调度、仿真生命周期的管理以及仿真场景的时间推进与状态演进，主要分为：（1）仿真编排控制中心，作为核心调度单元，负责协调各个功能组件的运行。它向下调用底层接口进行资源管理，向上提供仿真控制接口以响应用户指令。其中容器生命周期管理模块负责调用 Docker API，自动化执行仿真节点（卫星、地面站等）的镜像加载、容器创建、启动、停止及资源回收操作，支持大规模节点的批量部署；虚拟网络拓扑构建模块负责操作 Open vSwitch (OVS) 虚拟交换机，在宿主机内部构建虚拟网桥，并通过 Veth-pair（虚拟以太网对） 技术将容器的网络接口挂载到网桥上，建立节点间的逻辑连接。（2）平台仿真驱动器负责维护全局仿真时钟与场景演进。它周期性从数据库读取 STK 生成的链路状态数据，将物理世界的轨道运动映射为仿真环境中的拓扑变更事件，并通过调用网络编排接口，驱动底层链路的通断状态切换，以模拟高动态的物理环境。
    
    控制面负责仿真网络内的路由决策、资源分配与转发表项下发，遵循前文的分层多域协同管控架构。其中主控制器模拟部署于 GEO 层的全局逻辑控制节点。负责维护全网域间抽象拓扑视图，执行域间路由算法，计算跨域传输路径并进行端到端的 QoS 预算分解，向各域控制器下发宏观调度策略。域控制器模拟部署于 LEO/Detect 层的区域控制节点，采用分布式的 SDN 控制器集群实现。每个控制器独立管理一个逻辑划分的卫星自治域，负责实时感知域内链路状态，执行域内路由算法，并将路由策略转化为具体的 OpenFlow 流表项下发至数据面。
    
    数据面由大量轻量级容器构成的异构网络拓扑组成，负责具体的数据包处理与转发任务，主要包括：（1）节点载体层，其中卫星节点载体基于轻量级 Linux 镜像构建，内部集成 OVS 交换机，具备二层/三层交换能力；地面终端载体运行完整的 Linux 网络协议栈，支持部署 iPerf3、Ping 等真实业务工具，用于生成高保真的业务流量并进行性能检测。（2）链路模拟层用于模拟拓扑连接，利用 OVS 虚拟网桥 与 Veth-pair 技术构建容器间的逻辑通路，模拟卫星间的星间链路与星地链路的连通性；而信道特征模拟引入 Linux TC 模块挂载于虚拟网卡接口之上，由管理编排面驱动，实时注入带宽限制、传播时延与误码率，从而在逻辑链路上模拟空间物理信道的动态通信质量。
    
    物理场景仿真面负责为上层网络提供高精度的轨道动力学数据与时空态势信息。其中时变拓扑数据库作为物理场景数据与网络仿真平台之间的中间件，它存储由 STK 预生成的卫星轨道参数、链路通断状态及覆盖关系等时序数据，利用其高并发读写特性支持大规模拓扑数据的毫秒级查询。场景数据生成器负责解析 STK 导出的原始数据文件，对其进行清洗、格式化处理，并按照仿真时间片序列将其转化为标准化的链路状态数据存入数据库，为平台仿真驱动器提供物理层输入。

    图\ref{fig: 仿真平台运行流程}展示了仿真平台的整体运行流程。在启动初始化阶段，平台按照严格的时序依赖关系初始化各层资源，首先启动 Redis 服务，仿真编排控制中心将 STK 生成的离散化轨道数据全量导入数据库。在宿主机层面创建全局 OVS 网桥，并初始化基础流表规则。利用多线程池并发创建数百个卫星与终端容器，并通过 Veth-pair 技术挂载网络接口。同时，系统为终端节点自动分配 IP 地址，并将 IP-MAC 映射关系写入 Redis 地址解析库。随后启动 Ryu 域控制器集群及 GEO 主控制器，控制器加载初始配置后进入监听状态。
    \begin{figure}[!htb]
        \centering
        \includegraphics[width=0.95\linewidth]{fig/平台仿真流程图.pdf}
        \caption{仿真平台运行流程}
        \label{fig: 仿真平台运行流程}
    \end{figure}
    
    在初始化完成后，平台仿真驱动器启动主循环线程，维护内部逻辑仿真时钟 $T_{sim}$。该时钟对应 Redis 中的时间片索引，以 1s 为物理步长推进仿真场景。驱动器根据当前时钟 $T_{sim}$，从 Redis 拓扑库读取全网活跃链路快照，并重置 OVS 网桥流表状态，随后根据活跃链路快照，重新下发 add-flow 指令构建星间/星地逻辑连接。在链路连通后，驱动器调用 TC 接口，串行更新各虚拟网卡的传播时延与带宽限制规则。完成环境更新后，驱动器线程执行休眠操作以同步物理时间，随后逻辑时钟 $T_{sim}$ 自增，进入下一周期。
    
    业务动态交互阶段业务流量的处理独立于仿真循环，采用异步事件驱动机制。当测试脚本触发 Ping 或 iPerf 业务流时，首个数据包因 OVS 无匹配规则触发 Packet-In 消息。Ryu域控制器接收消息后，查询本地主机表。若判断为跨域通信，则通过 TCP Socket 请求主控制器计算跨域路径。Ryu域控制器再根据跨域路由信息计算域内路径，并下发 Flow-Mod 指令。后续数据包将直接由 OVS 转发，不受仿真驱动线程的影响。
    
    

    \section{关键功能模块设计与实现}
    本节严格遵循上一小节提出的总体架构，从物理数据预处理、数据面基础设施构建、管理面动态驱动以及控制面路由决策四个维度，阐述关键功能模块的设计与实现。
    
    \subsection{物理场景数据处理与存储实现}
    
    % 对应架构：物理场景仿真面
    % 核心内容：阐述如何将 STK 的连续轨道数据离散化为时间片快照，并设计 Redis 的分库存储结构（拓扑库、链路库、映射库）以支持高效查询。
    1、星座数据的离散化映射
    
    卫星网络仿真的物理驱动源来自于 STK 软件生成的轨道动力学数据。为了实现对全星座节点全生命周期的高效数据获取，平台摒弃了人工操作 STK 导出数据的传统低效模式，通过编写 MATLAB 自动化脚本调用 STK Object Model，驱动 STK 引擎执行轨道推演、覆盖分析及链路预算计算，并将计算结果自动格式化导出为标准中间文件。在此基础上，平台构建了一套轨道数据离散化处理流水线，将连续的物理场景转化为按时间片索引的离散拓扑快照。数据离散化的核心模型是将连续的仿真时间轴 $T$ 划分为一系列等间隔的时间片序列 $\{t_0, t_1, ..., t_n\}$，步长设为 $\Delta t$。具体处理流程如图\ref{fig: 星座数据的离散化映射}所示，包含三个关键阶段。
    \begin{figure}[!htb]
        \centering
        \includegraphics[width=0.8\linewidth]{fig/平台星座数据的离散化映射.pdf}
        \caption{星座数据的离散化映射流程图}
        \label{fig: 星座数据的离散化映射}
    \end{figure}
    
    首先利用 MATLAB 脚本初始化 STK 场景，加载定义的星座参数，再利用脚本遍历仿真周期内的每一时刻，提取全网节点（卫星-卫星、卫星-地面站）间的几何可视关系及链路距离，并按预定义格式自动生成 .txt 或 .csv 格式的原始链路通断报告。接着建立连续区间到离散索引的映射，解析程序读取导出的报告文件，遍历每条链路的可见时间窗口，将其离散化为整数时间片编号集合，构建时间片 ID $\to$ 活跃链路集合的倒排索引。 最后将处理后的离散拓扑快照按统一的数据结构序列化，写入中间件数据库，以支持后续仿真驱动器的高效读取。

    2、数据库存储
    
     在异构星座仿真过程中，数据访问同时呈现出显著的时空局部性与阶段性特征：一方面，仿真驱动器主要集中于对同一时间片内全网拓扑快照的批量读取、特定节点的链路映射查询以及静态地址解析；另一方面，在动态 QoS 路由与分层控制过程中，控制平面还需要获取当前时间片活跃链路的运行态状态信息，用于链路可用性判断、带宽接纳与路径评分。基于此特征，平台采用 Redis 逻辑分库机制，将静态/预处理数据与运行态动态状态隔离存储于独立的数据库命名空间，以提升访问效率、降低不同类型数据之间的耦合，并减少高并发访问场景下的索引干扰。各逻辑库的设计如表\ref{tab: 逻辑分库存储结构与典型用途}所示。
    \begin{table}[htbp]
        \centering
        \caption{Redis 逻辑分库存储结构与典型用途}
        \label{tab: 逻辑分库存储结构与典型用途}
        \begin{tabular}{lllll}
        \toprule
        \textbf{逻辑库} & \textbf{功能定位}  & \textbf{数据结构} & \textbf{典型用途} \\
        \midrule
        DB1 & 拓扑快照库  & List & 存储节点在各时间片的邻居列表 \\
        DB2 & 链路映射库  & Hash & 存储 邻居名到链路ID 映射 \\
        DB3 & 地址解析库  & String & 存储 IP到MAC 映射\\
        DB4 & 链路状态库  & Hash & 存储当前时间片活跃链路的时延、丢包率等信息 \\
        \bottomrule
        \end{tabular}
    \end{table}
    
    其中，拓扑快照库作为核心动态数据源，存储各节点随仿真时间推进而变化的邻接关系。键名遵循${Type}:{Name}$ 规范，实现了按节点类型和名称的快速索引，键值采用 List 结构，列表的索引位 $k$ 直接对应仿真时间片编号 $t_k$，元素为该时间片内所有可见邻居的标识符集合。示例如图\ref{fig: 键值示例-拓扑快照库}所示。
    \begin{figure}[!htb]
        \centering
        \includegraphics[width=0.7\linewidth]{fig/键值示例-拓扑快照库.png}
        \caption{拓扑快照库的数据结构示例}
        \label{fig: 键值示例-拓扑快照库}
    \end{figure}
    
    链路映射库用于存储节点间物理连接的静态属性，解决逻辑连接到物理端口的映射问题，键名为源节点名称，字段为邻居节点名称，值为对应的物理链路标识。例如，查询 L1205 到 L1206 的连接接口时，通过 “HGET L1205 L1206” 命令即可快速获取链路端口ID。示例如图\ref{fig: 键值示例-链路映射库}所示。
    \begin{figure}[!htb]
        \centering
        \includegraphics[width=0.3\linewidth]{fig/键值示例-链路映射库.png}
        \caption{链路映射库的数据结构示例}
        \label{fig: 键值示例-链路映射库}
    \end{figure}
    
    地址解析库存储网络层的辅助映射信息。在仿真初始化阶段，平台为每个地面终端分配IP，并将 IP-MAC 对应关系写入该库。在业务验证阶段，控制平面可直接查询该库以生成 ARP 代理响应，无需在仿真网络中进行广播探测，降低了控制信令开销。示例如图\ref{fig: 键值示例-地址解析库}所示。
    \begin{figure}[!htb]
        \centering
        \includegraphics[width=0.3\linewidth]{fig/键值示例-地址解析库.png}
        \caption{地址解析库的数据结构示例}
        \label{fig: 键值示例-地址解析库}
    \end{figure}

    链路状态库针对当前时间片的活跃链路记录运行时状态信息，其字段通常包括链路传播时延、丢包率、物理总容量、已用带宽、剩余带宽、链路利用率及采样时间戳等，用于支撑跨域路径剪枝、域内路径评估与链路可承载性判断。示例如图\ref{fig: 键值示例-链路状态库}所示。
    \begin{figure}[!htb]
        \centering
        \includegraphics[width=0.3\linewidth]{fig/键值示例-链路状态库.png}
        \caption{链路状态库的数据结构示例}
        \label{fig: 键值示例-链路状态库}
    \end{figure}
    
    
    \subsection{异构节点与链路载体的构建}
    % 对应架构：数据转发面（节点载体层）
    % 核心内容：（补全部分） 阐述如何利用 Docker 构建轻量级卫星/终端镜像（Alpine/OVS集成），以及如何通过 Veth-pair 和 Namespace 技术将容器网络接口映射到宿主机网桥，构建基础的静态物理网络环境。

    1、节点载体构建
    
    针对异构网络中数据转发、边缘接入与逻辑控制的功能需求，平台构建了三类定制化的 Linux 容器镜像，具体配置如表\ref{tab: Docker镜像清单}所示。在此基础上，为了规范容器与宿主机之间的连接，平台制定了网络接口映射与命名规范（见表\ref{tab: 网络接口映射与命名规范}），为后续 SDN 控制器的拓扑管理提供基础。
    \begin{table}[!htb]
        \centering
        \caption{ 容器镜像清单}
        \label{tab: Docker镜像清单}
        % 定义列对齐：左对齐 | 左对齐 | 居中 | 左对齐 | 左对齐
        \begin{tabular}{lllll}
        \toprule
        \textbf{镜像名称} & \textbf{基础环境} & \textbf{镜像大小} & \textbf{适用节点} & \textbf{核心组件} \\
        \midrule
        \texttt{ovs-tc}    & Alpine Linux  & $\sim$28 MB  & 卫星 & OVS、Traffic Control \\
        \texttt{alpine-tc} & Alpine Linux  & $\sim$13 MB  & 地面终端      & iPerf3、Ping、tcpDump \\
        \texttt{ryu-ctrl}  & Ubuntu/Debian & $\sim$737 MB & 域控制器      & Python3、Ryu、NetworkX \\
        \bottomrule
        \end{tabular}
    \end{table}


   （1）卫星节点仿真模型
   
   卫星节点负责数据路由转发与链路维护，模型采用 ovs-tc 镜像，支持 OpenFlow 与 TC 流控。在接口建模上，严格遵循表\ref{tab: 网络接口映射与命名规范}进行配置。其中，非复用的星间天线对应星间链路，由于任意时刻仅能参与一条链路构建，平台使用独立的 Veth-pair 接口（命名为 link\_{id}）模拟。该接口在宿主机端直接接入基础承载网桥 (br\_basic)，形成点对点物理连接。而对可复用的地天线对应星地链路，考虑到单根天线需同时连接多个地面站，卫星容器仅保留唯一的对地接口（命名为 link\_0）接入 br\_basic。该接口本身不具备复用能力，而是作为流量汇聚点，需配合下文所述的副网桥 (br\_sub) 实现多用户复用接入。

    \begin{table}[!htb]
        \centering
        \caption{网络接口映射与命名规范}
        \label{tab: 网络接口映射与命名规范}
        
        \begin{threeparttable}
            \resizebox{\textwidth}{!}{%
                \begin{tabular}{llllll}
                \toprule
                \textbf{节点类型} & \textbf{接口功能} & \textbf{天线特性} & \textbf{容器内接口名} & \textbf{宿主机接口名} & \textbf{物理挂载点} \\
                \midrule
                % 卫星节点部分
                \multirow{2}{*}{卫星节点} & 星间链路 (ISL) & 非复用 (点对点) & link\_\{id\} & \{节点名\}\_\{id\}  & br\_basic \\
                 & 星地链路 (GSL) & 复用型 (相控阵) & link\_0 & \{节点名\}\_0 & br\_basic \\
                \midrule
                % 地面终端部分
                地面终端 & 用户接入 & 非复用 (单天线) & link\_0 & \{节点名\}\_0 & br\_basic \\
                \midrule
                % 网桥互联部分
                \multirow{2}{*}{网桥互联} & 逻辑层入口 & N/A & -- & b\_\{节点名\}\_0  & br\_basic \\
                 & 逻辑层出口 & N/A & -- & s\_\{节点名\}\_0 & br\_sub \\
                \bottomrule
                \end{tabular}%
            }
            
            \begin{tablenotes}
				\item[] 注：\{id\} 为 Redis 数据库分配的全局唯一链路编号；b\_ 与 s\_ 分别代表 Basic 端与 Sub 端的补丁端口。
			\end{tablenotes}
        \end{threeparttable}
    \end{table}
    
    （2）地面终端仿真模型
    
    与卫星节点的交换架构不同，地面终端模型的设计侧重于真实业务流量的生成与动态接入切换。该载体基于 alpine-tc 容器镜像构建，剔除了冗余的 OVS 交换组件。依据表 \ref{tab: 网络接口映射与命名规范} 的规范，其单一网络接口直接绑定 Linux 原生 IP 协议栈，并物理接入 br\_basic。这种设计使得容器能够作为一个标准的 TCP/IP 终端，利用 iPerf3、Ping 等工具注入真实的端到端业务流量。

    2、链路载体构建
    
    为了支撑大规模异构节点的互联互通，并解决星间直连与星地动态接入在物理拓扑上的显著差异，平台在宿主机层面设计双层 OVS 网桥架构。如图 \ref{fig: 节点和链路载体模型} 所示，该架构由br\_basic与br\_sub 共同构成，实现了星间链路的直连与星地链路的复用。
    \begin{figure}[!htb]
        \centering
        \includegraphics[width=1\linewidth]{fig/平台节点和链路载体模型.pdf}
        \caption{节点和链路载体模型示意图}
        \label{fig: 节点和链路载体模型}
    \end{figure}

    (1) 星间链路仿真模型
    
    星间链路通常呈现为点对点连接的特性。因此，利用基础承载主网桥 (br\_basic)，将所有卫星节点的星间 Veth-pair 接口直接挂载于 br\_basic。星间链路如图中绿色实线所示，直接在 br\_basic 内部配置流表，来模拟不同链路之间的通断。

    （2）星地链路仿真模型
    
    星地链路具有单点对多点覆盖及高动态切换的特性。若沿用星间链路的直连模式，终端切换卫星时需频繁重构物理接口，系统开销较大。为此，平台引入接入复用副网桥 (br\_sub)，配合主网桥构建复用机制。如图 \ref{fig: 节点和链路载体模型} 中部所示，br\_basic 与 br\_sub 通过补丁端口（图中蓝色与橙色的垂直管道）相连。在该机制下，卫星的复用型对地天线与终端的接入天线均物理接入 br\_basic 保持不动。所有星地交互流量的流向如图中橙色虚线所示，被牵引 br\_sub，br\_sub可以根据实时轨道关系将流量切换至目标节点的管道接口。这种设计将物理接口的重构转化为流表的更新，从而在逻辑上实现了卫星单天线对多个地面终端的复用接入与平滑切换。
    
    3、场景举例说明

    如图\ref{fig: 场景举例说明}所示，构建一条跨越地面与空间的完整通信链路，终端节点包括F1（源端）、F2（宿端）；卫星节点包括L1（F1 的接入卫星）、L2（F2 的接入卫星）、L3（中间转发卫星）。整体的传输路径为：F1 $\rightarrow$  L1 $\rightarrow$  L3 $\rightarrow$  L2 $\rightarrow$  F2。
    \begin{figure}[!htb]
        \centering
        \includegraphics[width=1\linewidth]{fig/平台场景举例说明.pdf}
        \caption{场景举例说明}
        \label{fig: 场景举例说明}
    \end{figure}
    
    数据包从源端 F1 发出到达宿端 F2 的过程中，在宿主机底层经历了三个阶段，具体流向如图中不同颜色箭头所示：
    
    第一阶段为上行接入。地面终端 F1 生成的业务数据包经由容器接口 link\_0 发出，进入宿主机基础承载网桥 (br\_basic) 的物理接口 F1\_0。由于该流量属于星地接入性质，流表将其通过补丁端口 b\_F1\_0 牵引至上层的接入复用副网桥 (br\_sub)。在 br\_sub 内部，根据当前的星地接入关系，流量被转发至卫星 L1 对应的补丁端口 s\_L1\_0，随后经由 b\_L1\_0 回落至 br\_basic，最终进入卫星 L1 的容器接口 link\_0。
    
    第二阶段为星间转发。数据包进入卫星 L1 后，经容器内 OVS 网桥查找路由，从星间接口 link\_101 发出。流量再次进入 br\_basic，此时流表直接将其转发至下一跳卫星 L3 的物理接口 L3\_101。同理，L3 完成内部交换后，经由 link\_102 将数据包转发至卫星 L2。此阶段数据包全程在 br\_basic 物理层进行水平转发，模拟了星间链路的直连。
    
    第三阶段为下行落地。卫星 L2 将数据包转发至对地接口 link\_0。流量进入 br\_basic 后，再次触发星地映射逻辑，被牵引至 br\_sub。逻辑层根据路由策略，将数据包交换至地面终端 F2 对应的补丁端口，最终回落并送达 F2 的协议栈。
    
    该实例表明，通过主副网桥的配合，平台能够准确地将上层逻辑拓扑映射为底层的物理数据流，实现了复杂跨域路径的连通。
    
    \subsection{仿真驱动与动态拓扑管理实现}
    % 对应架构：管理编排面 + 数据转发面（链路模拟层）
    % 核心内容：阐述平台仿真驱动器如何维护全局时钟，通过“状态轮询与同步”机制驱动底层 OVS 流表的下发/清除（实现链路通断），以及如何调用 TC 注入时变信道参数（实现链路特征模拟）。
    仿真驱动器是连接物理轨道数据与虚拟网络环境的执行核心，负责维护全局仿真时钟并协调底层网元与上层控制器的状态同步。图\ref{fig: 仿真驱动与拓扑管理}描绘了驱动器在离散时间片内的完整执行流程。
    \begin{figure}[!htb]
        \centering
        \includegraphics[width=0.7\linewidth]{fig/平台驱动与拓扑管理流程图.pdf}
        \caption{仿真驱动与拓扑管理流程图}
        \label{fig: 仿真驱动与拓扑管理}
    \end{figure}
   % 1、基于离散时间片的全局时钟推进
   
   为了适配大规模星座的长时间跨度仿真，平台利用 Redis 存储预切分的离散时间片数据。驱动器运行在独立的主线程中，维护一个全局逻辑时钟 $T_{sim}$，采用读取、执行、步进的循环模式推进仿真。在每个周期开始，驱动器根据当前时钟 $T_{sim}$ 从 Redis 原子化读取对应的全网拓扑快照，随后依次执行拓扑重构、控制器同步与信道参数注入操作，确保当前时间片内的网络状态达成一致。任务完成后，驱动器执行间隔控制机制，为了给底层 OVS 网桥的流表刷新和上层 SDN 控制器的路由收敛预留稳定窗口，防止因状态更新频率过快导致控制平面拥塞。在完成等待后，逻辑时钟 $T_{sim}$ 自增，进入下一仿真周期。
   
   % 2. 基于差分的动态拓扑重构
   
   针对卫星网络拓扑频繁变化的特性，若在每个时间片都全量重置网络状态，将产生巨大的系统开销。为此，平台采用了基于差分的增量更新策略（如图中分支 A 所示）。驱动器在内存中缓存了上一时刻 $T_{sim-1}$ 的全网链路状态，在时刻 $T_{sim}$ 将当前拓扑快照与缓存状态进行比对。对于识别出的新增链路集合 ($L_{add}$)，驱动器向 OVS 下发 add-flow 指令，其中星间链路在 br\_basic 建立点对点转发，星地链路则协同配置 br\_basic 与 br\_sub 打通逻辑环路；对于失效链路集合 ($L_{del}$)，则下发 del-flow 指令清除对应流表项。通过仅处理变更部分，OVS 的指令交互量大幅降低，从而有效支持了大规模节点的连续仿真。
   
   % 3. 基于文件共享的控制器同步
   
   为了使 SDN 控制器能够感知底层流表的变化，平台建立了从仿真驱动层到控制平面的单向数据注入通道（如图中分支 C 所示）。驱动器在完成流表更新后，根据当前拓扑数据生成包含节点接入关系和邻接关系的标准文本文件，并利用 Docker API 的文件复制功能将上述文件直接注入到运行中的 Ryu 控制器容器指定目录中，实现了从仿真环境到控制逻辑的状态传导。
   
   % 4. 信道特征的物理层模拟
   在维护连接关系的同时，为了在虚拟网络中复现真实的物理信道特征，平台利用 Linux TC工具对容器接口进行动态配置（如图中分支 B 所示）。在每个仿真周期内，驱动器查询 Redis 链路映射库获取当前活跃链路的物理参数，并将其转化为 TC 指令：一方面根据链路距离 $d$ 计算传播时延 $t = d/c$，通过 netem delay 命令注入到对应的容器接口；另一方面根据链路预算结果，通过 tbf rate 和 netem loss 命令限制接口带宽并设置随机丢包率。通过上述机制，平台成功将预处理的 STK 轨道数据转化为动态的、具备真实物理特征的虚拟网络环境。
    
    \subsection{平台协同控制与路由的设计}
    % 对应架构：控制面
    % 核心内容：阐述 Ryu 域控制器的初始化与拓扑感知，以及基于 Socket 的跨域协同机制，实现从“业务感知”到“跨域路径计算”再到“流表下发”的完整控制闭环。
    本节阐述控制平面的工程实现。系统采用分层架构，底层由多个部署在容器中的 Ryu 域控制器负责域内流表下发，顶层部署全局主控制器负责跨域路径计算。两者通过 TCP Socket 接口交互，实现完整的路由控制。图\ref{fig: 协同控制与路由时序图}展示了过程的时序图。
    \begin{figure}[!htb]
        \centering
        \includegraphics[width=0.9\linewidth]{fig/平台协同控制与路由时序图.pdf}
        \caption{协同控制与路由时序图}
        \label{fig: 协同控制与路由时序图}
    \end{figure}

    网络完成初始化之后，Ryu 首先向全局主控制器发起 TCP Socket 连接以建立跨域通信通道；随后，底层卫星节点的 OVS 网桥主动向本域 Ryu 控制器发起 OpenFlow 握手并完成设备注册。在此基础上，仿真驱动器将生成的拓扑描述文件注入控制器容器，触发 Ryu 的内存刷新机制。其中，物理拓扑图 (topology\_info.txt) 记录了域内卫星间的邻接关系，控制器解析文件并利用 NetworkX 库构建有向图，节点为卫星 DPID，边为星间链路，作为域内路由算法的计算基础；终端位置表 (host\_info.txt) 则记录终端 MAC与接入卫星 DPID的映射关系，该表具有双重作用：既存储本域真实接入的终端位置，也用于存储跨域查询后获得的远程终端位置，即映射为本地边界卫星。

    当业务终端发送首个数据包时，因 OVS 查表失败，数据包被封装为 Packet-In 消息上送控制器。Ryu 解析 IP 头部并在本地终端位置表中查询目的 MAC 地址，根据查询结果执行分级路由计算。若本地表中存在目的地址记录，则属于域内通信，直接获取其对应的真实接入卫星 ID 作为路径终点。若表中无记录，则判定为域间通信，Ryu 向全局主控制器发送包含 $\{srcMac, dstMac, domainId\}$ 的 Socket 请求；主控制器根据宏观拓扑计算跨域路径并确定本域的出境边界卫星，返回 $\{dstMac, dpid, port\}$ 响应。Ryu 接收响应后，将目的 MAC 与 边界卫星 ID的映射关系追加写入 host\_info 表，随即以该边界卫星为终点计算域内路径。

    完成路径计算后，系统进入执行阶段。控制器根据计算出的路径序列，向沿途 OVS 下发 Flow-Mod 指令，建立端到端的转发规则。为了保障业务的连续性，控制器同时发送 Packet-Out 消息，指示 OVS 将当前缓存的首个数据包转发至下一跳。此后，该业务流的后续报文到达 OVS 后，将直接匹配已安装的流表进行硬件转发，不再经过控制器处理。

    \section{平台功能验证与结果分析}
    % 5.4.1 实验环境与场景配置
    % 目的：交代“在哪跑的”、“跑的什么业务”。
    % 核心内容：硬件/软件环境：宿主机的 CPU、内存、操作系统（Ubuntu 20.04?）、Docker 版本、OVS 版本、Ryu 版本。
    % 仿真拓扑参数：为了验证平台，您选取了什么星座？（例如：选取 Starlink 第一壳层的一个子集，或者典型的 Walker Delta 星座，如 6个轨道面 $\times$ 11颗卫星）。
    % 业务流量模型：使用 iPerf3 打什么流（UDP/TCP？带宽多少？），用 Ping 测什么？

    \subsection{实验环境与场景配置}
    为了验证仿真平台的功能完备性与性能指标，本章搭建了基于高性能服务器的实验环境，并构建了包含大规模异构节点的仿真场景及多维度的业务流量模型。
    
    1、软硬件环境配置
    
    实验平台部署于高性能服务器节点，软件环境基于 Ubuntu 22.04 LTS，利用 Docker 实现节点轻量化，通过 Open vSwitch 构建虚拟网络，并采用 Redis 存储时变拓扑。各核心组件配置详见表\ref{tab: 实验平台软硬件环境配置}。
    \begin{table}[!htb]
        \centering
        \caption{实验平台软硬件环境配置}
        \label{tab: 实验平台软硬件环境配置}
        \begin{threeparttable}
            \begin{tabular}{lll}
                \toprule
                \textbf{类别} & \textbf{组件} & \textbf{配置/版本} \\
                \midrule
                % 硬件部分
                \multirow{4}{*}{硬件环境} & CPU & Intel Xeon Gold 6145 @ 2.0GHz \\
                                         & 内存 & 251 GB (DDR4) \\
                                         & 存储 & SSD 固态存储 \\
                                         & 网络 & 千兆以太网 \\
                \midrule
                % 软件部分
                \multirow{4}{*}{软件环境} & 操作系统 & Ubuntu 22.04.5 LTS \\
                                         & 容器引擎 & Docker 20.10.7 \\
                                         & 虚拟交换机 & Open vSwitch 2.17.9 \\
                                         & 数据库 & Redis 6.0.16 \\
                \bottomrule
            \end{tabular}
        \end{threeparttable}
    \end{table}

    2、仿真拓扑参数
    
    仿真场景基于 STK 高精度轨道数据构建，涵盖 Detect/LEO/MEO/GEO 四层节点，具体参数严格遵循表 \ref{tab: 特定异构星座仿真参数配置}。管控层面部署多控制器集群实现分层协同，所有节点的时变位置、邻接关系及链路接口映射均经过预处理并存储于 Redis 数据库中，仿真驱动器通过读取数据库中的拓扑快照，以 驱动整个网络的拓扑演进。
    

    3、业务流量工具
    
    为全面评估平台的网络连通性、传输性能及 QoS 保障能力，设计了以下三类业务流量模型：（1）连通性测试：利用 Ping 工具检测节点间的网络可达性。通过记录 ICMP 报文的响应情况与往返时延，验证单域内及跨域路径的连通状态，以及在拓扑动态切换过程中的连接保持能力。（2）性能测试：利用 iPerf3 工具生成端到端业务流量。测试涵盖 TCP 与 UDP 两种协议模式，重点评估特定带宽限制下的吞吐量表现、抖动情况以及多跳路径带来的性能开销。（3）QoS 信道模拟：依据第四章（表 \ref{tab: 异构网络链路参数配置}）设定的异构网络链路模型，对虚拟接口动态注入带宽、传播时延与基础丢包约束。并在业务报文头中写入 DSCP 标记，并在转发面通过队列映射实现差异化调度。其中，链路物理约束由 Linux TC 机制注入，用于模拟信道特性；队列优先级由OVS QoS与多队列策略实现，用于体现不同 DSCP 业务在竞争条件下的调度差异。
    
    \subsection{平台基础功能验证}
    % 目的：证明 5.3 节写的那些机制（动态拓扑、跨域路由）是真的实现了。
    % 核心内容（图表党）：
    % 动态拓扑可视化的验证：放两张不同时间片的截图（比如 T=0s 和 T=60s），证明链路断了、重连了。
    % 端到端连通性验证：
    % Ping 测试截图：展示 F1 Ping F2，TTL 变化或 RTT 变化。
    % TraceRoute 截图：展示路径跳数，证明经过了 L1 -> L3 -> L2。
    % 协议交互验证（可选）：Wireshark 抓包截图，展示自定义的 Socket 报文或者 OpenFlow 的 Packet-In 消息，证明控制面在工作。
    为排除大规模场景下的干扰因素，本节选取了一个典型的跨域通信场景，重点验证仿真平台在环境初始化、动态拓扑驱动、跨域协同控制及端到端数据转发等层面的基础功能。实验场景选取从源节点 Detect0205（隶属 D02 域）到目的节点 F2（隶属 D01 域）的通信路径。该路径跨越两个逻辑自治域，涉及 LEO 层的多次星间跳变，能够完整覆盖分层多域架构的核心逻辑。

    1、网络基础环境初始化
    
    在业务运行前，首先对底层容器网络与控制平面的连接状态进行核查。通过 Docker 接口遍历路径上涉及的LEO 卫星节点，检查其 OVS 网桥状态及 OpenFlow 控制通道。如图\ref{fig: 容器与Ryu控制器建立连接}所示，节点的 OVS 网桥均已启动。隶属 D02 域的 L0207 与 L0206 节点连接至控制器 IP 172.17.0.11（Ryu02），而隶属 D01 域的 L0205 至 L0203 节点连接至控制器 IP 172.17.0.10（Ryu01）。证实了平台已成功将物理层的卫星节点正确划分至不同的逻辑管控域中，分层多域架构在部署层面已就绪。
    \begin{figure}[!htb]
        \centering
        \includegraphics[width=0.8\linewidth]{fig/平台截图/容器与Ryu控制器建立连接.png}
        \caption{容器与Ryu控制器建立连接}
        \label{fig: 容器与Ryu控制器建立连接}
    \end{figure}
    
    2、动态拓扑管理验证
    
    为了验证仿真平台对卫星网络高动态特性的支撑能力，实验重点分析了时间片切换前后底层 OVS 网桥转发逻辑的演进。实验选取 $T=0s$与 $T=33s$两个时刻，对比分析了负责星地接入的 br\_sub与负责星间转发的 br\_basic的流表状态。
    
    首先验证终端接入关系的逻辑变更，如图\ref{fig: br_sub关键端口映射}所示，在 br\_sub中，端口 1 对应地面终端 F2，端口 7 对应卫星 LEO0203，端口 8 对应卫星 LEO0202。对比图\ref{fig: br_sub流表变更}所示流表快照发现，在 $T=0s$ 时，流表项为 in\_port=1, actions=output:7，表明终端 F2 此时接入 LEO0203 节点；在 $T=33s$ 时，流表项自动更新为 in\_port=1, actions=output:8。这一变化证实，随着卫星过顶运动，仿真驱动器成功检测到接入条件变更，并向下发了流表更新指令，完成了终端 F2 从 LEO0203 到 LEO0202 的逻辑切换。
    \begin{figure}[!htb]
        \centering
        \subfloat[]{
            \includegraphics[width=0.5\linewidth]{fig/平台截图/br_sub关键端口映射表.png}
            \label{fig: br_sub关键端口映射}
        }
        \hfill
        \subfloat[]{
            \includegraphics[width=0.6\linewidth]{fig/平台截图/br_sub流表变更.png}
            \label{fig: br_sub流表变更}
        }
        \caption[]{br\_sub 的关键信息。\subref{fig: br_sub关键端口映射}端口映射；\subref{fig: br_sub流表变更}流表变更} 
        \label{fig: br_sub的关键信息}
    \end{figure}
    
    进一步验证底层物理链路的实体切换，图\ref{fig: br_basic关键端口映射}展示了 big\_br 的关键端口映射表，其中端口 9 对应卫星 LEO0203 的星地接口，端口 11 对应卫星 LEO0202 的星地接口。结合图\ref{fig: br_basic流表变更}的流表快照分析，在 $T=0s$ 时，流表显示仅有端口 9 (LEO0203) 处于转发状态；而当 $T=33s$ 时，端口 9 的流表被移除，取而代之的是端口 11 (LEO0202) 的转发规则。说明仿真平台成功驱动底层物理链路完成了从 LEO0203 到 LEO0202 的实体切换。
    \begin{figure}[!htb]
        \centering
        \subfloat[]{
            \includegraphics[width=0.5\linewidth]{fig/平台截图/br_basic关键端口映射表.png}
            \label{fig: br_basic关键端口映射}
        }
        \hfill
        \subfloat[]{
            \includegraphics[width=0.6\linewidth]{fig/平台截图/br_basic流表变更.png}
            \label{fig: br_basic流表变更}
        }
        \caption[]{br\_basic 的关键信息。\subref{fig: br_basic关键端口映射}端口映射；\subref{fig: br_basic流表变更}流表变更} 
        \label{fig: br_basic的关键信息}
    \end{figure}
    
     3、控制面协同路由机制验证
     
     实验截取了业务触发瞬间的全局控制日志与底层卫星节点的 OVS 流表状态，分析了从路由计算到策略执行的过程。如图\ref{fig: 跨域业务的请求与响应}所示，当源节点 D0205 向目的节点 F2 发起通信时，D2 的域控制器（Ryu02）检测到跨域需求，随即向全局主控制器发起路径查询。依据全局拓扑实时计算出域间路径序列 [D2, D1]，并决策选取 LEO0206 至 LEO0205 作为当前时间片的最佳物理链路，最终向域控制器下发了边界节点 206 转发至端口 8400的路由策略。
     \begin{figure}[!htb]
        \centering
        \includegraphics[width=0.8\linewidth]{fig/平台截图/跨域业务的请求与响应.png}
        \caption{跨域业务的请求与响应}
        \label{fig: 跨域业务的请求与响应}
    \end{figure}
    接收到全局指令后，Ryu02 计算出域内路径。如图\ref{fig: 卫星容器的转发规则}所示，接入节点 LEO0207 将流量直接导向下一跳端口 8736；而在边界卫星 LEO0206 的 OVS 流表中，观测到匹配该业务流的规则将数据从入口 8401 准确转发至指定的跨域出口 8400，将数据包转发至下一个域。
    \begin{figure}[!htb]
        \centering
        \includegraphics[width=0.8\linewidth]{fig/平台截图/容器内部br0流表.png}
        \caption{卫星容器的转发规则}
        \label{fig: 卫星容器的转发规则}
    \end{figure}

    4、跨域端到端连通性验证

    最后，利用 Ping 工具进行Detect0205到F2之间的连通性测试。根据仿真场景配置，TC参数设定为：接入链路之间的时延为4ms，LEO星间转发链路的时延 10ms，如图\ref{fig: 链路TC设置示例}所示。
    \begin{figure}[!htb]
        \centering
        \includegraphics[width=0.7\linewidth]{fig/平台截图/链路TC设置.png}
        \caption{链路TC设置示例}
        \label{fig: 链路TC设置示例}
    \end{figure}
    
    图 \ref{fig: 跨域端到端连通性验证} 展示了两个时间片下的实测 Ping 结果。图 \ref{fig: t0跨域端到端连通性验证} 展示了在 $T=0s$ 时，实际路径为 Detect0205 $\leftrightarrow$ ... $\leftrightarrow$ LEO0203 $\leftrightarrow$ F2，包含 2 段接入链路和 4 段星间链路，其理论的 RTT 为 $2 \times (4 + 4 \times 10 + 4) = 96.0ms$，实测 RTT 均值为 97.26ms，偏差仅为 1.26ms；图 \ref{fig: t33跨域端到端连通性验证} 展示了在 $T=33s$ 时，随着 F2 切换接入至 LEO0202，路径新增一跳星间链路，理论 RTT 增加至 116.0ms，实测 RTT 均值为 117.45ms，与理论预期高度吻合。
    \begin{figure}[!htb]
        \centering
        \subfloat[]{
            \includegraphics[width=0.7\linewidth]{fig/平台截图/t0-时间片端到端连通验证.png}
            \label{fig: t0跨域端到端连通性验证}
        }
        \hfill
        \subfloat[]{
            \includegraphics[width=0.7\linewidth]{fig/平台截图/t33-时间片端到端连通验证.png}
            \label{fig: t33跨域端到端连通性验证}
        }
        \caption[]{跨域端到端连通性验证。\subref{fig: t0跨域端到端连通性验证}$T=0s$；\subref{fig: t33跨域端到端连通性验证}$T=33s$} 
        \label{fig: 跨域端到端连通性验证}
    \end{figure}

    \subsection{综合性能验证与分析}
    % 目的：用数据说话，证明平台的价值以及第四章算法的工程有效性。
    % 核心内容（数据党）：
    % 平台资源开销（证明轻量化）：随着节点数量增加（10 -> 50 -> 100），CPU 和 内存的占用曲线。对比虚拟机（如果有数据的话，没有就只放 Docker 的）。
    % 业务传输性能（证明 QoS/信道模拟）：
    % 吞吐量测试：iPerf3 跑出来的带宽图。
    % 时延/丢包率：展示 TC 注入的误码率生效了。
    % 路由算法对比（回扣第四章）：
    % 在平台上运行您的“多维 QoS 路由算法” vs “传统 Dijkstra”。
    % 对比指标：端到端平均时延、丢包率、网络吞吐量。
    % 结论：证明您的算法在“真实”协议栈环境下依然优于传统算法。
    在完成平台基础连通性与动态拓扑切换功能验证后，本节旨在搭载真实 Linux 网络协议栈与底层 OVS 虚拟交换机的系统环境下，对本文提出的多域管控机制与 H-IGA 协同路由算法开展系统级的综合验证。本节按照机制有效性和验证算法性能评估的逻辑框架展开，重点看平台在动态拓扑下的控制闭环能力、多维 QoS 差异化保障效果、控制面路由开销，以及 H-IGA 算法在真实业务并发场景下的性能优势。

    1、跨域 QoS 路由控制闭环验证

    本小节以 H-IGA 算法组中的一条指令业务为例，验证系统跨域 QoS 路由的控制闭环机制。从控制平面时序分析，系统成功完成了链路状态感知、业务触发下发、域内按需求解、数据面收敛的全流程。仿真启动后，系统首先执行动态链路状态的覆盖性自检。日志实证表明，在当前时间片下，全网 1022 条活动星间链路实现了 100\% 的动态状态采集。在此基础上，全局根控制器采用按需计算策略，完成了链路状态收集和拓扑视图构建，并面向相关自治域同步了业务集合，确保了后续路由决策严格建立在实时动态状态与业务画像均已就绪的基础之上，如图 \ref{fig: 状态就绪} 所示
    \begin{figure}[!htb]
        \centering
        \includegraphics[width=0.7\linewidth]{fig/平台截图/状态就绪.png}
        \caption{动态链路状态采集与拓扑视图构建}
        \label{fig: 状态就绪}
    \end{figure}
    
    业务流触发后，域内控制器在跨域入口处向全局控制器发起 QoS 请求。全局控制器成功解析该请求，提取出相应的QoS指标。基于全局抽象拓扑生成域间路径（[D2, D1]）后，全局控制器在跨域边界节点（LEO0106）直接下发了包含 queue=0 动作的底层物理流表。随后，根控制器同步完成了跨域与域内链路的 QoS 预算分解，扣除跨域物理传输开销后，为局部域划定了 48.5 ms 的独立时延预算，并将最终决策回传至局部控制器。如图 \ref{fig: 跨域QoS路由请求解析与响应} 所示，该交互证明宏观跨域控制面不仅成功打通了跨域信令，更将抽象的 QoS 预算精确映射为了数据面可执行的硬件队列隔离动作。
    \begin{figure}[!htb]
        \centering
        \includegraphics[width=0.7\linewidth]{fig/平台截图/跨域QoS路由请求解析与相应.png}
        \caption{跨域QoS路由请求解析与响应}
        \label{fig: 跨域QoS路由请求解析与响应}
    \end{figure}

    局部域控制器在接收到跨域指令后，依据分配的局部 QoS 预算独立执行 H-IGA 路由求解。如图\ref{fig: 域内H-IGA路由求解结果}表明，D1 与 D2 域均完成了路径规划，且输出结果包含 hard\_ok=True 标识，验证了所生成物理路径满足时延与带宽等硬约束。
    \begin{figure}[!htb]
        \centering
        \includegraphics[width=0.7\linewidth]{fig/平台截图/域内动态求解结果.png}
        \caption{域内H-IGA路由求解结果}
        \label{fig: 域内H-IGA路由求解结果}
    \end{figure}
    
    为兼顾端到端路径的一致性与底层转发的精确性，本平台在控制面实现了业务级路由缓存与协议级流表安装的逻辑解耦。具体而言，路由规划与缓存以业务实体为唯一键值，而底层 OVS 转发规则仍基于精确的协议五元组下发。图\ref{fig: 业务级路径缓存复用与协议级流表安装} 清晰验证了这一过程，首次到达的 TCP 控制报文（proto=6）触发了真实的全局路由计算；随后，当同一业务的 UDP（proto=17）与 ICMP（proto=1）探测首包到达时，虽因底层缺乏对应协议流表而再次引发 Packet-In 事件上送控制器，但控制器直接命中了已构建的业务级路径缓存（cache\_hit=True），仅耗时 1 $\sim$ 2 ms 执行协议专属流表的本地安装，未再触发重复算路。该机制在不增加控制面计算开销的前提下，确保了不同探测协议严格复用同一物理路径，并消除了路由分歧现象。
    \begin{figure}[!htb]
        \centering
        \includegraphics[width=0.7\linewidth]{fig/平台截图/业务级路径缓存复用与协议级流表安装.png}
        \caption{业务级路径缓存复用与协议级流表安装}
        \label{fig: 业务级路径缓存复用与协议级流表安装}
    \end{figure}

    \begin{figure}[!htb]
        \centering
        \includegraphics[width=0.7\linewidth]{fig/平台截图/业务传输性能统计.png}
        \caption{业务传输性能统计}
        \label{fig: 业务传输性能统计}
    \end{figure}

   最终数据平面的端到端探测如图 \ref{fig: 业务传输性能统计} 所示，针对该业务实施的 UDP、ICMP 与 TCP 三类端到端探测均返回 success=true，说明该业务已在当前场景下完成正常收敛。其中，UDP 探测在1 Mbps 的设定发送速率下获得了 1.00 Mbps 的接收吞吐量，丢包率为 1.16\%；ICMP 探测测得平均往返时延为 116.76 ms，丢包率为0.00\%；TCP 探测作为补充性容量测量，得到同一路径上的 TCP goodput 为 7.00 Mbps。同时，系统级看门狗日志表明测试进程正常启动并在任务完成后按期安全退出，全程未触发资源过载熔断。

    \section{本章小结}
    本章针对传统的纯数值仿真或离散事件模拟难以真实反映网络协议栈交互过程的局限性，设计并实现了一个基于容器虚拟化技术的异构星座网络仿真平台。

    首先，在总体架构上，平台采用了控制与转发分离、仿真驱动与网络业务分离的设计思想，构建了以 STK 轨道动力学数据为物理驱动源、以 Docker 容器为网络载体的协同架构。其次，在核心功能实现方面，提出了一种基于时间片离散化的场景映射机制，利用 Redis 高并发内存数据库解决了大规模时变拓扑的实时同步问题。通过定制化的轻量级容器镜像与双层 OVS 虚拟网桥的配合，成功构建了支持星间点对点直连与星地高频切换复用的虚拟链路载体。同时，结合 Linux TC 模块注入了动态的物理信道特征，并打通了 SDN 域控制器与全局控制器的协同控制闭环，实现了从路由计算到底层流表下发的工程落地。

    最后，依托高性能服务器搭建了系统级的测试环境，对平台进行了全面的验证与分析。在基础功能验证中，确认了平台对动态拓扑切换与控制面连通性的精准支撑；在综合性能验证中，通过注入真实业务流量，充分验证了本文第三章提出的分级多域管控架构与第四章设计的 H-IGA 协同路由机制在真实网络协议交互下的工程可行性与多维 QoS 保障优势。
    
    \chapter{全文总结与展望}
    \section{全文总结}
    随着空间信息技术的快速发展，由GEO、MEO和LEO深度融合的异构星座网络已经成为未来天地一体化信息网络的重要演进方向。然而，节点规模庞大、拓扑高动态变化以及业务QoS需求的多样化，给传统的组网管控和路由机制带来了很大的挑战。针对这些问题，本文围绕异构星座的组网架构、多域协同路由算法以及基于容器的仿真验证平台开展了具体的研究工作，主要的成果可以总结为以下三个方面：

    （1）在组网与管控架构方面，针对异构星座拓扑变化快、管理开销大的问题，提出了一种基于SDN的分级多域协同管控架构。架构采用GEO作为全局主控、LEO进行域内协同的模式，将控制权合理分级。同时，设计了基于地理网格的虚拟节点映射机制和基于虚拟地址的双层报文封装方案。这种设计使得底层的物理拓扑不论如何高速运动，上层的逻辑拓扑都能保持相对静止，有效屏蔽了物理高动态变化对网络管控的影响，大幅降低了系统维护控制信令的开销。
    
    （2）在多域路由算法方面，针对多维QoS约束和局部链路易拥塞的问题，设计了一套域间宏观规划配合域内微观寻优的协同路由机制。在域间，由GEO控制器基于抽象拓扑进行跨域路径计算并分解QoS预算；在域内，提出了一种改进的遗传算法（H-IGA）。为了解决传统遗传算法在复杂网络中搜索盲目的问题，算法引入了基于几何距离的方向引导初始化策略，并在适应度函数中加入了非线性拥塞惩罚项。仿真结果表明，该算法能够主动规避热点链路，在满足时延和丢包率要求的同时，有效提高了网络的整体吞吐量。

    （3）在网络仿真平台建设与验证方面，为了克服传统数值仿真或离散事件模拟无法运行真实协议栈的缺点，设计并实现了一个基于容器虚拟化技术的高保真异构星座仿真平台。平台将STK生成的物理轨道数据作为底层驱动，利用Redis数据库进行状态同步，并结合Docker容器、OVS虚拟网桥和Linux TC模块，构建了能够真实反映时变拓扑和信道损伤的虚拟网络环境。通过在平台上注入iPerf3等真实业务流进行测试，充分验证了本文提出的架构与路由算法在实际工程协议交互下的可行性与性能优势。

    \section{后续工作展望}
    受限于研究时间、计算资源以及现有实验条件，本文的研究工作虽然取得了一定成果，但仍有一些细节需要进一步探讨和完善。面向未来更复杂的空天网络环境，下一步的研究可以从以下几个方向展开：

    （1）引入更智能的路由决策模型。本文在域内路由中采用的H-IGA算法虽然提升了寻优效率，但在面对未来动辄几万颗卫星组成的巨型星座时，启发式算法的计算开销依然是一个挑战。未来可以考虑引入深度强化学习（DRL）或图神经网络（GNN），让卫星节点具备自主学习和预测网络拥塞状态的能力，从而实现更加快速、端到端的智能路由调度。

    （2）开展更大规模的半实物及全实物验证。本文的仿真平台主要依赖Docker容器进行轻量化部署，虽然成功跑通了真实的Linux协议栈，但在计算能力和硬件接口上与真实的星载路由器仍存在一定差异。后续研究可以尝试接入实际的星载交换硬件设备，构建硬件与软件的混合测试床，从而更准确地评估算法在星上受限硬件环境中的执行效率和资源消耗。

    （3）增强网络的安全与容灾生存能力。本文在设计管控架构时，主要考虑了理想状态下节点的协同工作。但在实际的太空环境中，卫星节点随时可能因为空间辐射、硬件故障而宕机，甚至面临恶意的网络攻击。如何在全局主控制器短时失联，或部分管理域瘫痪的情况下，依靠底层节点的分布式回退机制来保障基础通信链路的生存能力，是异构组网研究中不可忽视的工作。
    
    % 致谢部分
	\acknowledgement
	TODO
	
	% 参考文献部分
	% \nocite{*}% 为了便于在示例文档中展示参考文献而设，正式撰写时需要注释掉
	% \bibliographystyle{DissertUESTC}% 自v25.03.12版本后，参考文献风格文件由用户设定的论文类型选项自动确定，无需在此手动设置
	\bibliography{ref}
	
	% 附录起始位置
	\appendix
	
	\achievement % 仅研究生用
    TODO
	
	% \section*{发表论文：}
	
	% \begin{enumerate}
	%     \item \textbf{作者1}，作者2*，作者3，作者4. Domain decomposition method based on integral equation for solution of scattering from very thin，conducting cavity. \emph{IEEE Transactions on Antennas and Propagation}，2014，62(10): 5344--5348. (\textbf{CCF评级}，\underline{中科院分区}，IF: 98.8)
	    
	% 	\setcounter{enumi}{98}
	    
	% 	\item \textbf{作者1}，作者2*，作者3，作者4. Domain decomposition method based on integral equation for solution of scattering from very thin，conducting cavity. \emph{IEEE Transactions on Antennas and Propagation}，2014，62(10): 5344--5348. (\textbf{CCF评级}，\underline{中科院分区}，IF: 98.8)
	% \end{enumerate}
	
	% \section*{发明专利：}
	

	% \begin{enumerate}
		
	% 	\item \textbf{作者1}，作者2*，作者3，作者4。一种基于xxxxx的方法: ZL201120846830.0. 2023--02--20.
		
	% \end{enumerate}
	
	% \section*{参与项目：}
	
	% \begin{itemize}
	% 	\item 项目号. 项目名称. 项目级别，2020.01--2022.12.
	% \end{itemize}

\end{document}